\begin{document}
\title{\bi{Uniform Laughlin gaps on large flat tori\\from finite Fock-space certificates}{由有限Fock空间证书推出\\大周期平坦环面上的Laughlin一致谱隙}}
\author{\normalsize prepared by Xiang-Jian Hou with advice from Ying-Hai Wu}
\date{\bi{Candidate proof for independent review\\Revised 16 October 2026; mathematical draft of 12 October 2026}{供独立复核的候选证明\\2026年10月16日修订；数学原稿日期为2026年10月12日}}
\maketitle
\begingroup\small
\noindent\bi{\textbf{Status.} This is an AI-developed candidate proof for independent review. AI agents made substantive contributions to the torus derivations, certificate calculations, verification programs, and manuscript. The preparation and advice credited above do not assert line-by-line human verification or endorsement of every claim.}
{\textbf{状态说明。} 本文是主要由AI发展、供独立复核的候选证明。AI代理实质参与了环面推导、证书计算、验证程序及文稿写作。上方列明的整理和建议不表示人类已逐行核验或背书每项结论。}
\par\smallskip
\noindent\bi{Two selected Lean endpoints have recorded kernel acceptance: a particle-number-uniform fermionic sector inequality and a bosonic full-Fock $1/3$ inequality with a bounded two-sided resolvent in $(0,1/3)$. The zero-mode counts, dedicated thermodynamic corollaries, explicit fermionic full-Fock spectral wrapper, and stronger bosonic constants are not covered by that selected acceptance. The mathematical arguments remain in the text; Appendix~\ref{app:reproducibility} states the verification scope and remaining review obligations. No completed whole-paper formalization or human peer review is claimed.}
{已有两条选定的Lean终点命题取得内核验收记录：对粒子数一致的费米扇区不等式，以及玻色全Fock空间的$1/3$不等式和$(0,1/3)$内有界双侧预解算子。零模计数、专门热力学推论、显式费米全Fock谱包装及更强玻色常数不在此次选定验收覆盖内。正文保留这些数学论证；附录~\ref{app:reproducibility}列明验证范围和待复核事项。本文不声称整篇形式化或人工同行评议已经完成。}
\par\endgroup
\begin{abstract}
\bi{We develop a magnetic-periodization transfer of finite Fock-space gap certificates to the full contact interactions on a flat torus. In the normalization with planar nonzero pair energy one, the bosonic $V_0$ Hamiltonian obeys $H_B^2\ge H_B/3$ whenever the shortest lattice vector is at least $20$ in units of $\sqrt2\ell_B$. The fermionic $V_1$ Hamiltonian obeys $H_F^2\ge\gamma H_F$ for every $\gamma<0.04616733319001$ once that shortest vector is sufficiently large. Both estimates hold on the entire corresponding Fock space, uniformly in particle number and boundary twists. The geometric argument retains all center-of-mass theta modes, controls the failure of finite-torus pair projectivity relative to the interaction energy, and periodizes the connected three-body coefficient through an absolutely convergent COM partial trace. Positive finite four-body certificates survive without an inverse Gram matrix. Fermionic errors are controlled by Pauli bounds, while bosonic errors require coherent-state occupation and thermal-kernel comparisons with a positive four-body reserve. At flux $M=2N$ or $M=3N$, the bounds apply above the twofold or threefold Laughlin zero-mode space. The results concern large actual cycles and the full untruncated interactions; they do not assert a shape-uniform bound for arbitrary thin tori.}
{本文把有限Fock空间谱隙证书通过磁周期化迁移到平坦环面上的完整接触相互作用。在平面两体非零能量为一的归一化下，只要最短格矢以$\sqrt2\ell_B$为单位计至少为$20$，玻色$V_0$ Hamiltonian便满足$H_B^2\ge H_B/3$。对每个$\gamma<0.04616733319001$，当最短格矢充分大时，费米$V_1$ Hamiltonian满足$H_F^2\ge\gamma H_F$。两种估计均在相应整个Fock空间上成立，对粒子数及边界twist一致。几何论证保留全部质心theta模式，以相互作用能量相对控制有限环面两体算符偏离投影子的误差，并通过绝对收敛的COM偏迹周期化三体连通系数。有限四体正证书无需逆Gram矩阵即可保留。费米误差由Pauli界控制；玻色误差则需要相干态占据数和thermal核比较，并保留正四体余量。在磁通$M=2N$或$M=3N$时，所得下界作用于二重或三重Laughlin零模空间之上。结果针对实际周期足够大及完整未截断相互作用，不声称对任意薄环面都有形状一致的下界。}
\end{abstract}

\ifChinese\tableofcontents\else{\small\tableofcontents}\fi
\clearpage
\section{\bi{Introduction}{引言}}\label{sec:introduction}
\bi{This manuscript is a candidate proof for independent review. Its mathematical conclusions and the narrower scope of selected machine checks are distinguished in Appendix~\ref{app:reproducibility}.}
{本文是供独立复核的候选证明。数学结论与选定机器检查的较窄范围在附录~\ref{app:reproducibility}中分别说明。}
\bi{The Laughlin wavefunctions provide explicit correlated states at fractional quantum Hall fillings, and contact pseudopotentials make them exact zero-energy states~\cite{Laughlin1983,Haldane1983}. A spectral-gap estimate asks a different question: can every orthogonal state be kept a positive energy above the entire zero-mode space as the particle number grows? On a torus this question also involves magnetic boundary conditions and a degenerate center-of-mass ground space. The familiar theta-function construction describes these zero modes~\cite{HaldaneRezayi1985Periodic,GreiterSchnellsThomale2016}, but their existence alone supplies no uniform lower bound on the positive spectrum.}
{Laughlin波函数在分数量子Hall填充处给出明确的关联态，而接触赝势使它们成为精确零能态~\cite{Laughlin1983,Haldane1983}。谱隙估计回答的是另一个问题：当粒子数增长时，能否让每个正交态与整个零模空间保持正的能量间隔？在环面上，这一问题还涉及磁边界条件与简并的质心基态空间。熟知的theta函数构造描述了这些零模~\cite{HaldaneRezayi1985Periodic,GreiterSchnellsThomale2016}，但零模的存在本身并不给出正谱的一致下界。}

\bi{Finite Fock-space certificates provide a useful way to separate local algebra from geometry. The fermionic certificate of OpenAI~\cite{OpenAI2026Fock} yields a uniform inequality for the full spherical interaction and its untruncated planar limit. Later work studies spherical stability~\cite{OpenAI2026Stability}, and Gao gives improved sphere-and-plane bounds through a semidefinite-program search~\cite{Gao2026ImprovedGap}. We deliberately retain the earlier seven-row certificate rather than claim its constant is the strongest available. For bosons we use the explicit 138-row certificate in the candidate proof prepared and curated by Xin Shen~\cite{Shen2026BosonicCandidate}. Only its finite planar coefficient inequalities are inputs here; neither its proposed global theorem nor its partial formalization is assumed. Those finite inequalities, and the fermionic ones, admit direct independent rational verification.}
{有限Fock空间证书把局部代数与几何分离。OpenAI的费米证书~\cite{OpenAI2026Fock}给出完整球面相互作用及其未截断平面极限的一致不等式。后续工作研究球面稳定性~\cite{OpenAI2026Stability}，Gao则通过半正定规划搜索提高球面与平面下界~\cite{Gao2026ImprovedGap}。本文有意保留较早的七行证书，不声称其常数是现有最强结果。玻色情形使用由Xin Shen整理的候选证明~\cite{Shen2026BosonicCandidate}中的明确138行证书。这里只把其有限平面系数不等式作为输入，不假设其候选整体定理或部分形式化已完成。这些有限不等式和费米不等式均可直接以独立有理算术验证。}

\bi{The geometric distinction is essential. The OpenAI source concerns spherical fermionic $V_1$ and its planar limit; Shen's candidate concerns spherical bosonic $V_0$, with a proposed sufficient flux threshold $Q\ge3088$~\cite{OpenAI2026Fock,Shen2026BosonicCandidate}. Our proposed torus thresholds concern the actual shortest lattice cycle and use $M=3N$ or $M=2N$ at Laughlin filling, in place of the spherical shifted relations. Gao's spherical lower bound $1/5$ is an energy bound at fermionic filling $1/3$, and supplies no torus gap theorem~\cite{Gao2026ImprovedGap}.}
{几何差别是必要的。OpenAI来源研究球面费米$V_1$及其平面极限；Shen的候选证明研究球面玻色$V_0$，提出充分磁通阈值$Q\ge3088$~\cite{OpenAI2026Fock,Shen2026BosonicCandidate}。本文拟议的环面阈值针对实际最短格矢，Laughlin填充关系为$M=3N$或$M=2N$，不采用球面的带移位关系。Gao的球面下界$1/5$是费米填充$1/3$处的能量下界，不提供环面谱隙定理~\cite{Gao2026ImprovedGap}。}

\bi{The torus transfer is not a replacement of spherical $SU(2)$ averaging by a formal magnetic-translation average. A finite torus pair interaction is generally not a projector, raw periodized orbitals are not orthogonal, and a finite local approximation does not control all center-of-mass modes. A three-body single-contraction identity must respect its magnetic phases and the order of the COM trace. The bosonic problem adds unbounded occupation numbers: the bounded-annihilator argument available for fermions cannot be reused. These are the specific analytic steps addressed in this manuscript.}
{环面迁移不是把球面$SU(2)$平均形式上换成磁平移平均。有限环面两体相互作用通常不是投影子，原始周期化轨道并不正交，有限局部近似也不能控制全部质心模式。三体单收缩恒等式必须正确保留磁相位与COM偏迹的次序。玻色问题还具有无界占据数，因此不能复用费米情形的有界湮灭算符论证。本文处理的正是这些具体解析步骤。}

\bi{Our common construction periodizes the planar Gaussian jets without renormalizing them. Parameter Landau-level integration then preserves relative matrix units, while a relative-jet identity controls the Gram operator on the entire output theta space. Positive finite certificates survive by congruence even when their Gram matrices have kernels. The resulting full-Fock estimates are assembled by two distinct error budgets: fermionic errors are paid for by the pair energy, whereas bosonic coherent-state comparisons pay for them jointly with the pair energy and a positive four-body term. Theorems~\ref{thm:boson} and~\ref{thm:fermion} state the resulting bounds.}
{共同构造是在不重新归一化的前提下周期化平面高斯jets。参数Landau能级积分保持相对matrix units，而相对jet恒等式控制整个输出theta空间上的Gram算符。有限正证书通过congruence保留，即使其Gram矩阵有核也如此。最终全Fock估计采用两套不同误差预算：费米误差由pair能量支付，玻色则通过相干态比较，由pair能量与正四体项共同支付。所得下界由定理~\ref{thm:boson}与~\ref{thm:fermion}陈述。}

\bi{The geometric condition is growth of the actual shortest noncontractible cycle at fixed magnetic length. Existing rigorous gap analyses for truncated thin-torus models~\cite{NachtergaeleWarzelYoung2021,WarzelYoung2022ThinTorus,WarzelYoung2023BulkEdge} concern different Hamiltonians or limiting regimes and are not used as gap inputs for the full interactions here. Pairing and coherent-state descriptions of the torus~\cite{OrtizNussinovDukelskySeidel2013,Fremling2013Coherent} inform the setting, but the all-COM frame estimate and the relative error bounds are proved below. All essential kernel, convergence, and occupation arguments are included in the manuscript. The finite exact datasets and executable checks accompany it separately from the finite-size numerical diagnostics.}
{几何条件是在固定磁长下实际最短非收缩周期增长。已有关于截断薄环面模型的严格谱隙分析~\cite{NachtergaeleWarzelYoung2021,WarzelYoung2022ThinTorus,WarzelYoung2023BulkEdge}涉及不同Hamiltonian或极限范围，不作为这里完整相互作用的谱隙输入。环面的pairing与相干态描述~\cite{OrtizNussinovDukelskySeidel2013,Fremling2013Coherent}提供背景，而全COM frame估计与相对误差界在下文证明。所有关键核、收敛及占据数论证都纳入正文和证明附录。有限精确数据及可执行核查随稿提供，并与有限尺寸数值诊断分开。}

\section{\bi{Models and main statements}{模型与主要结论}}\label{sec:models}
\bi{We consider the full lowest-Landau-level contact pseudopotential $V_0$ for bosons and $V_1$ for fermions on a flat magnetic torus. No hopping range, momentum sector, or COM sector is truncated. The energy convention is that the corresponding nonzero planar two-particle pseudopotential eigenvalue is one. Finite-torus pair eigenvalues are not reset to one.}
{本文研究平坦磁环面上玻色子的完整最低Landau能级接触赝势$V_0$及费米子的完整$V_1$。不截断跃迁范围、动量扇区或COM扇区。能量约定是对应平面两粒子赝势的非零本征值为一；不把有限环面两体本征值重新设为一。}

\bi{Let $T=\mathbb C/\Lambda$, $\operatorname{area}(\Lambda)=\pi M$, $M\in\mathbb N$, in the dimensionless coordinate $z=x/(\sqrt2\ell_B)$. Boundary twists are arbitrary and identical for all particles. The one-particle space $\mathcal H_1(T)$ has dimension $M$. The $N$-particle spaces are $\Sym^N\mathcal H_1(T)$ and $\bigwedge^N\mathcal H_1(T)$. The raw periodized orbitals $u_j^T(R)$ and the measure $\dd\mu_k(R)=k\dd^2R/\pi$ are fixed in Section~\ref{sec:geometry}. These orbitals retain their planar scale; they are neither normalized pointwise nor orthogonalized on the torus.}
{令$T=\mathbb C/\Lambda$，在无量纲坐标$z=x/(\sqrt2\ell_B)$中有$\operatorname{area}(\Lambda)=\pi M$、$M\in\mathbb N$。边界twist任意，且所有粒子相同。一粒子空间$\mathcal H_1(T)$维数为$M$。$N$粒子空间分别为$\Sym^N\mathcal H_1(T)$及$\bigwedge^N\mathcal H_1(T)$。原始周期化轨道$u_j^T(R)$与测度$\dd\mu_k(R)=k\dd^2R/\pi$在第~\ref{sec:geometry}节固定。这些轨道保持平面尺度，不作逐点归一化，也不在环面上正交化。}

\bi{Write $a_j(R)=a(u_j^T(R))$ for bosonic annihilation and $c_j(R)=c(u_j^T(R))$ for fermionic annihilation. The COM-descendant pair annihilators are}
{以$a_j(R)=a(u_j^T(R))$表示玻色湮灭，以$c_j(R)=c(u_j^T(R))$表示费米湮灭。Pair的COM后裔湮灭算符为}
\begin{align}
 B_p^B(R)&=\frac1{\sqrt2}\sum_{i+j=p}\sqrt{\frac{p!}{2^pi!j!}}\,a_i(R)a_j(R),\label{eq:B-pair}\\
 v_p^{F,T}(R)&=\sum_{\substack{i<j\\i+j=p+1}}(i-j)\sqrt{\frac{p!}{2^pi!j!}}\,u_i^T(R)\wedge u_j^T(R),
 \qquad B_p^F(R)=A(v_p^{F,T}(R)).\label{eq:F-pair}
\end{align}
\bi{In the bosonic case $v_p^{B,T}=B_p^{B\dagger}|\mathrm{vac}\rangle$. In particular, $B_0^B=a_0^2/\sqrt2$ and $v_0^{F,T}=-u_0^T\wedge u_1^T$. On the plane both $v_p$ families are orthonormal bases of the pair subspaces of relative degree zero and one, respectively. For $\sigma=B,F$ define}
{玻色情形中$v_p^{B,T}=B_p^{B\dagger}|\mathrm{vac}\rangle$。特别地，$B_0^B=a_0^2/\sqrt2$，$v_0^{F,T}=-u_0^T\wedge u_1^T$。在平面上，两族$v_p$分别是相对次数零和一的pair子空间的正交归一基。对$\sigma=B,F$，定义}
\begin{equation}
 h_\sigma=\int_T|v_0^{\sigma,T}(R)\rangle\langle v_0^{\sigma,T}(R)|\dd\mu_2(R),\qquad
 H_\sigma=\mathcal L_2(h_\sigma)=\int_T B_0^\sigma(R)^\dagger B_0^\sigma(R)\dd\mu_2(R).\label{eq:physical-H}
\end{equation}
\bi{This is the periodized contact definition of the full interactions in the stated normalization. The positive lift $\mathcal L_2$ is defined below. All operators are bounded on each fixed $M,N$ sector. The bosonic full Fock space is an infinite direct sum, so its statements first hold on the finite-particle algebraic core and then as spectral inequalities for the direct-sum self-adjoint operator. In particular, $H^2\ge\gamma H$ means $\|H\psi\|^2\ge\gamma\langle\psi,H\psi\rangle$ on the form domain of $H^2$.}
{这是所述归一化下完整相互作用的周期化接触定义。正提升$\mathcal L_2$在下文定义。在固定$M,N$扇区，所有算符均有界。玻色全Fock空间是无限直和，故先在有限粒子代数core上陈述不等式，再通过直和自伴算符的谱演算延拓。特别地，$H^2\ge\gamma H$表示在$H^2$的二次型定义域上有$\|H\psi\|^2\ge\gamma\langle\psi,H\psi\rangle$。}

\begin{theorem}[\bi{Bosonic torus gap}{玻色环面谱隙}]\label{thm:boson}
\bi{For every lattice with $\ell(\Lambda)\ge20$ and every boundary twist, the Hamiltonian~\eqref{eq:physical-H} on the entire bosonic Fock space satisfies}
{对所有满足$\ell(\Lambda)\ge20$的格子及任意边界twist，式~\eqref{eq:physical-H}的Hamiltonian在整个玻色Fock空间上满足}
\begin{equation}
 H_B^2\ge\frac13H_B.
\end{equation}
\bi{More generally, every $0<\gamma<\gamma_B^*$ admits a systole threshold independent of particle number and twist, where}
{更一般地，对每个$0<\gamma<\gamma_B^*$，都存在与粒子数及twist无关的最短周期阈值，其中}
\begin{align}
 \gamma_B^*&=1-\frac{764191703750823432978457}{1500000000000000000000000}
 -\frac{19531250}{129140163}\notag\\
 &\hspace{8mm}-10^{-6}\left[12(9/8)^9+272(9/8)^{17}\right]
 \simeq0.33724904107059384.
\end{align}
\bi{The length threshold $20$ is dimensionless, equivalently $20\sqrt2\ell_B$ in physical units. It is a conservative sufficient condition, not an optimized finite-size estimate.}
{长度阈值$20$是无量纲的，等价于物理长度$20\sqrt2\ell_B$。它是保守的充分条件，不是优化后的有限尺寸估计。}
\end{theorem}

\begin{theorem}[\bi{Fermionic torus gap}{费米环面谱隙}]\label{thm:fermion}
\bi{For every $0<\gamma<\gamma_F^*$ there is a finite $L_\gamma$ such that, for every lattice with $\ell(\Lambda)\ge L_\gamma$ and every boundary twist, the full fermionic Fock-space Hamiltonian satisfies}
{对每个$0<\gamma<\gamma_F^*$，存在有限$L_\gamma$，使对所有满足$\ell(\Lambda)\ge L_\gamma$的格子及任意边界twist，完整费米Fock空间Hamiltonian满足}
\begin{equation}
 H_F^2\ge\gamma H_F,\qquad
 \gamma_F^*=\frac{4616733319001}{10^{14}}=0.04616733319001.
\end{equation}
\bi{The threshold is independent of particle number and twist. No explicit numerical value of $L_\gamma$ is asserted here.}
{阈值与粒子数及twist无关。本文不声称得到$L_\gamma$的显式数值。}
\end{theorem}

\begin{corollary}[\bi{Laughlin fillings}{Laughlin填充} ]\label{cor:fillings}
\bi{At $M=2N$ for bosons and $M=3N$ for fermions, the zero-energy spaces have dimensions two and three, respectively. The preceding inequalities bound the spectrum above these whole ground spaces. In any compact family of torus shapes at fixed magnetic length, they give thermodynamic gap bounds $1/3$ and $1/25$, respectively, for sufficiently large $N$. On the bosonic square sequence $M=2N$, the explicit condition is already satisfied for $N\ge64$.}
{玻色在$M=2N$、费米在$M=3N$时，零能空间维数分别为二和三。上述不等式控制整个基态空间之上的谱。在固定磁长、环面形状属于紧集的任意族中，它们对充分大$N$分别给出$1/3$及$1/25$的热力学谱隙下界。对玻色方形序列$M=2N$，显式条件在$N\ge64$时已满足。}
\end{corollary}
\bi{The kernel assertion is proved in Section~\ref{sec:groundspaces}. The fermionic estimate is assembled in Section~\ref{sec:F-fermionic} and the bosonic estimate in Section~\ref{sec:B-bosonic}. Neither theorem includes a fixed short circumference or arbitrary degenerating thin-torus shapes. Moreover, on a sector with no zero mode, a full-Fock inequality excludes positive spectrum in $(0,\gamma)$; it does not assert a spacing $\gamma$ between the two lowest positive eigenvalues.}
{零空间结论在第~\ref{sec:groundspaces}节证明。费米估计在第~\ref{sec:F-fermionic}节汇总，玻色估计在第~\ref{sec:B-bosonic}节汇总。两定理都不覆盖固定短周长或任意退化的薄环面形状。此外，在没有零模的扇区，全Fock不等式排除$(0,\gamma)$中的正谱，并不声称最低两个正本征值之间相隔$\gamma$。}

\section{\bi{Magnetic periodization and local oscillator identities}{磁周期化与局部振子恒等式}}\label{sec:geometry}
\bi{Torus coherent states and their reproducing kernels provide the geometric background~\cite{Fremling2013Coherent}. We specify the magnetic phases, measures, and raw-jet normalization below, and derive the full-COM frame estimate rather than assume it from a localized coherent-state construction.}
{环面相干态及其再生核提供几何背景~\cite{Fremling2013Coherent}。下文明确磁相位、测度和原始jet的归一化，并推导全COM框架估计；这一估计不作为局域相干态构造的现成结论使用。}
\subsection{\bi{Line bundles, measures, and the parameter Landau levels}{线丛、测度与参数Landau能级}}
\bi{Let $\Lambda=\mathbb Z\omega_1+\mathbb Z\omega_2\subset\mathbb C$, with $\operatorname{Im}(\overline\omega_1\omega_2)=\pi M$, and write $T=\mathbb C/\Lambda$. Lengths are measured in the coordinate $z=x/(\sqrt2\ell_B)$, where $\ell_B$ is the magnetic length. The systole is $\ell(\Lambda)=\min_{\lambda\in\Lambda\setminus\{0\}}|\lambda|$. All limits below concern this actual shortest lattice vector. Fix arbitrary real boundary angles $\theta_1,\theta_2$, and set}
{令$\Lambda=\mathbb Z\omega_1+\mathbb Z\omega_2\subset\mathbb C$，其中$\operatorname{Im}(\overline\omega_1\omega_2)=\pi M$，并记$T=\mathbb C/\Lambda$。长度使用坐标$z=x/(\sqrt2\ell_B)$计量，$\ell_B$为磁长。定义最短周期$\ell(\Lambda)=\min_{\lambda\in\Lambda\setminus\{0\}}|\lambda|$。以下所有几何极限都针对这一实际最短格矢。固定任意实边界相位$\theta_1,\theta_2$，定义}
\begin{equation}
 \chi(m\omega_1+n\omega_2)=e^{i(m\theta_1+n\theta_2)}(-1)^{Mmn},\qquad \chi_k=\chi^k.
\end{equation}
\bi{For positive integer charge $k$, use the measure $\dd\mu_k(z)=k\dd^2z/\pi$ and magnetic translations}
{对正整数电荷$k$，使用测度$\dd\mu_k(z)=k\dd^2z/\pi$及磁平移}
\begin{align}
 (T_x^{(k)}f)(z)&=\exp\!\left[\frac{k}{2}(z\bar x-\bar z x)\right]f(z-x),\label{eq:translation}\\
 T_x^{(k)}T_y^{(k)}&=e^{-ik\operatorname{Im}(x\bar y)}T_{x+y}^{(k)},\qquad
 \chi_k(\lambda)\chi_k(\mu)=e^{-ik\operatorname{Im}(\lambda\bar\mu)}\chi_k(\lambda+\mu).
\end{align}
\bi{The charge-$k$ lowest Landau level (LLL) $\mathcal H_k(T)$ consists of Gaussian-holomorphic sections satisfying $T_\lambda^{(k)}f=\chi_k(\lambda)f$. Equivalently,}
{电荷$k$的最低Landau能级（LLL）$\mathcal H_k(T)$由满足$T_\lambda^{(k)}f=\chi_k(\lambda)f$的高斯全纯截面构成。等价地，}
\begin{equation}
 f(z+\lambda)=\chi_k(\lambda)^{-1}
 e^{k(z\bar\lambda-\bar z\lambda)/2}f(z).
\end{equation}
\bi{Its dimension is $kM$. The plane reproducing kernel in the same measure is $K_k(z,w)=\exp[k(z\bar w-|z|^2/2-|w|^2/2)]$. For a Gaussian times a polynomial define}
{其维数为$kM$。同一测度下的平面再生核为$K_k(z,w)=\exp[k(z\bar w-|z|^2/2-|w|^2/2)]$。对高斯函数乘多项式，定义}
\begin{equation}
 P_k f=\sum_{\lambda\in\Lambda}\chi_k(\lambda)^{-1}T_\lambda^{(k)}f,
 \qquad K_k^T=P_k K_k.\label{eq:periodization}
\end{equation}
\bi{These series and any fixed number of their derivatives converge absolutely. Reindexing gives the stated boundary law. Unfolding a fundamental-cell integral proves the reproducing property of $K_k^T$ with no additional area normalization. We do not assert that $P_k$ is a bounded map on the entire infinite-dimensional planar LLL. Its use is restricted to finite-degree vectors and the weighted, absolutely convergent series justified below.}
{这些级数及任意固定阶导数绝对收敛。重排指标得到上述边界规律，展开基本胞积分则证明$K_k^T$的再生性质，无需额外面积归一化。这里不声称$P_k$在整个无限维平面LLL上是有界映射；它只作用于有限次数向量及下文证明绝对收敛的带权级数。}

\bi{Let $e_j(z)=z^j e^{-|z|^2/2}/\sqrt{j!}$, $u_j(R)=T_R^{(1)}e_j$, and $u_j^T(R)=P_1u_j(R)$. A common-image unfolding gives, for finite Gaussian polynomials $f,g$,}
{令$e_j(z)=z^j e^{-|z|^2/2}/\sqrt{j!}$、$u_j(R)=T_R^{(1)}e_j$、$u_j^T(R)=P_1u_j(R)$。对有限高斯多项式$f,g$，展开共同image可得}
\begin{equation}
 \langle P_1f,P_1g\rangle_T
 =\sum_{\lambda\in\Lambda}\chi(\lambda)^{-1}
   \langle f,T_\lambda^{(1)}g\rangle_{\mathbb C}.\label{eq:inner-unfold}
\end{equation}
\bi{Indeed, write the two images as $\mu$ and $\mu+\lambda$. The cocycle and semicharacter factors cancel the common translation, and the translated cells tile the plane. Absolute Gaussian convergence permits this rearrangement. In particular, putting $b_\lambda(R)=\chi(\lambda)e^{i\operatorname{Im}(\lambda\bar R)}$,}
{具体地，把两个image写成$\mu$及$\mu+\lambda$；cocycle与semicharacter因子消去共同平移，而平移后的基本胞铺满平面。高斯绝对收敛保证这种重排合法。特别地，记$b_\lambda(R)=\chi(\lambda)e^{i\operatorname{Im}(\lambda\bar R)}$，则}
\begin{equation}
 u_j^T(R+\lambda)=b_\lambda(R)u_j^T(R),\qquad
 g^T_{ij}(R,S)=\sum_{\lambda\in\Lambda}b_\lambda(S)^{-1}g^{\rm pl}_{ij}(R,S+\lambda),\label{eq:gram-unfold}
\end{equation}
\bi{where $g_{ij}(R,S)=\langle u_i(R),u_j(S)\rangle$. The phase in the second identity follows from $T_\lambda u_j(S)=e^{-i\operatorname{Im}(\lambda\bar S)}u_j(S+\lambda)$.}
{其中$g_{ij}(R,S)=\langle u_i(R),u_j(S)\rangle$。第二个恒等式的相位来自$T_\lambda u_j(S)=e^{-i\operatorname{Im}(\lambda\bar S)}u_j(S+\lambda)$。}

\bi{For every fixed $R$, all $u_j^T(R)$ are physical one-particle LLL vectors in the variable $z$. Their dependence on the center $R$ is different: $u_j^T$ lies in the $j$th Landau level of the dual parameter line bundle. The annihilation amplitude is its conjugate and has the holomorphic-side automorphy, but is in the parameter LLL only for $j=0$. More generally, write $\mathcal H_{k,n}(T)$ for parameter charge $k$ and Landau level $n$, with $\mathcal H_k=\mathcal H_{k,0}$. A $k$-body center-of-mass (COM) vacuum seed has an amplitude in $\mathcal H_k$; its $n$th COM descendant has an amplitude in $\mathcal H_{k,n}$. A shared boundary multiplier alone is not the LLL differential condition.}
{对每个固定$R$，所有$u_j^T(R)$都是物理变量$z$的一粒子LLL向量。但它们对中心$R$的依赖不同：$u_j^T$位于对偶参数线丛的第$j$个Landau能级。湮灭振幅取其共轭，具有全纯侧的边界规律，但仅在$j=0$时属于参数LLL。一般记$\mathcal H_{k,n}(T)$为电荷$k$、第$n$级的参数空间，并令$\mathcal H_k=\mathcal H_{k,0}$。$k$体质心（COM）真空种子的振幅属于$\mathcal H_k$，其第$n$个COM后裔的振幅属于$\mathcal H_{k,n}$。共同边界乘子本身并不等于LLL微分条件。}

\subsection{\bi{The positive lift and exact COM integration}{正提升与精确COM积分}}
\bi{The positive normal-ordered lift is the Fock-space device used for the fermionic comparison in OpenAI~\cite[Section~2]{OpenAI2026Fock} and for the bosonic comparison in Shen~\cite[Section~1]{Shen2026BosonicCandidate}. We retain its vacuum normalization. The parameter Landau-level identities that follow supply the torus replacement for the sources' spherical averaging step.}
{正的正规序提升沿用OpenAI费米比较~\cite[第2节]{OpenAI2026Fock}及Shen玻色比较~\cite[第1节]{Shen2026BosonicCandidate}中的Fock空间方法，并保留其真空归一化。后续参数Landau能级恒等式给出环面上替代原方案球面平均步骤的构造。}
\bi{In either statistics, let $\mathcal L_k(|v\rangle\langle w|)=A(v)^\dagger A(w)$, where $A(v)^\dagger|\mathrm{vac}\rangle=v$ with normalized symmetric or antisymmetric tensors. This linear map preserves positivity, equals the input operator in the $k$-particle sector, and vanishes below that sector. For bosonic occupation vectors, $A(\boldsymbol\nu)=\prod_j a_j^{\nu_j}/\sqrt{\nu_j!}$; thus no extra $k!$ occurs. All few-body inequalities are lifted by this map, rather than estimated by counting particle subsets.}
{对两种统计，定义$\mathcal L_k(|v\rangle\langle w|)=A(v)^\dagger A(w)$，其中$A(v)^\dagger|\mathrm{vac}\rangle=v$，对称或反对称张量均取规范化约定。该线性映射保正，在$k$粒子扇区等于输入算符，在更低粒子数扇区为零。玻色占据向量满足$A(\boldsymbol\nu)=\prod_j a_j^{\nu_j}/\sqrt{\nu_j!}$，因此没有额外$k!$。所有少体不等式都通过这一映射提升，不通过粒子子集数估计。}

\bi{On dual parameter sections define}
{在对偶参数截面上定义}
\begin{equation}
 \mathfrak a_k=-\frac{\partial_R+k\bar R/2}{\sqrt k},\qquad
 \mathfrak a_k^\dagger=\frac{\partial_{\bar R}-kR/2}{\sqrt k},\qquad
 [\mathfrak a_k,\mathfrak a_k^\dagger]=1.\label{eq:parameter-ladder}
\end{equation}
\bi{They preserve the boundary law and are adjoints in the cell integral. Differentiating the Gaussian series gives}
{它们保持边界规律，并在基本胞积分中互为伴随。逐项微分高斯级数得到}
\begin{equation}
 \mathfrak a_1u_j^T=\sqrt j\,u_{j-1}^T,\qquad
 \mathfrak a_1^\dagger u_j^T=\sqrt{j+1}\,u_{j+1}^T.
\end{equation}
\bi{These canonical commutation relations act on the infinite parameter Landau-level space, not inside the finite physical LLL.}
{这些正则对易关系作用于无限维参数Landau能级空间，不是在有限维物理LLL内实现。}

\begin{lemma}[\bi{COM integration}{COM积分}]\label{lem:ladder}
\bi{Let $w,z$ be finite-degree $k$-body COM vacuum seeds of the chosen statistics. Define}
{设$w,z$为所选统计下有限次数的$k$体COM真空种子。定义}
\begin{equation}
 w_n^T=(\mathfrak a_k^\dagger)^nw^T/\sqrt{n!},\qquad
 z_m^T=(\mathfrak a_k^\dagger)^mz^T/\sqrt{m!}.
\end{equation}
\bi{Then}{则}
\begin{equation}
 \int_T|w_n^T(R)\rangle\langle z_m^T(R)|\dd\mu_k(R)
 =\delta_{nm}\int_T|w^T(R)\rangle\langle z^T(R)|\dd\mu_k(R).\label{eq:ladder-integral}
\end{equation}
\bi{Different relative seeds need not be orthogonal; equal COM levels retain their full relative matrix unit.}
{不同相对种子无需正交；相同COM级数保留它们之间完整的相对matrix unit。}
\end{lemma}
\begin{proof}
\bi{Evaluate the operator identity between fixed physical bras and kets. Integration by parts on the parameter line bundle has no boundary term, because the two boundary multipliers cancel. Move lowering operators across the integral, use $[\mathfrak a_k,(\mathfrak a_k^\dagger)^n]=n(\mathfrak a_k^\dagger)^{n-1}$, and annihilate the vacuum seeds. Unequal levels vanish; equal levels contribute $n!$, cancelled by the two normalizations. Polarization gives the identity for distinct seeds.}
{在固定物理bra与ket之间检验该算符恒等式。两个边界乘子相消，因此参数线丛上的积分分部没有边界项。把降低算符移过积分，使用$[\mathfrak a_k,(\mathfrak a_k^\dagger)^n]=n(\mathfrak a_k^\dagger)^{n-1}$并湮灭真空种子。不等的级数给零，相等的级数给$n!$，恰被两侧规范化抵消。极化即给出不同种子之间的恒等式。}
\end{proof}

\bi{Let $C_R e_j=u_j^T(R)$ and let $C_R^{(k)}$ denote its symmetric or exterior $k$th tensor action. If $X$ is an operator on the physical relative space, define}
{令$C_Re_j=u_j^T(R)$，$C_R^{(k)}$表示其第$k$次对称张量或外幂作用。对物理相对空间上的算符$X$，定义}
\begin{equation}
 \Omega_k(X)=\int_T C_R^{(k)}
 \bigl(X\otimes|0_{\CM}\rangle\langle0_{\CM}|\bigr)
 C_R^{(k)\dagger}\dd\mu_k(R).\label{eq:Omega}
\end{equation}
\bi{For finite support this is a positive map, including when the raw-jet Gram matrix has a kernel. For every finite-degree anchored operator $A$, Lemma~\ref{lem:ladder} gives}
{对有限支撑，这是保正映射，即使原始jet的Gram矩阵有核也成立。对任意有限次数anchored算符$A$，引理~\ref{lem:ladder}给出}
\begin{equation}
 \int_T C_R^{(k)}A C_R^{(k)\dagger}\dd\mu_2(R)
 =\Omega_k\!\left(\frac2k\Tr_{\CM}A\right).\label{eq:torus-partial-trace}
\end{equation}
\bi{The corresponding plane twirl is}
{对应的平面twirl为}
\begin{equation}
 \int_{\mathbb C}T_R^{\otimes k}A T_R^{\otimes k\dagger}\dd\mu_2(R)
 =\frac2k I_{\CM}\otimes\Tr_{\CM}A.\label{eq:plane-twirl}
\end{equation}
\bi{To verify the normalization, the normalized COM displacement is $h=\sqrt k R$. Writing $h=(q+ip)/\sqrt2$, Weyl conjugation sends the COM position kernel $A(x,y)$ to $e^{ip(x-y)}A(x-q,y-q)$. Integrating first in $p$ and then in $q$, with measure $\dd q\dd p/(2\pi)$, gives identity in COM and trace in that factor. The change of measure $\dd\mu_2(R)=(2/k)\dd^2h/\pi$ supplies the factor $2/k$. This argument also applies to relative-operator-valued kernels; it does not scalarize multiplicity spaces.}
{为核对归一化，规范化COM位移是$h=\sqrt kR$。写$h=(q+ip)/\sqrt2$，Weyl共轭把COM位置核$A(x,y)$变为$e^{ip(x-y)}A(x-q,y-q)$。按测度$\dd q\dd p/(2\pi)$先积$p$再积$q$，得到COM恒等算符及该因子上的迹。测度变换$\dd\mu_2(R)=(2/k)\dd^2h/\pi$给出$2/k$。该论证同样适用于取值为相对算符的核，不把重数空间误作标量。}

\subsection{\bi{Uniform raw-jet bounds and the convergent domain}{一致原始jet界与收敛定义域}}
\begin{lemma}[\bi{Finite and high jets}{有限阶与高阶jet}]\label{lem:jets}
\bi{For fixed $J$, define}{对固定$J$，定义}
\begin{equation}
 e_J(\Lambda)=\sup_R\|(\langle u_i^T(R),u_j^T(R)\rangle)_{0\le i,j\le J}-I\|.
\end{equation}
\bi{There is a constant $C_J$, independent of flux, twist, particle number, and lattice shape, such that}
{存在与磁通、twist、粒子数及格子形状无关的常数$C_J$，使}
\begin{equation}
 e_J(\Lambda)\le C_J\sum_{\lambda\ne0}(1+|\lambda|)^{2J}e^{-|\lambda|^2/2}
 \longrightarrow0\quad(\ell(\Lambda)\to\infty).\label{eq:finite-jet}
\end{equation}
\bi{For all lattices with $\ell(\Lambda)\ge1$, a universal constant $C$ also satisfies}
{对所有满足$\ell(\Lambda)\ge1$的格子，还存在统一常数$C$使}
\begin{equation}
 \sup_R\|u_j^T(R)\|^2\le C(1+j),\qquad j\ge0.\label{eq:high-jet}
\end{equation}
\end{lemma}
\begin{proof}
\bi{Equation~\eqref{eq:inner-unfold} expresses the finite Gram matrix as its identity image plus oscillator displacement matrices. Each fixed matrix element is a Gaussian times a polynomial of degree at most $i+j$, proving~\eqref{eq:finite-jet} by a finite matrix norm bound. For the diagonal element use $\langle e_j,T_\lambda e_j\rangle=e^{-t/2}L_j(t)$, $t=|\lambda|^2$. Unitarity bounds its modulus by one, whereas the finite Laguerre series gives}
{式~\eqref{eq:inner-unfold}把有限Gram矩阵写成恒等image加上振子位移矩阵。每个固定矩阵元都是高斯函数乘次数至多$i+j$的多项式，有限矩阵范数估计即证明~\eqref{eq:finite-jet}。对角元使用$\langle e_j,T_\lambda e_j\rangle=e^{-t/2}L_j(t)$，其中$t=|\lambda|^2$。酉性给出其模至多为一，而有限Laguerre级数给出}
\begin{equation}
 |L_j(t)|\le\sum_{r=0}^j\frac{(jt)^r}{(r!)^2}\le e^{2\sqrt{jt}}.
\end{equation}
\bi{For $t\ge64j$, the product with $e^{-t/2}$ is at most $e^{-t/4}$. The points with $|\lambda|<8\sqrt j$ number at most $C(1+j)$ by disjoint discs of radius $1/2$; the exterior Gaussian sum is uniformly bounded. The same packing argument bounds every fixed polynomial-Gaussian lattice sum uniformly in shape and makes it vanish as the systole grows.}
{在$t\ge64j$时，与$e^{-t/2}$相乘后的上界为$e^{-t/4}$。通过半径$1/2$的不交圆盘，$|\lambda|<8\sqrt j$内的格点数至多为$C(1+j)$；外部高斯和一致有界。同一packing论证对每个固定多项式乘高斯的格点和给出与形状无关的控制，并使其随最短周期增大而消失。}
\end{proof}
\bi{On the span of $e_0,\ldots,e_L$, the squared norm of $C_R$ is at most the sum of the squared column norms, hence at most $C(L+1)^2$. Its $k$th tensor action has squared norm at most $C_k(L+1)^{2k}$. Consequently all anchored degree blocks with trace norms bounded by a polynomial times $\rho^L$, $0<\rho<1$, remain absolutely summable after periodization. Equations~\eqref{eq:torus-partial-trace} and~\eqref{eq:plane-twirl} extend to precisely the exponentially weighted seeds used below by finite-degree truncation. This convergence is first established at fixed flux and particle sector; particle-number uniformity of the final gap comes from the relative energy bounds, not from a coarse many-body norm estimate.}
{在$e_0,\ldots,e_L$张成的空间中，$C_R$的范数平方至多为列范数平方之和，因此不超过$C(L+1)^2$。其第$k$次张量作用的范数平方不超过$C_k(L+1)^{2k}$。于是，迹范数被多项式乘$\rho^L$（$0<\rho<1$）控制的anchored次数块在周期化后仍绝对可和。通过有限次数截断，式~\eqref{eq:torus-partial-trace}与~\eqref{eq:plane-twirl}延拓到下文实际使用的指数带权种子。先在固定磁通和粒子扇区建立这一收敛；最终谱隙对粒子数的一致性来自相对能量界，不来自粗略的多体范数估计。}

\subsection{\bi{A relative jet frame retaining every COM mode}{保留全部COM模式的相对jet frame}}
\begin{lemma}[\bi{Full-COM relative frame}{全COM相对frame}]\label{lem:relative-frame}
\bi{Let $a,b$ be positive integers, $t=a+b$, and $\nu=ab/t$. On $\mathcal H_a(T)\otimes\mathcal H_b(T)$ define the covariant normal jet along the diagonal by}
{设$a,b$为正整数，$t=a+b$、$\nu=ab/t$。在$\mathcal H_a(T)\otimes\mathcal H_b(T)$上，沿对角线定义协变法向jet}
\begin{equation}
 J_dF(C)=(\nu^dd!)^{-1/2}
 \left.\partial_r^d\left\{e^{\nu|r|^2/2}
 F(C+br/t,C-ar/t)\right\}\right|_{r=0}.\label{eq:J-definition}
\end{equation}
\bi{For a fixed finite set $\mathcal D\subset\{0,\ldots,D\}$ the output is $\bigoplus_{d\in\mathcal D}\mathcal H_t(T)$, retaining all $tM$ COM modes in each summand. Its Gram operator satisfies the exact identity}
{对固定有限集合$\mathcal D\subset\{0,\ldots,D\}$，输出为$\bigoplus_{d\in\mathcal D}\mathcal H_t(T)$，每个直和项保留全部$tM$个COM模式。其Gram算符满足精确恒等式}
\begin{equation}
 J_dJ_e^\dagger
 =\sum_{\delta\in\Lambda}\chi_a(\delta)^{-1}
 M_{de}^{(\nu)}(\delta)\,T_{a\delta/t}^{(t)},\qquad
 M_{de}^{(\nu)}(\delta)=\langle d|T_\delta^{(\nu)}|e\rangle.\label{eq:exact-JJ}
\end{equation}
\bi{Every output translation in this formula is unitary on the same fixed-twist space. In particular,}
{该式中的每个输出平移都在同一固定twist空间上为酉算符。特别地，}
\begin{equation}
 \|JJ^\dagger-I\|\le\epsilon_{a,b,D}(\Lambda)
 \le C_{a,b,D}\sum_{\delta\ne0}(1+|\delta|)^{2D}e^{-\nu|\delta|^2/2}\longrightarrow0,
 \qquad J^\dagger J\le(1+\epsilon_{a,b,D})I.\label{eq:J-bound}
\end{equation}
\end{lemma}
\begin{proof}
\bi{The weighted Jacobi change $C=(aR+bS)/t$, $r=R-S$ gives}
{带权Jacobi变换$C=(aR+bS)/t$、$r=R-S$给出}
\begin{equation}
 a|R|^2+b|S|^2=t|C|^2+\nu|r|^2,\qquad
 \dd\mu_a(R)\dd\mu_b(S)=\dd\mu_t(C)\dd\mu_\nu(r).
\end{equation}
\bi{On the plane this is a unitary tensor factorization. In the normalized relative oscillator basis $e_d^{(\nu)}(r)=(\sqrt\nu r)^d e^{-\nu|r|^2/2}/\sqrt{d!}$, $J_d$ extracts the $d$th coefficient, so $J_dJ_e^\dagger=\delta_{de}I$. This plane calculation defines the local diagonal jet on the torus; no global COM-times-relative torus factorization is assumed.}
{在平面上，这是酉的张量分解。相对振子的规范化基为$e_d^{(\nu)}(r)=(\sqrt\nu r)^d e^{-\nu|r|^2/2}/\sqrt{d!}$，$J_d$提取第$d$个系数，故$J_dJ_e^\dagger=\delta_{de}I$。这一平面计算定义了环面对角线上的局部jet，并未假设环面全局分解成独立的COM与相对环面。}

\bi{Expand both input reproducing kernels using~\eqref{eq:periodization}. Under the Jacobi unitary $U$, their images obey}
{利用~\eqref{eq:periodization}展开两个输入再生核。在Jacobi酉变换$U$下，它们的image满足}
\begin{equation}
 U(T_\lambda^{(a)}\otimes T_\mu^{(b)})U^\dagger
 =T_{(a\lambda+b\mu)/t}^{(t)}\otimes T_{\lambda-\mu}^{(\nu)}.
\end{equation}
\bi{Set $\delta=\lambda-\mu$ and $\eta=a\delta/t$. There is no extra phase in this factorization: the magnetic bilinear form and Gaussian quadratic form split by the same change of variables. The remaining cocycles cancel explicitly:}
{令$\delta=\lambda-\mu$及$\eta=a\delta/t$。这一分解没有额外相位，因为磁双线性型与高斯二次型按同一变量变换分解。余下的cocycle逐式相消：}
\begin{align}
 \chi_a(\mu+\delta)^{-1}\chi_b(\mu)^{-1}
 &=\chi_t(\mu)^{-1}\chi_a(\delta)^{-1}e^{-ia\operatorname{Im}(\mu\bar\delta)},\\
 T_{\mu+\eta}^{(t)}&=e^{ia\operatorname{Im}(\mu\bar\delta)}T_\mu^{(t)}T_\eta^{(t)}.
\end{align}
\bi{Summing the common image $\mu$ recovers the charge-$t$ torus kernel. Taking relative oscillator matrix elements yields~\eqref{eq:exact-JJ}. For any $\gamma\in\Lambda$, the commutator of $T_\gamma^{(t)}$ and $T_\eta^{(t)}$ has phase $e^{-2ia\operatorname{Im}(\gamma\bar\delta)}=1$, since $\operatorname{Im}(\gamma\bar\delta)\in\pi M\mathbb Z$ and $aM$ is an integer. Thus the fractional lattice translation preserves the output boundary conditions.}
{对共同image $\mu$求和恢复电荷$t$的环面核，取相对振子矩阵元即得~\eqref{eq:exact-JJ}。对任意$\gamma\in\Lambda$，$T_\gamma^{(t)}$与$T_\eta^{(t)}$的交换相位为$e^{-2ia\operatorname{Im}(\gamma\bar\delta)}=1$，因为$\operatorname{Im}(\gamma\bar\delta)\in\pi M\mathbb Z$且$aM$为整数。因此这一分数格点平移保持输出边界条件。}

\bi{The zero image is $\delta_{de}I$ on the entire output theta space. The oscillator formula bounds each nonzero image by}
{零image是在整个输出theta空间上的$\delta_{de}I$。振子公式对每个非零image给出}
\begin{equation}
 |M_{de}^{(\nu)}(\delta)|\le e^{-\nu|\delta|^2/2}\sqrt{d!e!}
 \sum_{l=0}^{\min(d,e)}\frac{(\sqrt\nu|\delta|)^{d+e-2l}}{l!(d-l)!(e-l)!}.\label{eq:displacement-polynomial}
\end{equation}
\bi{The triangle inequality and a finite matrix norm estimate prove~\eqref{eq:J-bound}. No factor $tM$ occurs. Tensoring with the identity of any spectator Hilbert space leaves the norm unchanged; this is the form used for Fock-valued annihilation amplitudes.}
{逐项三角不等式及有限矩阵范数估计证明~\eqref{eq:J-bound}，不会产生$tM$因子。与任意spectator Hilbert空间上的恒等算符作张量积不改变范数；下文正是将该形式用于取值为Fock空间的湮灭振幅。}
\end{proof}

\subsection{\bi{The pair almost-projector and the common four-body frame}{两体近投影子与共同四体frame}}
\bi{For relative degree $r=0$ in the bosonic case and $r=1$ in the fermionic case, define $Q\psi(R)=\langle v_0^T(R),\psi\rangle$. The target is all of $\mathcal H_2(T)$. On normalized symmetric tensors $Q=J_0$; on normalized antisymmetric tensors $Q=J_1$, because our seed is $v_0^F=-e_0\wedge e_1=(z_1-z_2)/\sqrt2$ and the relative charge is $1/2$. Thus neither case has an additional $\sqrt2$. Restriction to the appropriate statistics does not change $J_rJ_r^\dagger$, since the relative jet already has that parity.}
{玻色情形的相对次数为$r=0$，费米情形为$r=1$。定义$Q\psi(R)=\langle v_0^T(R),\psi\rangle$，其目标是整个$\mathcal H_2(T)$。在规范化对称张量上$Q=J_0$；在规范化反对称张量上$Q=J_1$，因为本文种子为$v_0^F=-e_0\wedge e_1=(z_1-z_2)/\sqrt2$且相对电荷为$1/2$。两种情形均无额外$\sqrt2$。限制到相应统计不改变$J_rJ_r^\dagger$，因为相对jet本身已经具有该奇偶性。}
\begin{equation}
 QQ^\dagger-I=\sum_{\delta\ne0}\chi(\delta)^{-1}
 L_r(|\delta|^2/2)e^{-|\delta|^2/4}T_{\delta/2}^{(2)},\qquad
 \epsilon_2^{(r)}=\sum_{\delta\ne0}|L_r(|\delta|^2/2)|e^{-|\delta|^2/4}.\label{eq:pair-fullCOM}
\end{equation}
\bi{Here $L_0(x)=1$ and $L_1(x)=1-x$. Since $h_T=Q^\dagger Q$, the relative estimate is the exact algebraic consequence}
{这里$L_0(x)=1$、$L_1(x)=1-x$。由于$h_T=Q^\dagger Q$，相对估计是如下精确代数恒等式的推论：}
\begin{equation}
 h_T^2-h_T=Q^\dagger(QQ^\dagger-I)Q,\qquad
 -\epsilon_2h_T\le h_T^2-h_T\le\epsilon_2h_T,\qquad
 \mathcal L_2(h_T^2)\ge(1-\epsilon_2)H.\label{eq:pair-relative}
\end{equation}
\bi{For $\epsilon_2<1$, this also controls every nonzero pair eigenvalue in $[1-\epsilon_2,1+\epsilon_2]$ and proves that $Q$ is onto. It cannot miss a high-COM small singular value.}
{当$\epsilon_2<1$时，它同时把全部两体非零本征值控制在$[1-\epsilon_2,1+\epsilon_2]$中，并证明$Q$满射，不会漏掉高COM处的小奇异值。}

\bi{Both pair families obey $v_p^T=(\mathfrak a_2^\dagger)^pv_0^T/\sqrt{p!}$. Hence Lemma~\ref{lem:ladder} gives $\int_TB_p^\dagger B_q\dd\mu_2=\delta_{pq}H$. The amplitude $B_p$ has parameter level $p$, but $B_0$ is in the parameter LLL. For $F(R,S)=B_0(R)B_0(S)\psi$, apply Lemma~\ref{lem:relative-frame} with $a=b=2$. The relative ladder binomial expansion gives}
{两种pair族均满足$v_p^T=(\mathfrak a_2^\dagger)^pv_0^T/\sqrt{p!}$。因此引理~\ref{lem:ladder}给出$\int_TB_p^\dagger B_q\dd\mu_2=\delta_{pq}H$。振幅$B_p$处于参数level $p$，而$B_0$属于参数LLL。对$F(R,S)=B_0(R)B_0(S)\psi$，以$a=b=2$应用引理~\ref{lem:relative-frame}。相对梯级的二项展开给出}
\begin{equation}
 J_dF=C_d\psi,\qquad C_d=\sum_{p=0}^d(-1)^{d-p}\sqrt{\binom dp}\,2^{-d/2}B_pB_{d-p},\qquad
 T_{4,D}=\int_T\sum_{d=0}^DC_d^\dagger C_d\dd\mu_4.\label{eq:four-frame}
\end{equation}
\bi{Odd $d$ vanish because pair annihilators commute. The tensor-valued frame estimate yields}
{由于pair湮灭算符彼此对易，奇数$d$项消失。取值于张量空间的frame估计给出}
\begin{equation}
 T_{4,D}\le(1+\epsilon_{2,2,D})S_4.\label{eq:four-Bessel}
\end{equation}
\bi{We use $D=22$ for fermions and $D=16$ for bosons. The fermionic proof converts the resulting error into a multiple of $H$ using bounded annihilators; the bosonic proof retains and pays for an $S_4$ contribution using occupation inequalities. These are different steps, and the latter never uses a bounded bosonic annihilator.}
{费米采用$D=22$，玻色采用$D=16$。费米证明利用有界湮灭算符把所得误差转为$H$的倍数；玻色证明保留$S_4$误差，并用占据数不等式支付。两者是不同步骤，后者从不把玻色湮灭算符当成有界算符。}

\section{\bi{The exact three-body connected operator}{精确三体连通算符}}\label{sec:connected}
\bi{The separation by contraction number is Wick normal ordering~\cite{Wick1950}. Its use to reduce the Laughlin gap to two-, three-, and four-body coefficients follows OpenAI~\cite[Section~2]{OpenAI2026Fock} and the bosonic adaptation in Shen~\cite[Section~1]{Shen2026BosonicCandidate}. On a finite torus the double-contraction coefficient must be recomputed from the physical pair operator; it cannot be replaced by the spherical projector coefficient.}
{按收缩数分离各项使用Wick正规序~\cite{Wick1950}。以此把Laughlin谱隙归结为二体、三体和四体系数的方法沿用OpenAI~\cite[第2节]{OpenAI2026Fock}及Shen的玻色推广~\cite[第1节]{Shen2026BosonicCandidate}。有限环面上须从物理两体算符重新计算双收缩系数，不能直接代入球面投影子的系数。}
\bi{Write $\sigma=B,F$ for bosons and fermions. Normal ordering the square of the full Hamiltonian gives}
{用$\sigma=B,F$分别表示玻色与费米。对完整Hamiltonian的平方作正规序，得到}
\begin{equation}
 H_\sigma^2=\mathcal L_2(h_\sigma^2)+S_3^\sigma+S_4^\sigma,\qquad
 S_4^\sigma=\int_{T^2}(B_0(R)B_0(S))^\dagger B_0(R)B_0(S)\dd\mu_2(R)\dd\mu_2(S).\label{eq:normal-square}
\end{equation}
\bi{The double-contraction term is $\mathcal L_2(h_\sigma^2)$, not $H_\sigma$. The zero-contraction term $S_4^\sigma$ is positive. The remaining term has one internal contraction. Its exact periodization requires both an unfolding identity and a COM partial trace.}
{双收缩项是$\mathcal L_2(h_\sigma^2)$，不是$H_\sigma$。零收缩项$S_4^\sigma$非负。余下项具有一条内部收缩。对它作精确周期化，需要同时使用展开恒等式和COM偏迹。}

\subsection{\bi{The CCR and CAR kernels and their phases}{CCR、CAR核及其相位}}
\bi{For bosons set $a_R=a(u_0^T(R))$ and $g_{RS}=\langle u_0^T(R),u_0^T(S)\rangle$. The canonical commutation relations (CCR) give}
{玻色情形记$a_R=a(u_0^T(R))$及$g_{RS}=\langle u_0^T(R),u_0^T(S)\rangle$。正则对易关系（CCR）给出}
\begin{equation}
 a_R^2a_S^{\dagger2}=a_S^{\dagger2}a_R^2+4g_{RS}a_S^\dagger a_R+2g_{RS}^2.
\end{equation}
\bi{Since $B_0(R)=a_R^2/\sqrt2$, the three-body integrand is $g_{RS}a_R^{\dagger2}a_S^\dagger a_Ra_S^2$. Define the unnormalized physical vectors $t_{RRS}=a_R^{\dagger2}a_S^\dagger|\mathrm{vac}\rangle$ and $t_{RSS}=a_R^\dagger a_S^{\dagger2}|\mathrm{vac}\rangle$. Its coefficient operator is therefore}
{由于$B_0(R)=a_R^2/\sqrt2$，三体被积项为$g_{RS}a_R^{\dagger2}a_S^\dagger a_Ra_S^2$。定义未另行规范化的物理向量$t_{RRS}=a_R^{\dagger2}a_S^\dagger|\mathrm{vac}\rangle$及$t_{RSS}=a_R^\dagger a_S^{\dagger2}|\mathrm{vac}\rangle$，其系数算符为}
\begin{equation}
 \mathsf S_3^B=\int_{T^2}g_{RS}|t_{RRS}\rangle\langle t_{RSS}|\dd\mu_2(R)\dd\mu_2(S).
\end{equation}
\bi{The double contraction is consistent with $\langle v_R,v_S\rangle=g_{RS}^2$: its Fock coefficient is $\frac12g_{RS}^2a_R^{\dagger2}a_S^2$, exactly the kernel of $\mathcal L_2(h_B^2)$.}
{双收缩与$\langle v_R,v_S\rangle=g_{RS}^2$一致：其Fock系数为$\frac12g_{RS}^2a_R^{\dagger2}a_S^2$，恰是$\mathcal L_2(h_B^2)$的核。}

\bi{For fermions take $v_R=a\wedge b$, $v_S=c\wedge d$, $g_{xy}=\langle x,y\rangle$, and denote single-particle annihilation by $a(x)$. With $B_R=a(b)a(a)$ and $B_S=a(d)a(c)$, the canonical anticommutation relations (CAR) give}
{费米情形取$v_R=a\wedge b$、$v_S=c\wedge d$、$g_{xy}=\langle x,y\rangle$，以$a(x)$表示一粒子湮灭算符。令$B_R=a(b)a(a)$、$B_S=a(d)a(c)$，正则反对易关系（CAR）给出}
\begin{align}
 B_RB_S^\dagger={}&\det(g)I-g_{ac}a(d)^\dagger a(b)+g_{bc}a(d)^\dagger a(a)\notag\\
 &+g_{ad}a(c)^\dagger a(b)-g_{bd}a(c)^\dagger a(a)+B_S^\dagger B_R.
\end{align}
\bi{After multiplying by $B_R^\dagger$ on the left and $B_S$ on the right, the four single-contraction terms are}
{左乘$B_R^\dagger$、右乘$B_S$后，四个单收缩项为}
\begin{align}
 &-g_{ac}[a(d)B_R]^\dagger[a(b)B_S],\qquad
 +g_{bc}[a(d)B_R]^\dagger[a(a)B_S],\notag\\
 &+g_{ad}[a(c)B_R]^\dagger[a(b)B_S],\qquad
 -g_{bd}[a(c)B_R]^\dagger[a(a)B_S].\label{eq:four-CAR}
\end{align}
\bi{Integrating them gives $S_3^F$. This calculation fixes the signs without assuming that the moving torus jets are orthogonal.}
{将它们积分即得$S_3^F$。这一计算不假设运动中的环面jets正交，并明确固定了全部符号。}

\bi{In each of these bosonic or fermionic coefficient kernels, the exterior rank-one operator $E(R,S)$ contains one $S$-centered orbital in the ket and two in the bra. Equation~\eqref{eq:gram-unfold} thus gives}
{在上述玻色或费米系数核中，外部rank-one算符$E(R,S)$的ket含一个以$S$为中心的轨道，bra含两个。于是式~\eqref{eq:gram-unfold}给出}
\begin{equation}
 E(R,S+\lambda)=b_\lambda(S)^{-1}E(R,S),\qquad
 g^T_{ij}(R,S)E(R,S)=\sum_{\lambda\in\Lambda}g^{\rm pl}_{ij}(R,S+\lambda)E(R,S+\lambda).
\end{equation}
\bi{The phases agree term by term, not merely in modulus. The sum and the $S$-cell integral therefore unfold to $S\in\mathbb C$, while $R$ remains in a cell and every exterior orbital remains periodized. Fixed-order planar contractions are polynomial-Gaussian, and exterior jets change only by unit phases under cell translations, so the unfolding is absolutely convergent. One internal edge leaves no unaccounted winding variable.}
{这里相位逐项相同，不只是模相同。因此求和与$S$基本胞积分展开为$S\in\mathbb C$，而$R$仍在基本胞内，所有外部轨道仍保持周期化。固定阶平面收缩是多项式乘高斯函数，外部jets在胞平移下只改变单位相位，所以展开绝对收敛。一条内部边不会留下尚未处理的winding变量。}

\subsection{\bi{Anchored operators and the order of the limiting operations}{Anchored算符与取极限的次序}}
\bi{On the planar three-particle space put $A_B=\tfrac12n_0(n_0-1)$ and $A_F=n_0n_1$, and define}
{在平面三粒子空间上令$A_B=\tfrac12n_0(n_0-1)$、$A_F=n_0n_1$，定义}
\begin{equation}
 K_\sigma=A_\sigma(H_{3,\sigma}^{\rm pl}-I).\label{eq:anchored-K}
\end{equation}
\bi{The unfolded kernel is $\int_TC_R^{(3)}K_\sigma C_R^{(3)\dagger}\dd\mu_2(R)$. To see the subtraction in~\eqref{eq:anchored-K}, normal order $A_\sigma H_{3,\sigma}^{\rm pl}$ at the fixed first center. The four-body term vanishes on three particles, and the integrated double contraction is $A_\sigma$, since the planar pair interaction is the projector onto the pair space. The single contraction is what remains. For bosons one can also read the seed directly as}
{展开后的核为$\int_TC_R^{(3)}K_\sigma C_R^{(3)\dagger}\dd\mu_2(R)$。为理解~\eqref{eq:anchored-K}中的减项，在固定第一个中心处正规序$A_\sigma H_{3,\sigma}^{\rm pl}$。四体项在三粒子空间上为零，积分后的双收缩项是$A_\sigma$，因为平面两体相互作用是pair空间的投影子。余下的就是单收缩。玻色种子还可直接写成}
\begin{equation}
 K_B=\int_{\mathbb C}e^{-|\beta|^2/2}|t_{00\beta}\rangle\langle t_{0\beta\beta}|\dd\mu_2(\beta).
\end{equation}
\bi{For example, with $E(0,\beta)=|t_{00\beta}\rangle\langle t_{0\beta\beta}|$, the precise covariance is}
{例如，取$E(0,\beta)=|t_{00\beta}\rangle\langle t_{0\beta\beta}|$，精确协变关系为}
\begin{align}
 C_R^{(3)}E(0,\beta)C_R^{(3)\dagger}
 &=e^{i\operatorname{Im}(R\bar\beta)}E_T(R,R+\beta),\notag\\
 g^{\rm pl}_{00}(R,R+\beta)
 &=e^{-|\beta|^2/2+i\operatorname{Im}(R\bar\beta)}.
\end{align}
\bi{Translation covariance supplies the same phase in this seed and the contraction kernel. Neither anchored operator is generally self-adjoint, but its integral is; replacing it by its Hermitian part has no effect.}
{平移协变性使这一种子与收缩核具有相同的相位。Anchored算符一般不自伴，但其积分自伴；取其Hermitian部分不改变结果。}

\bi{The required infinite sums have explicit bounds. In the bosonic total-degree $l>0$ block, $A_B$ has support only on $|0,0,l\rangle$, and}
{所需无限和具有明确的界。在玻色总次数$l>0$的块中，$A_B$仅支撑于$|0,0,l\rangle$，并且}
\begin{equation}
 (H_{3,B}^{\rm pl}-I)|0,0,l\rangle
 =2\,2^{-l/2}B_l^\dagger|e_0\rangle_{1p},\qquad
 \|K_{B,l}\|_1\le2\sqrt3\,2^{-l/2}.\label{eq:anchored-B-bound}
\end{equation}
\bi{The degree-zero scalar is $6$. The vector on the right is created from the one-particle $e_0$ state, not the vacuum; its norm is bounded by $\sqrt3$ because pair insertion into three bosons is $\sqrt3$ times a symmetric projection. In the fermionic degree $L=l+1$ block, $l\ge2$,}
{次数零处的标量为$6$。右侧从一粒子$e_0$态而非真空出发；其范数不超过$\sqrt3$，因为向三玻色子空间插入pair是$\sqrt3$乘对称投影。费米次数$L=l+1$、$l\ge2$的块满足}
\begin{align}
 K_{F,L}&=|0,1,l\rangle\langle0,1,l|(H_{3,F}^{\rm pl}-I),\notag\\
 \|K_{F,L}\|_1&\le[\sqrt{2l}+l-1]2^{-l/2}.\label{eq:anchored-F-bound}
\end{align}
\bi{Here the $01$ pair contributes exactly the subtracted identity. The other two pair coefficients have magnitudes $\sqrt l\,2^{-(l-1)/2}$ and $(l-1)2^{-l/2}$; wedging a normalized pair with a normalized orbital has norm at most one. Each block has rank at most one, so these operator-norm bounds are also trace-norm bounds.}
{其中$01$ pair恰给出被减去的恒等项；其余两个pair系数的模分别为$\sqrt l\,2^{-(l-1)/2}$及$(l-1)2^{-l/2}$。规范化pair与规范化轨道的楔积范数不超过一。每块秩至多一，因此这些算符范数界同时也是迹范数界。}

\bi{The legitimate order is: truncate total degree; perform the unitary relative-COM decomposition in each finite block; periodize its seeds and descendants; integrate the parameter ladders; then remove the degree cutoff. Equal descendant levels retain all relative matrix units, including off-diagonal ones. Lemma~\ref{lem:jets} bounds the squared threefold periodization norm by $C(L+1)^6$, which is summable against~\eqref{eq:anchored-B-bound} or~\eqref{eq:anchored-F-bound}. Thus}
{合法次序是：先截断总次数；在每个有限块中作相对—COM酉分解；周期化种子及后裔；积分参数梯级；最后去掉次数截断。相等后裔级数保留全部相对matrix units，包括非对角项。引理~\ref{lem:jets}把三次周期化的范数平方控制为$C(L+1)^6$，它与~\eqref{eq:anchored-B-bound}或~\eqref{eq:anchored-F-bound}相乘仍可和。因此}
\begin{equation}
 \mathsf S_3^\sigma=\Omega_3\!\left(\frac23\Tr_{\CM}K_\sigma\right).\label{eq:connected-partial-trace}
\end{equation}
\bi{An integrated planar spectral formula cannot replace $K_\sigma$ pointwise. For example, the fermionic anchored operator has a nonzero matrix element on $|0,1,3\rangle$, whereas the relative highest seed of the corresponding total degree is zero. COM descendants are removed only by the integral, not beforehand.}
{不能用积分后的平面谱公式逐中心替换$K_\sigma$。例如，费米anchored算符在$|0,1,3\rangle$上有非零矩阵元，而对应总次数的相对最高种子为零。COM后裔只能经积分消去，不能事先删除。}

\subsection{\bi{Relative eigenvalues from compressed exchanges}{由压缩交换得到相对本征值}}
\bi{Compressed pair exchanges are the recoupling mechanism behind the fermionic three-body spectrum in OpenAI~\cite[Section~3]{OpenAI2026Fock} and its bosonic counterpart in Shen~\cite[Section~1]{Shen2026BosonicCandidate}. Here we derive the planar oscillator coefficients in the normalization needed by the torus transfer. The anchored unfolding and the convergence argument above are additional geometric obligations; a planar spectral formula alone does not establish them.}
{压缩pair交换是OpenAI费米三体谱~\cite[第3节]{OpenAI2026Fock}及Shen玻色对应结果~\cite[第1节]{Shen2026BosonicCandidate}中的重耦合机制。这里在环面迁移所需的归一化下推导平面振子系数。上文的anchored展开及收敛论证是额外的几何环节，不能仅由平面谱公式代替。}
\bi{Let $U$ be the planar one-particle LLL and $V_\sigma$ its relative-degree-zero bosonic or relative-degree-one fermionic pair space. Define pair insertion $W_3^B(v\otimes u)=B(v)^\dagger u$ and $W_3^F(v\otimes u)=v\wedge u$. Let $\iota:V_\sigma\otimes U\to U^{\otimes3}$ be the normalized tensor embedding, with symmetry or antisymmetry in the first pair, and $S_{ij}=\iota^\dagger\tau_{ij}\iota$ the compression of the ambient permutation. The bare permutation need not preserve the domain. Then}
{令$U$为平面一粒子LLL，$V_\sigma$为其中相对次数零的玻色pair空间或相对次数一的费米pair空间。定义pair插入$W_3^B(v\otimes u)=B(v)^\dagger u$及$W_3^F(v\otimes u)=v\wedge u$。令$\iota:V_\sigma\otimes U\to U^{\otimes3}$为前两个因子对称或反对称的规范化张量嵌入，并令$S_{ij}=\iota^\dagger\tau_{ij}\iota$为环境置换的压缩。裸置换无需保持定义域。于是}
\begin{equation}
 W_3^{B\dagger}W_3^B=I+S_{13}+S_{23},\qquad
 W_3^{F\dagger}W_3^F=I-S_{13}-S_{23},\qquad
 W_3^\sigma W_3^{\sigma\dagger}=H_{3,\sigma}^{\rm pl}.\label{eq:compressed-exchanges}
\end{equation}
\bi{The pair COM and spectator oscillator have one relative oscillator after separating their total COM. Normalized relative highest seeds are}
{把pair的COM振子与spectator振子的总COM分离后，只有一个相对振子。规范化相对最高种子为}
\begin{align}
 \xi_z^\sigma&=\sum_{p=0}^zs_{zp}^\sigma v_p^\sigma\otimes e_{z-p},\qquad
 w_z^\sigma=W_3^\sigma\xi_z^\sigma,\notag\\
 s_{zp}^B&=(-1)^p\sqrt{\binom zp}\,3^{-p/2}(2/3)^{(z-p)/2},\qquad
 s_{zp}^F=(-1)^{z-p}\sqrt{\binom zp}\,3^{-p/2}(2/3)^{(z-p)/2}.
\end{align}
\bi{Each compressed exchange is scalar on this one-dimensional relative level, since it intertwines COM translations and preserves total degree. For bosons, its $v_0\otimes e_z$ component receives only the $p=z$ contribution. The ratio of seed coefficients is $(-1)^z2^{-z/2}$ and the ordered pair tensor coefficient is $2^{-z/2}$, giving $b_z^B=(-1/2)^z$. For fermions the two contributions are $p=z-1,z$. With antisymmetric continuation of the planar pair coefficients,}
{每个压缩交换在这一维相对级数上为标量，因为它交织COM平移并保持总次数。玻色情形的$v_0\otimes e_z$分量仅收到$p=z$项；种子系数比为$(-1)^z2^{-z/2}$，有序pair张量系数为$2^{-z/2}$，故$b_z^B=(-1/2)^z$。费米情形有$p=z-1,z$两项。对平面pair系数作反对称延拓后，}
\begin{equation}
 b_z^F=\frac{-s_{z,z-1}^Fv_{z-1}^F(0,z)+s_{z,z}^Fv_z^F(1,z)}{2s_{z,0}^F}
 =(-1)^{z-1}\frac{3z-1}{2^{z+1}}.
\end{equation}
\bi{Out-of-range terms vanish; the expression also gives $z=0,1$ directly. Therefore $W_3^{\sigma\dagger}W_3^\sigma$ has eigenvalue $1+q_z^\sigma$, where}
{区间外项取零；该表达也直接给出$z=0,1$。因此$W_3^{\sigma\dagger}W_3^\sigma$的本征值为$1+q_z^\sigma$，其中}
\begin{equation}
 q_z^B=2(-1/2)^z,\qquad q_z^F=(3z-1)(-1/2)^z.\label{eq:q-values}
\end{equation}
\bi{Using~\eqref{eq:compressed-exchanges},}
{由~\eqref{eq:compressed-exchanges}可得}
\begin{equation}
 (H_{3,\sigma}^{\rm pl})^2-H_{3,\sigma}^{\rm pl}
 =I_{\CM}\otimes\sum_{z\ge0}q_z^\sigma|w_z^\sigma\rangle\langle w_z^\sigma|.
\end{equation}
\bi{On the other hand, $\int A_\sigma(R)\dd\mu_2(R)=H_{3,\sigma}^{\rm pl}$ and the latter commutes with simultaneous translations. Equation~\eqref{eq:plane-twirl} applied to~\eqref{eq:anchored-K} identifies $(2/3)\Tr_{\CM}K_\sigma$ with this relative coefficient. Combining with~\eqref{eq:connected-partial-trace} proves the following proposition.}
{另一方面，$\int A_\sigma(R)\dd\mu_2(R)=H_{3,\sigma}^{\rm pl}$，且后者与同时平移对易。把式~\eqref{eq:plane-twirl}用于~\eqref{eq:anchored-K}，便识别出$(2/3)\Tr_{\CM}K_\sigma$恰为该相对系数。结合~\eqref{eq:connected-partial-trace}，即证明下述命题。}

\begin{proposition}[\bi{Exact connected coefficients}{精确连通系数}]\label{prop:connected}
\bi{Let $D_z^\sigma(R)=A(w_z^{\sigma,T}(R))$. In either statistics, on each finite particle sector,}
{令$D_z^\sigma(R)=A(w_z^{\sigma,T}(R))$。对两种统计，在每个有限粒子扇区上，}
\begin{align}
 S_3^B&=\int_T\sum_{z\ge0}2(-1/2)^zD_z^{B\dagger}D_z^B\dd\mu_3,\label{eq:boson-threebody}\\
 S_3^F&=\int_T\sum_{z\ge0}(3z-1)(-1/2)^zD_z^{F\dagger}D_z^F\dd\mu_3.\label{eq:fermion-threebody}
\end{align}
\bi{The series are defined by the absolutely convergent anchored construction above. For bosons $w_1^B=0$; for fermions $w_0^F=w_1^F=w_3^F=0$. These are algebraic zero vectors because $\|w_z^\sigma\|^2=1+q_z^\sigma=0$, so periodization preserves their vanishing. All other negative odd-degree terms must be controlled, rather than discarded.}
{级数通过上文绝对收敛的anchored构造定义。玻色有$w_1^B=0$；费米有$w_0^F=w_1^F=w_3^F=0$。它们是代数零向量，因为$\|w_z^\sigma\|^2=1+q_z^\sigma=0$，所以周期化保持其为零。其他负的奇次数项都必须控制，不能直接丢弃。}
\end{proposition}

\section{\bi{The fermionic comparison and its uniform error budget}{费米比较不等式及其一致误差预算}}
\label{sec:F-proof}\label{sec:F-fermionic}

\bi{We now complete the proof of Theorem~\ref{thm:fermion}. Throughout this section, $H=H_{\mathrm F}$ is the full $V_1$ Hamiltonian with the normalization fixed in the model definition. We abbreviate $B_p=B_p^F$, $D_z=D_z^F$, $C_3=S_3^F$, $S_4=S_4^F$, $h_T=h_F$, and $\varepsilon_2^{\mathrm F}=\epsilon_2^{(1)}$. We use the raw periodized jets $c_j(R)=c(u_j^T(R))$. The geometric estimates and the exact center-of-mass (COM) integration identity were established in Lemmas~\ref{lem:jets}, \ref{lem:relative-frame}, and~\ref{lem:ladder}. The remaining task is to compare the three- and four-body terms with a fixed sum of squares. All errors will be bounded by a scalar multiple of $H$.}{本节完成定理~\ref{thm:fermion}的证明。以下记 $H=H_{\mathrm F}$，表示按模型定义归一化的完整 $V_1$ 哈密顿量；简记 $B_p=B_p^F$、$D_z=D_z^F$、$C_3=S_3^F$、$S_4=S_4^F$、$h_T=h_F$ 和 $\varepsilon_2^{\mathrm F}=\epsilon_2^{(1)}$，并使用未经正交化的周期化 jets $c_j(R)=c(u_j^T(R))$。几何估计和精确的质心积分恒等式已由引理~\ref{lem:jets}、\ref{lem:relative-frame}及~\ref{lem:ladder}给出。余下任务是用一组固定的正平方比较三体和四体项，并将所有误差控制为 $H$ 的标量倍数。}

\bi{The finite coefficient data are those of the Fock-space certificate in \cite{OpenAI2026Fock}. Their role here is entirely algebraic: Appendix~\ref{app:F-certificate} prints all seven integer rows and gives the rational calculations proving the finite inequalities used below. No spherical spectral-gap statement is an input to this argument.}{有限系数取自文献~\cite{OpenAI2026Fock}的 Fock 空间证书，其在本证明中的作用完全是代数性的。附录~\ref{app:F-certificate}列出全部七行整数，并给出下文有限不等式的有理算术验证。本论证不以任何球面谱隙定理为前提。}

\subsection{\bi{A four-body frame error paid for by the pair energy}{以对能量支付四体框架误差}}
\label{subsec:F-four-frame}

\bi{The zero-contraction term in the normal-ordered square of $H$ is}{将 $H$ 的平方正规序后，其零收缩项为}
\begin{equation}
 S_4=\int_{T\times T}(B_0(R)B_0(S))^\dagger
                  (B_0(R)B_0(S))\,\dd\mu_2(R)\dd\mu_2(S)\ge0.
 \label{eq:F-S4}
\end{equation}
\bi{For a Fock-space vector $\psi$, the function $B_0(R)B_0(S)\psi$ lies in $\mathcal H_2(T)\otimes\mathcal H_2(T)$ with values in the remaining Fock space. Apply the finite relative-jet map with charges $a=b=2$ and relative degrees $0\le d\le22$. Its components are}{对任意 Fock 空间向量 $\psi$，函数 $B_0(R)B_0(S)\psi$ 属于以剩余 Fock 空间为值域的 $\mathcal H_2(T)\otimes\mathcal H_2(T)$。取电荷 $a=b=2$，并对相对阶数 $0\le d\le22$ 应用有限相对 jet 映射，其各分量为}
\begin{equation}
 C_d(R)=\sum_{p=0}^{d}(-1)^{d-p}
          \sqrt{\binom dp}\,2^{-d/2}B_p(R)B_{d-p}(R),
 \qquad
 T_4=\int_T\sum_{d=0}^{22}C_d(R)^\dagger C_d(R)\,\dd\mu_4(R).
 \label{eq:F-Cd}
\end{equation}
\bi{Pair annihilators commute because each has even fermionic degree. Consequently $C_d=0$ for odd $d$, by exchanging $p$ and $d-p$. Lemma~\ref{lem:relative-frame}, applied to its tensor product with the identity on the spectator space, gives}{对湮灭算符均具有偶数费米次数，故两两对易。交换 $p$ 与 $d-p$ 可知奇数 $d$ 时 $C_d=0$。将引理~\ref{lem:relative-frame}中的映射与旁观者空间上的恒等算符作张量积，得到}
\begin{equation}
 T_4\le(1+\varepsilon_4^{\mathrm F}(\Lambda))S_4,
 \qquad
 \varepsilon_4^{\mathrm F}(\Lambda)=\epsilon_{2,2,22}(\Lambda)
       \xrightarrow[\ell(\Lambda)\to\infty]{}0.
 \label{eq:F-four-bessel}
\end{equation}
\bi{This use of the frame bound keeps every COM theta mode and does not introduce the dimension of the spectator space. A second estimate converts its error into pair energy.}{这一框架估计保留全部质心 theta 模式，且不引入旁观者空间的维数。下面的第二个估计将其误差转化为对能量。}

\bi{The method has a spherical antecedent in the fixed-block relative-energy estimate of OpenAI~\cite[Section~7, vanishing-weight lemma and fixed-block-energy estimate]{OpenAI2026Stability}: first bound the finite left pair annihilator, then average the remaining pair energy. The torus proof below supplies its raw-Gram correction and uses the exact parameter pair integral.}
{此方法的球面前源是OpenAI的固定块相对能量估计~\cite[第7节，vanishing-weight引理与fixed-block-energy估计]{OpenAI2026Stability}：先控制有限左pair湮灭算符，再平均剩余pair能量。下文的环面证明给出原始Gram矩阵修正，并使用精确参数pair积分。}
\begin{lemma}[\bi{Finite left-factor estimate}{有限左因子估计}]
\label{lem:F-left-factor}
\bi{For $0\le p\le22$, and for every center $R$,}{对任意 $0\le p\le22$ 及中心 $R$，有}
\begin{equation}
 \|B_p(R)\|^2\le(p+2)(1+e_{23}(\Lambda))^2.
 \label{eq:F-left-factor}
\end{equation}
\bi{In particular,}{特别地，}
\begin{equation}
 T_4\le C_4^{\mathrm F}(\Lambda)H,
 \qquad C_4^{\mathrm F}(\Lambda)=5152(1+e_{23}(\Lambda))^2.
 \label{eq:F-four-relative}
\end{equation}
\end{lemma}
\begin{proof}
\bi{Let $C_R:e_j\mapsto u_j^T(R)$ be restricted to the first $p+2$ orbitals. Its Gram matrix is bounded above by $(1+e_{23})I$. The planar pair vector $v_p$ has norm one, so the periodized vector $(\bigwedge^2 C_R)v_p$ has norm at most $1+e_{23}$. A two-form on a space of dimension at most $p+2$ has a canonical orthogonal decomposition $\sum_{a=1}^{r}\sigma_a f_{2a-1}\wedge f_{2a}$, where $r\le(p+2)/2$ and $\sum_a|\sigma_a|^2$ is its squared norm. The corresponding annihilator has norm at most $\sum_a|\sigma_a|$, since each product of two normalized one-particle annihilators is a contraction. Cauchy--Schwarz therefore bounds its squared norm by $r\sum_a|\sigma_a|^2$, which implies the slightly weaker bound \eqref{eq:F-left-factor}.}{将映射 $C_R:e_j\mapsto u_j^T(R)$ 限制于前 $p+2$ 个轨道，其 Gram 矩阵不超过 $(1+e_{23})I$。平面对向量 $v_p$ 的范数为一，故周期化向量 $(\bigwedge^2 C_R)v_p$ 的范数不超过 $1+e_{23}$。维数不超过 $p+2$ 的空间上的二阶外积向量可作标准正交分解 $\sum_{a=1}^{r}\sigma_a f_{2a-1}\wedge f_{2a}$，其中 $r\le(p+2)/2$，且 $\sum_a|\sigma_a|^2$ 等于该向量的范数平方。每个归一化单粒子湮灭算符都是压缩算符，因此相应对湮灭算符的范数不超过 $\sum_a|\sigma_a|$。由 Cauchy--Schwarz 不等式，其范数平方不超过 $r\sum_a|\sigma_a|^2$，从而得到略弱的估计~\eqref{eq:F-left-factor}。}

\bi{The coefficients in \eqref{eq:F-Cd} have squared sum one. Thus, for every $\psi$,}{式~\eqref{eq:F-Cd}中的系数平方和为一。因此对任意 $\psi$，}
\begin{align}
 \|C_d(R)\psi\|^2
 &\le\sum_{p=0}^{d}\|B_p(R)B_{d-p}(R)\psi\|^2\notag\\
 &\le(1+e_{23})^2\sum_{p=0}^{d}(p+2)
                      \|B_{d-p}(R)\psi\|^2.
 \label{eq:F-four-pointwise}
\end{align}
\bi{Integrating, use $\dd\mu_4=2\dd\mu_2$ and the exact identity $\int_T B_q^\dagger B_q\,\dd\mu_2=H$ for every $q$, from Lemma~\ref{lem:ladder}. This gives}{积分后，使用 $\dd\mu_4=2\dd\mu_2$ 以及引理~\ref{lem:ladder}对每个 $q$ 给出的精确恒等式 $\int_T B_q^\dagger B_q\,\dd\mu_2=H$，得到}
\begin{equation}
 T_4\le2(1+e_{23})^2
       \sum_{d=0}^{22}\sum_{p=0}^{d}(p+2)H
     =5152(1+e_{23})^2H.
\end{equation}
\bi{Every estimate used a bounded left factor followed by an exact pair-energy integral, so the result is independent of particle number.}{各步均先控制有界的左因子，再使用精确的对能量积分，因此结果与粒子数无关。}
\end{proof}

\bi{Combining \eqref{eq:F-four-bessel} and \eqref{eq:F-four-relative} yields the form that will enter the final budget:}{合并式~\eqref{eq:F-four-bessel}和~\eqref{eq:F-four-relative}，得到最终预算所需的形式：}
\begin{equation}
 S_4\ge\frac{T_4}{1+\varepsilon_4^{\mathrm F}}
      \ge T_4-C_4^{\mathrm F}(\Lambda)\varepsilon_4^{\mathrm F}(\Lambda)H.
 \label{eq:F-four-payment}
\end{equation}
\bi{The last inequality follows from $1-(1+x)^{-1}=x/(1+x)\le x$ for $x\ge0$. In particular, the error does not multiply $S_4$, a particle-counting operator, or an identity with a size-dependent coefficient.}{最后一步使用 $x\ge0$ 时 $1-(1+x)^{-1}=x/(1+x)\le x$。特别地，该误差不乘以 $S_4$、粒子数计数算符，或带有尺寸依赖系数的恒等算符。}

\subsection{\bi{Pauli control of the complete negative three-body tail}{完整三体负尾的 Pauli 控制}}
\label{subsec:F-tail}

\bi{The exact three-body identity \eqref{eq:fermion-threebody} contains all relative degrees. Its negative tail cannot be controlled by extending a fixed-degree Gram estimate to arbitrarily high jets. The high-jet estimate of Lemma~\ref{lem:jets}, $\sup_R\|u_j^T(R)\|^2\le C(1+j)$ for $\ell(\Lambda)\ge1$, supplies the uniform majorant needed for the following Pauli estimate.}{精确三体恒等式~\eqref{eq:fermion-threebody}包含全部相对阶数，故不能将固定阶 Gram 估计外推至任意高阶 jets 来控制其负尾。引理~\ref{lem:jets}的高阶 jet 估计在 $\ell(\Lambda)\ge1$ 时给出 $\sup_R\|u_j^T(R)\|^2\le C(1+j)$，为下述 Pauli 估计提供所需的一致控制项。}

\bi{Recall the coefficients and annihilators in the exact three-body identity:}{回顾精确三体恒等式中的系数与湮灭算符：}
\begin{equation}
 \begin{aligned}
 q_z&=(3z-1)(-1/2)^z,\qquad
 D_z(R)=\sum_{p=0}^{z}s_{z,p}\,c_{z-p}(R)B_p(R),\\
 s_{z,p}&=(-1)^{z-p}\sqrt{\binom zp}\,3^{-p/2}(2/3)^{(z-p)/2}.
 \end{aligned}
 \label{eq:F-Dz}
\end{equation}
\bi{The coefficients have squared sum one, and $D_0,D_1,D_3$ vanish as algebraic wedge identities. All even terms are nonnegative. Define}{这些系数的平方和为一，且外积代数恒等式给出 $D_0=D_1=D_3=0$。所有偶数阶项均非负。定义}
\begin{equation}
 M_z(\Lambda)=\sup_{R\in T,\,0\le j\le z}\|u_j^T(R)\|^2,
 \qquad
 \delta_\Lambda=\frac32\sum_{\substack{z\ge17\\z\ \mathrm{odd}}}
                     |q_z|(z+1)M_z(\Lambda).
 \label{eq:F-delta}
\end{equation}
\bi{The canonical anticommutation relations give $\|c_j(R)\|=\|u_j^T(R)\|$. Cauchy--Schwarz in the coefficient index therefore gives}{正则反对易关系给出 $\|c_j(R)\|=\|u_j^T(R)\|$。对系数指标使用 Cauchy--Schwarz 不等式，便有}
\begin{equation}
 \|D_z(R)\psi\|^2
 \le\sum_{p=0}^{z}\|c_{z-p}(R)B_p(R)\psi\|^2
 \le M_z(\Lambda)\sum_{p=0}^{z}\|B_p(R)\psi\|^2.
 \label{eq:F-Pauli-tail}
\end{equation}
\bi{After integrating with $\dd\mu_3=(3/2)\dd\mu_2$, the exact pair-frame identity pays for each of the $z+1$ summands by $H$. Thus the negative tail satisfies}{按 $\dd\mu_3=(3/2)\dd\mu_2$ 积分后，精确的对框架恒等式将 $z+1$ 个求和项中的每一项均控制为 $H$，故负尾满足}
\begin{equation}
 \int_T\sum_{\substack{z\ge17\\z\ \mathrm{odd}}}
       q_zD_z^\dagger D_z\,\dd\mu_3\ge-\delta_\Lambda H.
 \label{eq:F-tail-bound}
\end{equation}
\bi{For each fixed $z$, $M_z(\Lambda)\to1$. Lemma~\ref{lem:jets} gives $M_z\le C(1+z)$ uniformly for $\ell\ge1$, while $|q_z|$ decays geometrically. Dominated convergence therefore yields}{对每个固定 $z$，有 $M_z(\Lambda)\to1$。引理~\ref{lem:jets}在 $\ell\ge1$ 时给出统一上界 $M_z\le C(1+z)$，而 $|q_z|$ 呈几何衰减。于是由控制收敛定理，}
\begin{equation}
 \lim_{\ell(\Lambda)\to\infty}\delta_\Lambda
 =\frac32\sum_{\substack{z\ge17\\z\ \mathrm{odd}}}
           (3z-1)(z+1)2^{-z}
 =\sum_{m=8}^{\infty}(9m^2+12m+3)4^{-m}
 =\frac{61}{4096}.
 \label{eq:F-tail-limit}
\end{equation}
\bi{The middle expression uses $z=2m+1$; the last equality follows by twice differentiating the geometric series beginning at $m=8$. Both the convergence and the resulting constant are uniform in particle number, boundary twist, and lattice shape subject to the stated systole condition.}{中间一式使用 $z=2m+1$；最后一式可由从 $m=8$ 开始的几何级数作两次求导得到。在所述最短周期条件下，该收敛及所得常数均与粒子数、边界 twist 和格子形状无关。}

\subsection{\bi{Decoding the seven rows and the finite coupled spaces}{七行数据及有限耦合空间的解码}}
\label{subsec:F-data}

\bi{Set $P=10^7$. For row $\rho$, the printed integers are $t_\rho$, $\ell_\rho$, and $a_{pj}^{\rho}$; missing entries are zero. Define}{令 $P=10^7$。对第 $\rho$ 行，所列整数为 $t_\rho$、$\ell_\rho$ 和 $a_{pj}^{\rho}$，未列项均为零。定义}
\begin{equation}
 \lambda_\rho=\frac{\ell_\rho}{P},\qquad
 \alpha_{pj}^{\rho}=\frac{a_{pj}^{\rho}}{P}
       \sqrt{\frac{2^{p-t_\rho}t_\rho!(p+j-t_\rho)!}{p!j!}},
 \qquad
 \eta=\sum_{\rho=1}^{7}\lambda_\rho^2
      =\frac{93527408868499}{10^{14}}.
 \label{eq:F-row-decode}
\end{equation}
\bi{The seven $t$ values are $0,1,1,2,5,5,7$. Every nonzero entry satisfies $0\le p\le7$, $0\le j\le8$, and $0\le p+j-t_\rho\le8$. These support conditions bound all auxiliary one-particle modes by eight, although their coupled four-body coefficients require jets up to degree twenty-four.}{七个 $t$ 值依次为 $0,1,1,2,5,5,7$。每个非零项均满足 $0\le p\le7$、$0\le j\le8$ 及 $0\le p+j-t_\rho\le8$。这些支撑条件使辅助单粒子模式的编号均不超过八，但其耦合四体系数需要最高二十四阶的 jets。}

\bi{To specify the coefficients without referring to a spherical limit, take $0<u<1$, put $v=\sqrt{1-u^2}$, and define $s^{(u)}_{z,T,p}$ by the exact polynomial identity}{为不借助球面极限直接指定系数，取 $0<u<1$，令 $v=\sqrt{1-u^2}$，并由以下精确多项式恒等式定义 $s^{(u)}_{z,T,p}$：}
\begin{equation}
 \frac{(uX-vY)^z(vX+uY)^{T-z}}{\sqrt{z!(T-z)!}}
 =\sum_{p=0}^{T}s^{(u)}_{z,T,p}
       \frac{X^pY^{T-p}}{\sqrt{p!(T-p)!}}.
 \label{eq:F-coupling}
\end{equation}
\bi{A coefficient is zero if any of its degree indices is outside its indicated range. The change of variables is orthogonal and preserves the normalized-monomial inner product. For a fixed $T$, the rows indexed by $z$ consequently form an orthogonal matrix. The factor $(vX+uY)^{T-z}/\sqrt{(T-z)!}$ is the normalized COM descendant. This identifies the algebraic transformation with the ladder normalization of Lemma~\ref{lem:ladder}.}{当次数指标超出所示范围时，系数定义为零。该变量变换正交，并保持归一化单项式内积，因此固定 $T$ 时，以 $z$ 为行标的系数构成正交矩阵。因子 $(vX+uY)^{T-z}/\sqrt{(T-z)!}$ 正是归一化的质心后裔，所以此代数变换与引理~\ref{lem:ladder}的梯级归一化一致。}

\bi{For the three-body comparison, set $u=1/\sqrt3$ and abbreviate $s=s^{(1/\sqrt3)}$. Define}{三体比较中取 $u=1/\sqrt3$，简记 $s=s^{(1/\sqrt3)}$，并定义}
\begin{equation}
 e_z=\sum_{\rho=1}^{7}\sum_{T=\max(z,t_\rho)}^{15}
 \left[
 2\lambda_\rho s_{z,T,t_\rho}
       \sum_{p+j=T}\alpha_{pj}^{\rho}s_{z,T,p}
 +\left(\sum_{p+j=T}\alpha_{pj}^{\rho}s_{z,T,p}\right)^2
 \right].
 \label{eq:F-ez}
\end{equation}
\bi{The exact rational calculation in Appendix~\ref{app:F-three-arithmetic} proves the thirteen inequalities}{附录~\ref{app:F-three-arithmetic}中的精确有理计算证明以下十三个不等式：}
\begin{equation}
 e_z<\frac32q_z,
 \qquad z\in\{2,4,5,\ldots,15\}.
 \label{eq:F-three-certificate}
\end{equation}
\bi{The excluded indices $0,1,3$ correspond to identically zero wedge images and require no scalar comparison.}{排除的指标 $0,1,3$ 对应恒为零的外积像，无需作标量比较。}

\bi{For four bodies, let $\mathcal R_D=\{1,3,5,\ldots\}\cap[1,D]$, where $1\le D\le23$. In the formal orthonormal input basis $v_p\otimes(e_j\wedge e_k)$, $j<k$, define}{四体部分令 $\mathcal R_D=\{1,3,5,\ldots\}\cap[1,D]$，其中 $1\le D\le23$。在形式上的正交归一输入基 $v_p\otimes(e_j\wedge e_k)$、$j<k$ 中定义}
\begin{equation}
 S_r(D,T;p,j,k)=\sqrt2\,
       s^{(1/\sqrt2)}_{r,j+k,j}
       s^{(1/\sqrt2)}_{D-r,T-r,p}.
 \label{eq:F-S}
\end{equation}
\bi{This is zero if $r>j+k$ and is extended antisymmetrically in $j,k$. The first coupling separates the relative degree $r$ of the two single-particle factors; odd $r$ is precisely their antisymmetry condition. The second coupling separates their pair COM from that of $v_p$. Its relative degree is $D-r$, and $T-D$ is the total COM descendant degree. The highest-weight wedge columns are}{当 $r>j+k$ 时该系数为零，并关于 $j,k$ 作反对称延拓。第一次耦合分离两个单粒子因子的相对阶数 $r$，而 $r$ 为奇数正对应其反对称性；第二次耦合将它们的对质心与 $v_p$ 的质心分开，其相对阶数为 $D-r$，总质心后裔阶数为 $T-D$。最高权外积列为}
\begin{equation}
 w_{D,r}=\sum_{\substack{p+j+k=D\\j<k}}
     S_r(D,D;p,j,k)\,v_p\wedge e_j\wedge e_k,
 \qquad W_D=(w_{D,r})_{r\in\mathcal R_D},\qquad G_D=W_D^\dagger W_D.
 \label{eq:F-wG}
\end{equation}
\bi{Each column has physical orbital degree $D+1$, because $v_p$ has degree $p+1$. With $i=p+j-t_\rho$, $k=q+l-t_\rho$, and $T=p+j+k=q+l+i$, set}{由于 $v_p$ 的轨道次数为 $p+1$，每列的物理轨道总次数为 $D+1$。令 $i=p+j-t_\rho$、$k=q+l-t_\rho$ 及 $T=p+j+k=q+l+i$，并设}
\begin{equation}
 (E_D)_{rs}=-\sum_{\rho=1}^{7}
  \sum_{\substack{(p,j),(q,l)\in\operatorname{supp}(a^\rho)\\T\ge D}}
  \alpha_{pj}^{\rho}\alpha_{ql}^{\rho}
  S_r(D,T;p,j,k)S_s(D,T;q,l,i).
 \label{eq:F-E}
\end{equation}
\bi{The entry pairs are ordered; exchanging them proves that $E_D$ is real symmetric. Their support implies $T\le23$. Let $\Pi_1$ denote the projection onto the copy index $r=1$, and fix $\varepsilon_{\mathrm F}=3\times10^{-6}$. The four-body certificate is}{两个输入项为有序取值；交换它们即可证明 $E_D$ 为实对称矩阵。支撑条件给出 $T\le23$。以 $\Pi_1$ 表示副本指标 $r=1$ 上的投影，并固定 $\varepsilon_{\mathrm F}=3\times10^{-6}$。四体证书为}
\begin{equation}
 G_D(2\Pi_1-E_D+\varepsilon_{\mathrm F}I)G_D\ge0,
 \qquad 1\le D\le23.
 \label{eq:F-four-certificate}
\end{equation}
\bi{Here $\Pi_1$ has no relation to the boundary semicharacter $\chi$. Appendix~\ref{app:F-four-arithmetic} gives a radical-free rational reconstruction of every matrix in \eqref{eq:F-four-certificate}, together with an exact positivity test.}{此处 $\Pi_1$ 与边界半特征 $\chi$ 无关。附录~\ref{app:F-four-arithmetic}给出式~\eqref{eq:F-four-certificate}中每个矩阵的无根式有理重建及精确正性检验。}

\subsection{\bi{Normal ordering of the positive comparison operator}{正比较算符的正规序展开}}
\label{subsec:F-squares}

\bi{Using the fixed coefficients above, define the torus operators}{利用上述固定系数，定义环面算符}
\begin{equation}
 F_\rho(R)=\lambda_\rho B_{t_\rho}(R)
       +\sum_{p,j}\alpha_{pj}^{\rho}
           c_{p+j-t_\rho}(R)^\dagger c_j(R)B_p(R),
 \qquad
 K_T=\int_T\sum_{\rho=1}^{7}F_\rho(R)^\dagger F_\rho(R)\,\dd\mu_2(R)\ge0.
 \label{eq:F-square}
\end{equation}
\bi{Every term ends in a pair annihilator. In particular, if $H\psi=0$, the exact pair-frame identity implies $B_p(R)\psi=0$ for every $p$ and almost every $R$, and continuity then gives this equality for all $R$. Hence $F_\rho(R)\psi=0$. The comparison thus preserves the zero-mode ideal exactly.}{每一项的右端均为对湮灭算符。特别地，若 $H\psi=0$，精确对框架恒等式推出对每个 $p$ 及几乎处处的 $R$ 均有 $B_p(R)\psi=0$；再由连续性可知该式对所有 $R$ 成立。因此 $F_\rho(R)\psi=0$，比较算符精确保留零模理想。}

\bi{For one row, suppress $\rho$ and $R$, and put $i=p+j-t$, $k=q+l-t$. Write $g_{ik}=\langle u_i^T,u_k^T\rangle$ and $O_{p,j,k}=c_kc_jB_p$. Applying $c_ic_k^\dagger=g_{ik}-c_k^\dagger c_i$ yields the exact decomposition}{先对一行省略 $\rho,R$，并令 $i=p+j-t$、$k=q+l-t$。记 $g_{ik}=\langle u_i^T,u_k^T\rangle$ 及 $O_{p,j,k}=c_kc_jB_p$。应用 $c_ic_k^\dagger=g_{ik}-c_k^\dagger c_i$，得到精确分解}
\begin{align}
 (F^\dagger F)_2={}&\lambda^2 B_t^\dagger B_t,
 \label{eq:F-normal2}\\
 (F^\dagger F)_3={}&
 \sum_{p,j}\lambda\alpha_{pj}
 \bigl[(c_iB_t)^\dagger c_jB_p+(c_jB_p)^\dagger c_iB_t\bigr]
 +\sum_{p,j;q,l}\alpha_{pj}\alpha_{ql}g_{ik}
                      (c_jB_p)^\dagger c_lB_q,
 \label{eq:F-normal3}\\
 (F^\dagger F)_4={}&-
 \sum_{p,j;q,l}\alpha_{pj}\alpha_{ql}
                         O_{p,j,k}^\dagger O_{q,l,i}.
 \label{eq:F-normal4}
\end{align}
\bi{Only the single contraction in \eqref{eq:F-normal3} contains the local Gram matrix. Let $K_3^{(0)}$ be its integral with $g_{ik}$ replaced by $\delta_{ik}$, including the displayed cross terms, and let $R_{3,\mathrm{metric}}$ be the difference. Let $K_4$ be the integral of \eqref{eq:F-normal4}, summed over rows. The pair-frame identity then gives}{只有式~\eqref{eq:F-normal3}中的单次收缩含有局部 Gram 矩阵。将其中 $g_{ik}$ 换为 $\delta_{ik}$ 并连同所示交叉项积分，记为 $K_3^{(0)}$；差值记为 $R_{3,\mathrm{metric}}$。将式~\eqref{eq:F-normal4}按全部行求和并积分，记为 $K_4$。由对框架恒等式，}
\begin{equation}
 K_T=\eta H+K_3^{(0)}+R_{3,\mathrm{metric}}+K_4.
 \label{eq:F-K-decomposition}
\end{equation}

\bi{We identify $K_3^{(0)}$ explicitly. Before wedging, its coefficient matrix is on the formal input space $v_p\otimes e_j$. The Kronecker condition $i=k$ says exactly that $p+j=q+l=T$. In a cross term the second input is $v_t\otimes e_{T-t}$, so total degree is again preserved. At fixed $T$, change to the orthonormal coupled basis \eqref{eq:F-coupling}. A relative index $z$ has COM descendant degree $T-z$. The integral in Lemma~\ref{lem:ladder} kills matrix units with unequal descendant degrees and replaces each equal-descendant diagonal block by its highest-weight matrix unit. This traces the COM factor but does not delete any relative matrix entries. The diagonal coefficient that remains is exactly \eqref{eq:F-ez}. Hence}{下面明确识别 $K_3^{(0)}$。在取外积之前，其系数矩阵作用于形式输入空间 $v_p\otimes e_j$。Kronecker 条件 $i=k$ 恰等价于 $p+j=q+l=T$；交叉项的另一输入为 $v_t\otimes e_{T-t}$，故同样保持总次数。固定 $T$ 后，变换到式~\eqref{eq:F-coupling}的正交耦合基，相对指标 $z$ 的质心后裔阶数为 $T-z$。由引理~\ref{lem:ladder}，积分消去后裔阶数不同的矩阵单位，并将后裔阶数相同的对角分块化为相应最高权矩阵单位。这是对质心因子取迹，并未额外删除相对空间矩阵元。留下的对角系数恰为式~\eqref{eq:F-ez}，因此}
\begin{equation}
 K_3^{(0)}=\int_T\sum_{z=0}^{15}e_zD_z(R)^\dagger D_z(R)\,\dd\mu_2(R).
 \label{eq:F-K3}
\end{equation}
\bi{Using \eqref{eq:F-three-certificate}, $\dd\mu_3=(3/2)\dd\mu_2$, the three zero wedge identities, and the nonnegativity of even degrees $z\ge16$, the exact three-body term satisfies}{使用式~\eqref{eq:F-three-certificate}、$\dd\mu_3=(3/2)\dd\mu_2$、三个零外积恒等式，以及 $z\ge16$ 时偶数阶项的非负性，精确三体项满足}
\begin{equation}
 C_3\ge K_3^{(0)}-\delta_\Lambda H.
 \label{eq:F-three-comparison}
\end{equation}

\begin{lemma}[\bi{The finite Gram error is relative to $H$}{有限 Gram 误差相对于 $H$ 的控制}]
\label{lem:F-metric}
\bi{Let $n_\rho$ be the number of nonzero entries in row $\rho$, and define}{令 $n_\rho$ 为第 $\rho$ 行非零项的个数，并定义}
\begin{equation}
 C_\alpha^{\mathrm F}
 =\sum_{\rho=1}^{7}n_\rho
        \max_{0\le i\le8}
        \sum_{p+j-t_\rho=i}(\alpha_{pj}^{\rho})^2
 =\frac{297842353828763087}{16800000000000000}.
 \label{eq:F-Calpha}
\end{equation}
\bi{Then, as quadratic-form inequalities on the fermionic Fock space,}{则在费米 Fock 空间上，以二次型不等式的意义有}
\begin{equation}
 -C_\alpha^{\mathrm F}e_8(1+e_8)H
 \le R_{3,\mathrm{metric}}
 \le C_\alpha^{\mathrm F}e_8(1+e_8)H.
 \label{eq:F-metric-bound}
\end{equation}
\end{lemma}
\begin{proof}
\bi{For a fixed row and center, put $X_i=\sum_{p+j-t_\rho=i}\alpha_{pj}^{\rho}c_jB_p$. The metric remainder is $\sum_{i,k}(g_{ik}-\delta_{ik})X_i^\dagger X_k$. Since $\|g-I\|\le e_8$, its absolute quadratic form on $\psi$ is at most $e_8\sum_i\|X_i\psi\|^2$. Cauchy--Schwarz gives}{固定一行及一个中心，令 $X_i=\sum_{p+j-t_\rho=i}\alpha_{pj}^{\rho}c_jB_p$。度量余项为 $\sum_{i,k}(g_{ik}-\delta_{ik})X_i^\dagger X_k$。由于 $\|g-I\|\le e_8$，其在 $\psi$ 上的二次型绝对值至多为 $e_8\sum_i\|X_i\psi\|^2$。由 Cauchy--Schwarz 不等式，}
\begin{align}
 \sum_i\|X_i\psi\|^2
 &\le\left(\max_i\sum_{p+j-t_\rho=i}(\alpha_{pj}^{\rho})^2\right)
             \sum_{(p,j)\in\operatorname{supp}(a^\rho)}
                            \|c_jB_p\psi\|^2\notag\\
 &\le(1+e_8)
       \left(\max_i\sum_{p+j-t_\rho=i}(\alpha_{pj}^{\rho})^2\right)
             \sum_{(p,j)\in\operatorname{supp}(a^\rho)}\|B_p\psi\|^2.
 \label{eq:F-X-bound}
\end{align}
\bi{The second inequality uses $\|c_j\|^2=g_{jj}\le1+e_8$. Integrating each term with $\dd\mu_2$ gives $\langle\psi,H\psi\rangle$ exactly; there are $n_\rho$ terms. Summing over rows proves \eqref{eq:F-metric-bound}. Finally, squaring \eqref{eq:F-row-decode} removes every radical, so the value in \eqref{eq:F-Calpha} is computed using rational arithmetic alone.}{第二个不等式使用 $\|c_j\|^2=g_{jj}\le1+e_8$。按 $\dd\mu_2$ 积分后，每项恰给出 $\langle\psi,H\psi\rangle$，而该行共有 $n_\rho$ 项。对全部行求和即得~\eqref{eq:F-metric-bound}。最后，对式~\eqref{eq:F-row-decode}取平方消去全部根式，因此式~\eqref{eq:F-Calpha}中的数值完全由有理算术得到。}
\end{proof}

\subsection{\bi{Kernel-preserving four-body transfer and the count 1222}{保核四体迁移及计数常数 1222}}
\label{subsec:F-four-transfer}

\bi{The four-body term \eqref{eq:F-normal4} has no contraction and therefore no local Gram correction. Change each fixed-$T$ formal input block to the coupled basis \eqref{eq:F-S}. Integrating COM descendants by Lemma~\ref{lem:ladder} leaves its copy matrix $E_D$. If $w_{D,r}^T(R)=(\bigwedge^4C_R)w_{D,r}$, this gives the exact identity}{四体项~\eqref{eq:F-normal4}没有收缩，因此也没有局部 Gram 修正。将每个固定 $T$ 的形式输入分块变换到耦合基~\eqref{eq:F-S}，再按引理~\ref{lem:ladder}积分质心后裔，留下的副本矩阵为 $E_D$。令 $w_{D,r}^T(R)=(\bigwedge^4C_R)w_{D,r}$，则有精确恒等式}
\begin{align}
 K_4&=\int_T\sum_{D=1}^{23}\sum_{r,s\in\mathcal R_D}
        (E_D)_{rs}A(w_{D,r}^T(R))^\dagger A(w_{D,s}^T(R))\,\dd\mu_2(R)
        \notag\\
    &=\frac12\int_T\sum_{D=1}^{23}\sum_{r,s\in\mathcal R_D}
        (E_D)_{rs}A(w_{D,r}^T(R))^\dagger A(w_{D,s}^T(R))\,\dd\mu_4(R).
 \label{eq:F-K4}
\end{align}
\bi{At copy index $r=1$, the second two-particle factor also has relative degree one. The remaining coupling has relative degree $d=D-1$, and its coefficient is the normalized binomial coefficient in \eqref{eq:F-Cd}. Therefore}{当副本指标 $r=1$ 时，第二个双粒子因子的相对阶数同样为一。余下耦合的相对阶数为 $d=D-1$，系数恰为式~\eqref{eq:F-Cd}中的归一化二项式系数。因此}
\begin{equation}
 A(w_{D,1}^T(R))=C_{D-1}(R),\qquad
 T_4=2\int_T\sum_{D=1}^{23}
                  A(w_{D,1}^T(R))^\dagger A(w_{D,1}^T(R))\,\dd\mu_2(R).
 \label{eq:F-target-copy}
\end{equation}
\bi{Both $\dd\mu_4=2\dd\mu_2$ and $D=d+1$ are needed here; together they produce $2\Pi_1$ in the certificate. For even $D$, the $r=1$ column is antisymmetric under exchange of its two pairs, while wedging two degree-two vectors is symmetric, so that column vanishes. The formulas do not require this column, or any other column, to be nonzero or independent.}{这里必须同时使用 $\dd\mu_4=2\dd\mu_2$ 和 $D=d+1$，二者共同产生证书中的 $2\Pi_1$。若 $D$ 为偶数，$r=1$ 列关于两对交换反对称，而两个二阶外积向量的外积关于该交换对称，故该列为零。以上公式并不要求此列或其他列非零，也不要求列向量线性无关。}

\bi{To transfer the certificate, let $A_D=2\Pi_1-E_D+\varepsilon_{\mathrm F}I$. The inequalities $G_DA_DG_D\ge0$ and $W_DA_DW_D^\dagger\ge0$ are equivalent: both test the same quadratic form on $\operatorname{ran}(G_D)=\operatorname{ran}(W_D^\dagger)$. Applying the finite linear map $\bigwedge^4C_R$ by congruence preserves positivity exactly. Applying the positive lift $\mathcal L_4$ and integrating then yields}{为迁移证书，令 $A_D=2\Pi_1-E_D+\varepsilon_{\mathrm F}I$。不等式 $G_DA_DG_D\ge0$ 与 $W_DA_DW_D^\dagger\ge0$ 等价，因为二者检验的是 $\operatorname{ran}(G_D)=\operatorname{ran}(W_D^\dagger)$ 上的同一个二次型。对其施以有限线性映射 $\bigwedge^4C_R$ 的合同变换，会精确保留正性。再应用保正提升 $\mathcal L_4$ 并积分，即得}
\begin{equation}
 K_4\le T_4+\varepsilon_{\mathrm F}R_4,\qquad
 R_4=\int_T\sum_{D=1}^{23}\sum_{r\in\mathcal R_D}
           A(w_{D,r}^T(R))^\dagger A(w_{D,r}^T(R))\,\dd\mu_2(R).
 \label{eq:F-R4}
\end{equation}
\bi{In particular, the proof never inverts $G_D$ and never perturbs its zero eigenvalues. It transports a positive physical four-body operator, rather than asserting positivity of $A_D$ on the entire copy space.}{特别地，证明既不求 $G_D$ 的逆，也不扰动其零本征值；迁移的是正的物理四体算符，而非声称 $A_D$ 在整个副本空间上非负。}

\bi{The allowance $R_4$ also has a particle-number-independent pair-energy bound. For fixed $D$, the coefficient rows $S_r(D,D;p,j,k)$ are orthonormal in the formal input coordinates $(p,j,k)$ with $p+j+k=D$ and $j<k$. This follows by applying the two orthogonal changes of variables in \eqref{eq:F-S} in succession. Bessel's inequality remains valid when those coordinates are arbitrary Hilbert-space vectors, so}{余项 $R_4$ 同样具有与粒子数无关的对能量界。固定 $D$ 时，系数行 $S_r(D,D;p,j,k)$ 在满足 $p+j+k=D$、$j<k$ 的形式输入坐标 $(p,j,k)$ 中正交归一；这可由依次施行式~\eqref{eq:F-S}中的两个正交变量变换得到。将这些坐标替换为任意 Hilbert 空间向量后，Bessel 不等式仍然成立，因此}
\begin{align}
 \sum_{r\in\mathcal R_D}\|A(w_{D,r}^T(R))\psi\|^2
 &\le\sum_{\substack{p+j+k=D\\j<k}}
                       \|c_k(R)c_j(R)B_p(R)\psi\|^2\notag\\
 &\le(1+e_{24})^2
       \sum_{\substack{p+j+k=D\\j<k}}\|B_p(R)\psi\|^2.
 \label{eq:F-Bessel-R4}
\end{align}
\bi{The last step uses the one-particle CAR norm bound twice. It does not assume that physical torus jets or their wedge images are orthogonal. For $n=j+k\ge1$, there are exactly $\lceil n/2\rceil$ nonnegative pairs $j<k$ with that sum, and $p=D-n$ is determined. Hence}{最后一步两次使用单粒子 CAR 范数界，并不假设物理环面 jets 或它们的外积像正交。对 $n=j+k\ge1$，满足该和的非负整数对 $j<k$ 恰有 $\lceil n/2\rceil$ 个，而 $p=D-n$ 随之确定。因此}
\begin{equation}
 n_D=\sum_{n=1}^{D}\left\lceil\frac n2\right\rceil
     =\left\lfloor\frac{(D+1)^2}{4}\right\rfloor,
 \qquad \sum_{D=1}^{23}n_D=1222.
 \label{eq:F-1222-count}
\end{equation}
\bi{Integrating every pair term in \eqref{eq:F-Bessel-R4} with the exact pair-frame identity proves}{对式~\eqref{eq:F-Bessel-R4}中的每个对项应用精确对框架恒等式并积分，即得}
\begin{equation}
 R_4\le1222(1+e_{24}(\Lambda))^2H,
 \qquad
 K_4\le T_4+1222\varepsilon_{\mathrm F}(1+e_{24}(\Lambda))^2H.
 \label{eq:F-four-final}
\end{equation}

\subsection{\bi{Closing the fermionic gap estimate}{完成费米谱隙估计}}
\label{subsec:F-budget}

\bi{We collect the estimates as quadratic forms. The exact normal-ordering identity is $H^2=\mathcal L_2(h_T^2)+C_3+S_4$. The all-COM contact-frame estimate gives $\mathcal L_2(h_T^2)\ge(1-\varepsilon_2^{\mathrm F}(\Lambda))H$; it does not replace $h_T^2$ by $h_T$. Combining this bound with \eqref{eq:F-three-comparison} and \eqref{eq:F-four-payment} yields}{现以二次型形式汇总各估计。精确正规序恒等式为 $H^2=\mathcal L_2(h_T^2)+C_3+S_4$。覆盖全部质心模式的接触框架估计给出 $\mathcal L_2(h_T^2)\ge(1-\varepsilon_2^{\mathrm F}(\Lambda))H$，并未将 $h_T^2$ 直接换为 $h_T$。将此界与式~\eqref{eq:F-three-comparison}和~\eqref{eq:F-four-payment}合并，得到}
\begin{align}
 H^2\ge K_3^{(0)}+T_4+
  \bigl[1-\varepsilon_2^{\mathrm F}-\delta_\Lambda
                  -C_4^{\mathrm F}\varepsilon_4^{\mathrm F}\bigr]H.
 \label{eq:F-budget-first}
\end{align}
\bi{Next, \eqref{eq:F-K-decomposition}, \eqref{eq:F-metric-bound}, and \eqref{eq:F-four-final} imply}{接着，由式~\eqref{eq:F-K-decomposition}、\eqref{eq:F-metric-bound}及~\eqref{eq:F-four-final}，}
\begin{equation}
 K_3^{(0)}+T_4
 \ge K_T-
 \bigl[\eta+C_\alpha^{\mathrm F}e_8(1+e_8)
              +1222\varepsilon_{\mathrm F}(1+e_{24})^2\bigr]H.
 \label{eq:F-budget-second}
\end{equation}
\bi{Consequently}{因此}
\begin{align}
 H^2&\ge K_T+\Gamma_{\mathrm F}(\Lambda)H
       \ge\Gamma_{\mathrm F}(\Lambda)H,\label{eq:F-final-gap}\\
 \Gamma_{\mathrm F}(\Lambda)
 &=1-\eta-\delta_\Lambda
      -1222\varepsilon_{\mathrm F}(1+e_{24}(\Lambda))^2
      -\varepsilon_2^{\mathrm F}(\Lambda)\notag\\
 &\hspace{1.5em}-C_\alpha^{\mathrm F}e_8(\Lambda)(1+e_8(\Lambda))
      -C_4^{\mathrm F}(\Lambda)\varepsilon_4^{\mathrm F}(\Lambda).
 \label{eq:F-Gamma}
\end{align}
\bi{Every term in this scalar budget has been paid for by the same operator $H$. Fixed-degree Gram errors and the two finite relative-frame errors tend to zero uniformly as the actual shortest nonzero lattice vector tends to infinity. Equation~\eqref{eq:F-tail-limit} gives the remaining limit. Therefore}{此标量预算中的每一项均以同一个算符 $H$ 支付。随着实际最短非零格矢趋于无穷，固定阶 Gram 误差及两个有限相对框架误差均一致趋于零；剩余尾项的极限由式~\eqref{eq:F-tail-limit}给出。因此}
\begin{equation}
 \lim_{\ell(\Lambda)\to\infty}\Gamma_{\mathrm F}(\Lambda)
 =1-\frac{93527408868499}{10^{14}}-\frac{61}{4096}
       -1222\frac{3}{10^6}
 =\frac{4616733319001}{10^{14}}
 =:\gamma_{\mathrm F}^{*}.
 \label{eq:F-gamma-star}
\end{equation}
\bi{Given $0<\gamma<\gamma_{\mathrm F}^{*}$, uniform convergence supplies a systole threshold $L_\gamma$ for which $\Gamma_{\mathrm F}(\Lambda)\ge\gamma$. Then \eqref{eq:F-final-gap} proves $H^2\ge\gamma H$ on every particle sector and hence on the full fermionic Fock space. A positive eigenvalue $E$ satisfies $E^2\ge\gamma E$, so $E\ge\gamma$. This proves the gap assertion of Theorem~\ref{thm:fermion}. The argument supplies an asymptotic threshold, not an explicit numerical value of $L_\gamma$.}{任给 $0<\gamma<\gamma_{\mathrm F}^{*}$，一致收敛保证存在最短周期阈值 $L_\gamma$，使 $\Gamma_{\mathrm F}(\Lambda)\ge\gamma$。于是式~\eqref{eq:F-final-gap}在每个粒子扇区乃至整个费米 Fock 空间上给出 $H^2\ge\gamma H$。任意正本征值 $E$ 均满足 $E^2\ge\gamma E$，故 $E\ge\gamma$，从而证明定理~\ref{thm:fermion}的谱隙结论。本论证给出渐近阈值的存在性，而未提供 $L_\gamma$ 的显式数值。}

\bi{In particular, $1/25<\gamma_{\mathrm F}^{*}$. At $M=3N$, the common zero-mode argument identifies the three-dimensional Laughlin ground space, and the same inequality gives a gap at least $1/25$ once the systole is sufficiently large. Particle number, magnetic flux, and twist do not enter the error constants. The geometric condition remains essential: neither increasing area alone nor choosing two long but nearly parallel lattice generators ensures a large shortest cycle.}{特别地，$1/25<\gamma_{\mathrm F}^{*}$。当 $M=3N$ 时，公共零模论证确定了三维 Laughlin 基态空间，而同一不等式在最短周期充分大后给出至少 $1/25$ 的谱隙。粒子数、磁通量和 twist 均不进入误差常数。但几何条件不可省略：仅面积增大，或选取两个很长却近乎平行的格子生成元，都不能保证最短周期很长。}

\section{\bi{The bosonic bound: positive comparison and a finite certificate}{玻色不等式：正比较与有限证书}}
\label{sec:B-bosonic}

\bi{We now prove the bosonic assertion of Theorem~\ref{thm:boson}. Throughout this section, $H=H_B$ denotes the full contact interaction, with planar nonzero two-particle eigenvalue one. We also abbreviate the bosonic connected and zero-contraction terms by $S_3$ and $S_4$, and write $H_k^{\rm pl}=H_{k,B}^{\rm pl}$. The exact identities established in the preceding sections reduce the problem to controlling the negative three-body tail and the errors of a finite sum-of-squares certificate. Bosonic occupation numbers are unbounded, so those errors must be compared with the positive operators $H$ and $S_4$, rather than with a particle-number-dependent multiple of the identity. We obtain that comparison from coherent-state occupation inequalities and positive Gaussian kernels.}
{本节证明定理~\ref{thm:boson} 的玻色结论。全节记 $H=H_B$ 为完整接触相互作用，并取平面两粒子非零本征值为一的归一化。同时将玻色连通项和零收缩项简称为 $S_3$ 和 $S_4$，并记 $H_k^{\rm pl}=H_{k,B}^{\rm pl}$。前文的精确恒等式将问题化为控制三体负尾和有限平方和证书的误差。玻色占据数无界，因此这些误差须与正算符 $H$ 和 $S_4$ 比较，而不能与依赖粒子数的恒等算符倍数比较。我们通过相干态占据数不等式和正高斯核建立所需比较。}

\bi{All operator inequalities in the proof are first understood on a fixed particle sector. The physical torus lowest-Landau-level space is finite dimensional, and every such sector is finite dimensional. The passage to the full bosonic Fock space is made at the end of the section. The periodized orbitals remain the raw jets of Section~\ref{sec:geometry}; none is orthogonalized or pointwise renormalized in the definitions of the interaction or the certificate.}
{证明中的算符不等式先在固定粒子数扇区上理解。物理环面最低朗道能级空间有限维，每个固定粒子数扇区也有限维。到全玻色 Fock 空间的延拓在本节末尾完成。周期化轨道始终采用第~\ref{sec:geometry} 节的原始 jets；定义相互作用和证书时既不将其正交化，也不逐点重新归一化。}

\subsection{\bi{Coherent occupation inequalities}{相干态占据数不等式}}
\label{sec:B-occupation}
\bi{The coherent-occupation comparison is adapted from the bosonic spherical candidate of Shen~\cite{Shen2026BosonicCandidate}. The constants $k_1,k_2,c,s$ below are torus reproducing-kernel quantities. Their bounds and the positive thermal transfer must therefore be established for this geometry, independently of that source's finite-sphere interval verification.}
{相干态占据数比较借鉴Shen的玻色球面候选证明~\cite{Shen2026BosonicCandidate}。下文的$k_1,k_2,c,s$是环面再生核给出的量，故须在本几何中另证其界及正热核迁移，不能从来源的有限球面区间验证直接推出。}

\bi{Let $K_j^T$ be the charge-$j$ parameter reproducing kernel, and set}
{记 $K_j^T$ 为电荷 $j$ 参数空间的再生核，并置}
\begin{equation}
 k_j(R)=K_j^T(R,R),\qquad
 s=\sup_{R\in T}k_1(R),\qquad
 c=\sup_{R\in T}\frac{k_2(R)}{k_1(R)^2}.
 \label{eq:B-density-constants}
\end{equation}
\bi{In the large-systole regime used below, $k_1$ is strictly positive. Write $B_R=B_0^B(R)$ and $\beta_R=u_0^T(R)/\sqrt{k_1(R)}$. The latter orbital is normalized only for the purpose of applying the following occupation estimates.}
{在下文使用的大最短周期范围内，$k_1$ 严格为正。\par 记 $B_R=B_0^B(R)$，并令 $\beta_R=u_0^T(R)/\sqrt{k_1(R)}$。这里对轨道的归一化仅用于下述占据数估计。}

\begin{lemma}[\bi{One and two coherent spectators}{一个和两个相干旁观粒子}]
\label{lem:B-occupation}
\bi{For a normalized coherent orbital $\beta=\beta_R$, let $n_\beta=a(\beta)^\dagger a(\beta)$. Then}
{对规范化相干轨道 $\beta=\beta_R$，记 $n_\beta=a(\beta)^\dagger a(\beta)$。则}
\begin{equation}
 \frac12n_\beta(n_\beta-1)\le cH,\qquad
 n_\beta\le cH+I.
 \label{eq:B-one-occupation}
\end{equation}
\bi{For any two such orbitals $\beta,\gamma$, including nonorthogonal ones,}
{对任意两个这样的轨道 $\beta,\gamma$，包括不正交的情形，均有}
\begin{equation}
 \bigl(a(\gamma)a(\beta)\bigr)^\dagger
       a(\gamma)a(\beta)\le4cH+2I.
 \label{eq:B-two-occupation}
\end{equation}
\end{lemma}

\begin{proof}
\bi{The two-particle evaluation map $Q$ has $h_T=Q^\dagger Q$, and $Q\psi$ is a charge-two parameter lowest-Landau-level section. Its reproducing property gives $|Q\psi(R)|^2\le k_2(R)\|Q\psi\|^2$. Consequently, $|v_0(R)\rangle\langle v_0(R)|\le k_2(R)h_T$. Applying the positive lift $\mathcal L_2$ yields $B_R^\dagger B_R\le k_2(R)H$. Since}
{两粒子求值映射 $Q$ 满足 $h_T=Q^\dagger Q$，且 $Q\psi$ 是电荷二参数最低朗道能级截面。再生性质给出 $|Q\psi(R)|^2\le k_2(R)\|Q\psi\|^2$，故 $|v_0(R)\rangle\langle v_0(R)|\le k_2(R)h_T$。施加正提升 $\mathcal L_2$ 得 $B_R^\dagger B_R\le k_2(R)H$。又因为}
\begin{equation}
 B_R=\frac{k_1(R)}{\sqrt2}\,a(\beta_R)^2,
 \qquad
 B_R^\dagger B_R
 =\frac{k_1(R)^2}{2}n_{\beta_R}(n_{\beta_R}-1),
 \label{eq:B-normalized-coherent}
\end{equation}
\bi{the first inequality in \eqref{eq:B-one-occupation} follows. The scalar inequality $n\le n(n-1)/2+1$ for every nonnegative integer $n$ gives the second inequality by the spectral calculus for $n_\beta$. In particular, the same statement applies in the output sector after $B_R$ has removed two particles.}
{式~\eqref{eq:B-one-occupation} 的第一个不等式随即成立。对每个非负整数 $n$，有标量不等式 $n\le n(n-1)/2+1$；对 $n_\beta$ 作谱演算即得第二个不等式。特别地，它同样适用于 $B_R$ 移除两个粒子后的输出扇区。}

\bi{For the two-orbital estimate, put $A=a(\gamma)a(\beta)$ and $\omega=\langle\beta,\gamma\rangle$, so $|\omega|\le1$. The canonical commutation relations give}
{为证明双轨道估计，令 $A=a(\gamma)a(\beta)$，并记 $\omega=\langle\beta,\gamma\rangle$，故 $|\omega|\le1$。正则对易关系给出}
\begin{align}
 A^\dagger A
 &=\frac12(n_\beta n_\gamma+n_\gamma n_\beta)
    -\operatorname{Re}\!\left(\omega a(\beta)^\dagger a(\gamma)\right)
 \notag\\
 &\le\frac12(n_\beta^2+n_\gamma^2)
       +\frac12(n_\beta+n_\gamma)
 \le4cH+2I.
 \label{eq:B-two-occupation-proof}
\end{align}
\bi{The first inequality uses $(n_\beta-n_\gamma)^2\ge0$ and quadratic-form Cauchy--Schwarz on the cross term. For the second, apply $n_x^2\le2cH+n_x$ and then $n_x\le cH+I$ for $x=\beta,\gamma$. This proves the lemma without a bound on an annihilation operator's norm.}
{第一个不等式使用 $(n_\beta-n_\gamma)^2\ge0$，并对交叉项使用二次型 Cauchy--Schwarz 不等式。第二个不等式先对 $x=\beta,\gamma$ 使用 $n_x^2\le2cH+n_x$，再使用 $n_x\le cH+I$。因此整个证明不需要假设湮灭算符具有一致有界的范数。}
\end{proof}

\bi{Replacing a normalized spectator by the raw coherent jet introduces at most a factor $s$; replacing two spectators introduces at most $s^2$. No additional factor is attached to the pair anchor. Indeed, its subsequent integration is governed by the exact identities}
{将一个规范化旁观轨道换回原始相干 jet，至多引入因子 $s$；两个旁观轨道则至多引入 $s^2$。配对锚点不再附加任何因子，因为随后对它的积分由以下精确恒等式控制：}
\begin{equation}
 \int_T B_\alpha^\dagger B_\alpha\,\dd\mu_2(\alpha)=H,
 \qquad
 \int_T B_\alpha^\dagger H B_\alpha\,\dd\mu_2(\alpha)=S_4.
 \label{eq:B-anchor-identities}
\end{equation}
\bi{The second identity follows by substituting the coherent integral for $H$ and commuting the pair annihilators. Thus the occupation estimates pay for an arbitrary spectator population using precisely the positive zero-contraction term already present in $H^2$.}
{第二个恒等式通过代入 $H$ 的相干态积分并交换配对湮灭算符得到。因此，占据数估计恰好利用 $H^2$ 中已有的正零收缩项来控制任意大的旁观粒子占据数。}

\subsection{\bi{Positive thermal kernels and their exact torus transfer}{正热核及其精确环面移植}}
\bi{The planar coherent-comparison input is adapted from Shen~\cite{Shen2026BosonicCandidate}. The exact transfer of positive thermal kernels to the magnetic quotient is established below using the convergent domain of Section~\ref{sec:geometry}; it is not a consequence of the source's spherical Schur averages.}
{平面相干态比较输入借鉴Shen~\cite{Shen2026BosonicCandidate}。下文利用第~\ref{sec:geometry}节的收敛域建立正热核到磁商空间的精确迁移；这一迁移不是来源中球面Schur平均的直接推论。}
\label{sec:B-thermal}

\bi{The auxiliary spaces in this subsection are planar spaces. Let $U$ be the planar one-particle Bargmann space and let $V$ be the pair subspace with relative degree zero, with orthonormal center-of-mass basis $v_p$. A normalized coherent vector $|\sqrt2\alpha\rangle_V$ represents the pair vector $B_\alpha^\dagger|0\rangle$ on the plane. The insertion maps are}
{本小节中的辅助空间均为平面空间。记 $U$ 为平面单粒子 Bargmann 空间，$V$ 为相对次数为零的配对子空间，其质心正交基为 $v_p$。规范化相干向量 $|\sqrt2\alpha\rangle_V$ 表示平面上的配对向量 $B_\alpha^\dagger|0\rangle$。定义插入映射}
\begin{align}
 W_3 &:V\otimes U\longrightarrow\Sym^3U,
 & W_3(v_p\otimes u)&=B_p^\dagger u,
 \notag\\
 W_4 &:V\otimes\Sym^2U\longrightarrow\Sym^4U,
 & W_4(v_p\otimes u)&=B_p^\dagger u.
 \label{eq:B-insertion-maps}
\end{align}
\bi{These maps use the normalized symmetric-tensor to Fock-space identification. In particular, $W_3=\sqrt3P_{\rm sym}\iota_3$, where $\iota_3$ is the isometric embedding of the distinguished pair and the spectator before symmetrization. Moreover,}
{这些映射采用规范化的对称张量到 Fock 空间识别。特别地，$W_3=\sqrt3P_{\rm sym}\iota_3$，其中 $\iota_3$ 是对称化之前将指定配对与旁观粒子嵌入的等距映射。此外，}
\begin{equation}
 W_kW_k^\dagger=H_k^{\rm pl},\qquad
 \|W_3\|\le\sqrt3,\qquad \|W_4\|\le\sqrt6.
 \label{eq:B-insertion-norms}
\end{equation}
\bi{The first identity is the sum over all physical pairs. The norm bounds follow because $H_k^{\rm pl}$ is the sum of $\binom{k}{2}$ pair projections. Each $W_k$ intertwines the total lowering and raising operators, since inserting a pair is compatible with simultaneous magnetic translation. Hence in the planar center-of-mass decomposition it has the form $I_{\CM}\otimes W_{k,\rm rel}$, including the full relative multiplicity space.}
{第一个恒等式来自对所有物理粒子对求和。范数上界则因为 $H_k^{\rm pl}$ 是 $\binom{k}{2}$ 个配对投影之和。每个 $W_k$ 都交织总降阶和升阶算符，因为插入配对与同时磁平移相容。因此，在平面质心分解下它具有 $I_{\CM}\otimes W_{k,\rm rel}$ 的形式，并完整保留相对空间的重数。}

\bi{For clarity, the center-of-mass factors used in the transfer are recorded explicitly. The planar twirl of a trace-class anchored seed $\sigma$ and its torus counterpart are}
{为明确移植所用的质心因子，列出迹类锚定种子 $\sigma$ 的平面平均及其环面对应式：}
\begin{align}
 \int_{\mathbb C}T_\alpha^{\otimes k}\sigma
      T_\alpha^{\otimes k\dagger}\,\dd\mu_2(\alpha)
 &=\frac2k I_{\CM}\otimes\Tr_{\CM}\sigma,
 \label{eq:B-planar-twirl}\\
 \int_T C_\alpha^{(k)}\sigma C_\alpha^{(k)\dagger}
                  \,\dd\mu_2(\alpha)
 &=\Omega_k\!\left(\frac2k\Tr_{\CM}\sigma\right).
 \label{eq:B-torus-seed}
\end{align}
\bi{For finite-degree seeds, the second identity is the integrated ladder identity of Lemma~\ref{lem:ladder}. The factor $2/k$ in the first identity comes from the normalized center-of-mass displacement $h=\sqrt{k}\alpha$, so that $\dd\mu_2(\alpha)=(2/k)\dd^2h/\pi$. The Weyl average traces out the center-of-mass oscillator and leaves the relative operator unchanged, rather than replacing it by a scalar. We justify the extension of \eqref{eq:B-torus-seed} to the thermal seeds below. No global relative-coordinate decomposition on the torus is assumed.}
{对有限次数种子，第二个恒等式即引理~\ref{lem:ladder} 的积分梯级恒等式。第一个恒等式中的 $2/k$ 来自规范化质心位移 $h=\sqrt{k}\alpha$，此时 $\dd\mu_2(\alpha)=(2/k)\dd^2h/\pi$。Weyl 平均对质心振子取偏迹，保留相对算符本身，而不是将其替换为标量。下文将证明式~\eqref{eq:B-torus-seed} 对热种子的延拓。整个论证不假设环面上存在全局相对坐标分解。}

\subsubsection{\bi{One spectator}{一个旁观粒子}}
\label{sec:B-thermal-one}

\bi{For $a>0$, define a positive one-particle seed by}
{对 $a>0$，定义正的单粒子种子}
\begin{equation}
 R_a=\int_{\mathbb C}e^{-a|\beta|^2}
        |u_0(\beta)\rangle\langle u_0(\beta)|\,\dd\mu_1(\beta)
     =\sum_{j=0}^{\infty}(1+a)^{-j-1}|e_j\rangle\langle e_j|.
 \label{eq:B-thermal-seed}
\end{equation}
\bi{The equality follows by expanding the normalized coherent vector in the basis $e_j$: the angular integral makes the matrix diagonal, and the radial Gaussian moment gives $(1+a)^{-j-1}$. In particular, $R_a$ is positive and trace class. Set}
{将规范化相干向量在基 $e_j$ 中展开即可证明等式：角向积分消去非对角项，径向高斯矩给出 $(1+a)^{-j-1}$。因此 $R_a$ 为正迹类算符。令}
\begin{equation}
 \sigma_3=|0\rangle_V\langle0|\otimes R_a,
 \qquad \widetilde\sigma_3=W_3\sigma_3W_3^\dagger.
 \label{eq:B-three-seed}
\end{equation}
\bi{In the planar simultaneous-translation integral, introduce the coherent coordinates $u=\sqrt2\alpha$ and $v=\alpha+\beta$, followed by the unitary Jacobi change of variables}
{在平面同时平移积分中，先引入相干坐标 $u=\sqrt2\alpha$ 和 $v=\alpha+\beta$，再作酉 Jacobi 变量变换}
\begin{equation}
 \zeta=\frac{\sqrt2u+v}{\sqrt3},\qquad
 \xi=\frac{-u+\sqrt2v}{\sqrt3}.
 \label{eq:B-three-jacobi}
\end{equation}
\bi{The measure is $\dd\mu_2(\alpha)\dd\mu_1(v)=\dd^2u\dd^2v/\pi^2$, and the Gaussian weight becomes $\exp[-(3a/2)|\xi|^2]$. The $\zeta$ integral resolves the center-of-mass identity, while \eqref{eq:B-thermal-seed} evaluates the $\xi$ integral. Thus, with $r=2/(2+3a)$, the auxiliary twirl is exactly $I_{\CM}\otimes r^{n_\xi+1}$. Intertwining with $W_3$ and comparing with \eqref{eq:B-planar-twirl} gives}
{测度为 $\dd\mu_2(\alpha)\dd\mu_1(v)=\dd^2u\dd^2v/\pi^2$，高斯权重变为 $\exp[-(3a/2)|\xi|^2]$。对 $\zeta$ 的积分给出质心恒等算符的分解，对 $\xi$ 的积分则由式~\eqref{eq:B-thermal-seed} 求得。因此，令 $r=2/(2+3a)$，辅助空间中的平均恰为 $I_{\CM}\otimes r^{n_\xi+1}$。利用 $W_3$ 的交织性质并与式~\eqref{eq:B-planar-twirl} 比较，得到}
\begin{equation}
 \frac23\Tr_{\CM}\widetilde\sigma_3
   =W_{3,\rm rel}r^{n_\xi+1}W_{3,\rm rel}^\dagger.
 \label{eq:B-three-partial-trace}
\end{equation}

\bi{Apply $\Omega_3$ to this planar relative coefficient and use \eqref{eq:B-torus-seed}. Substituting the positive coherent integral for $R_a$ then gives an explicit physical coefficient. Indeed, simultaneous translation by $\alpha$ sends the spectator coherent vector at offset $\beta$ to the periodized vector at $\alpha+\beta$, multiplied by $\exp[-i\operatorname{Im}(\alpha\overline\beta)]$. This phase cancels in the rank-one projector. Positive lifting therefore gives the exact identity}
{对这一平面相对系数施加 $\Omega_3$，并使用式~\eqref{eq:B-torus-seed}。随后将 $R_a$ 的正相干积分代入，即得显式物理系数。具体而言，同时平移 $\alpha$ 将偏移量为 $\beta$ 的旁观相干向量变为位于 $\alpha+\beta$ 的周期化向量，并乘上相位 $\exp[-i\operatorname{Im}(\alpha\overline\beta)]$。该相位在秩一投影中抵消。因此，正提升给出精确恒等式}
\begin{align}
 \mathcal G_3(a)
 &:={\mathcal L}_3\Omega_3
       \left(W_{3,\rm rel}r^{n_\xi+1}W_{3,\rm rel}^\dagger\right)
 \notag\\
 &=\int_{\alpha\in T}\dd\mu_2(\alpha)
   \int_{\beta\in\mathbb C}\dd\mu_1(\beta)\,
   e^{-a|\beta|^2}B_\alpha^\dagger
   a\bigl(u_0^T(\alpha+\beta)\bigr)^\dagger
   a\bigl(u_0^T(\alpha+\beta)\bigr)B_\alpha.
 \label{eq:B-one-spectator-transfer}
\end{align}
\bi{The spectator offset ranges over the whole plane. Restricting that integral to one torus cell without retaining its images would change the coefficient and is not part of the transfer. Lemma~\ref{lem:B-occupation} bounds the spectator factor by $s(cH+I)$, and $\int_{\mathbb C}e^{-a|\beta|^2}\dd\mu_1(\beta)=1/a$. The exact anchor identities \eqref{eq:B-anchor-identities} imply}
{旁观粒子的偏移量遍历整个平面。若将积分限制在环面的单个基胞而不保留其镜像，就会改变系数；移植中并不进行这种替换。引理~\ref{lem:B-occupation} 将旁观因子控制为 $s(cH+I)$，且 $\int_{\mathbb C}e^{-a|\beta|^2}\dd\mu_1(\beta)=1/a$。再由精确锚点恒等式~\eqref{eq:B-anchor-identities} 得}
\begin{equation}
 \mathcal G_3(a)\le\frac{s}{a}(cS_4+H).
 \label{eq:B-three-thermal-bound}
\end{equation}

\subsubsection{\bi{Two spectators}{两个旁观粒子}}
\label{sec:B-thermal-two}

\bi{On $V\otimes\Sym^2U$, use the seed}
{在 $V\otimes\Sym^2U$ 上使用种子}
\begin{equation}
 \sigma_4=|0\rangle_V\langle0|\otimes[R_a\otimes R_a]_{\rm sym},
 \qquad \widetilde\sigma_4=W_4\sigma_4W_4^\dagger,
 \label{eq:B-four-seed}
\end{equation}
\bi{where $[\cdot]_{\rm sym}$ denotes compression to the normalized symmetric tensor space. Set $u=\sqrt2\alpha$, $v=\alpha+\beta$, and $w=\alpha+\gamma$. The unitary Jacobi coordinates are}
{其中 $[\cdot]_{\rm sym}$ 表示压缩到规范化对称张量空间。令 $u=\sqrt2\alpha$、$v=\alpha+\beta$、$w=\alpha+\gamma$。酉 Jacobi 坐标为}
\begin{equation}
 \zeta=\frac{\sqrt2u+v+w}{2},\qquad
 \xi_1=\frac{v-w}{\sqrt2},\qquad
 \xi_2=\frac{u}{\sqrt2}-\frac{v+w}{2}.
 \label{eq:B-four-jacobi}
\end{equation}
\bi{The measure is preserved, and the weight separates because}
{该变换保持测度，而权重因下式分离：}
\begin{equation}
 |\beta|^2+|\gamma|^2=|\xi_1|^2+2|\xi_2|^2.
 \label{eq:B-four-gaussian}
\end{equation}
\bi{The center-of-mass coherent integral gives the identity; the two remaining Gaussian moments give $r_1^{n_1+1}r_2^{n_2+1}$ with $r_1=(1+a)^{-1}$ and $r_2=(1+2a)^{-1}$. Spectator interchange sends $\xi_1$ to $-\xi_1$ and fixes $\xi_2$, so the symmetric spectator space is the subspace with even $n_1$. Equation~\eqref{eq:B-planar-twirl}, now with $k=4$, therefore yields}
{质心相干积分给出恒等算符；其余两个高斯矩给出 $r_1^{n_1+1}r_2^{n_2+1}$，其中 $r_1=(1+a)^{-1}$、$r_2=(1+2a)^{-1}$。交换两个旁观粒子将 $\xi_1$ 变为 $-\xi_1$，并保持 $\xi_2$，故旁观粒子的对称空间对应 $n_1$ 为偶数的子空间。于是，在式~\eqref{eq:B-planar-twirl} 中取 $k=4$ 得}
\begin{equation}
 \frac12\Tr_{\CM}\widetilde\sigma_4
 =W_{4,\rm rel}
       \left[r_1^{n_1+1}r_2^{n_2+1}\right]W_{4,\rm rel}^\dagger.
 \label{eq:B-four-partial-trace}
\end{equation}

\bi{Transfer this relative coefficient by $\Omega_4$ and then lift it. The resulting positive integral is}
{先通过 $\Omega_4$ 移植这一相对系数，再作正提升，所得正积分为}
\begin{align}
 \mathcal G_4(a)
 &:={\mathcal L}_4\Omega_4
   \left(W_{4,\rm rel}
       [r_1^{n_1+1}r_2^{n_2+1}]W_{4,\rm rel}^\dagger\right)
 \notag\\
 &=\frac12\int_{\alpha\in T}\dd\mu_2(\alpha)
       \int_{\beta,\gamma\in\mathbb C}
           \dd\mu_1(\beta)\dd\mu_1(\gamma)\,
       e^{-a(|\beta|^2+|\gamma|^2)}
       B_\alpha^\dagger A_{\alpha;\beta\gamma}^\dagger
                  A_{\alpha;\beta\gamma}B_\alpha,
 \label{eq:B-two-spectator-transfer}\\
 A_{\alpha;\beta\gamma}
 &=a\bigl(u_0^T(\alpha+\gamma)\bigr)
       a\bigl(u_0^T(\alpha+\beta)\bigr).
 \notag
\end{align}
\bi{The factor $1/2$ in this integral has a different origin from the factor $1/2$ in \eqref{eq:B-four-partial-trace}. The latter is the center-of-mass factor $2/k$. The former comes from the symmetric-tensor identification}
{该积分前的 $1/2$ 与式~\eqref{eq:B-four-partial-trace} 中的 $1/2$ 来源不同：后者是质心因子 $2/k$，前者则来自对称张量识别}
\begin{equation}
 P_{\rm sym}^{(2)}\bigl(|\beta\rangle\otimes|\gamma\rangle\bigr)
 \longleftrightarrow
 \frac{a(\beta)^\dagger a(\gamma)^\dagger|0\rangle}{\sqrt2}.
 \label{eq:B-sym-factor}
\end{equation}
\bi{Both factors have therefore been included, once each. The two offsets in \eqref{eq:B-two-spectator-transfer} again range over the entire plane. Applying \eqref{eq:B-two-occupation} to the unnormalized jets costs at most $s^2(4cH+2I)$; the two Gaussian marginals give $1/a^2$. Together with \eqref{eq:B-anchor-identities}, this proves}
{因此，这两个因子均已分别计入一次。式~\eqref{eq:B-two-spectator-transfer} 中的两个偏移量也都遍历整个平面。将式~\eqref{eq:B-two-occupation} 应用于未归一化 jets，至多付出 $s^2(4cH+2I)$；两个高斯边缘积分给出 $1/a^2$。结合式~\eqref{eq:B-anchor-identities}，得到}
\begin{equation}
 \mathcal G_4(a)\le\frac{s^2}{a^2}(2cS_4+H).
 \label{eq:B-four-thermal-bound}
\end{equation}

\subsubsection{\bi{Convergence and preservation of positive order}{收敛与正序的保持}}
\label{sec:B-thermal-convergence}

\bi{The preceding formulas can be obtained without treating infinite periodization as a bounded operator on the entire planar Bargmann space. To see this, truncate each anchored seed by total degree $d$. Before insertion, the degree-$d$ traces obey}
{上述公式的建立不需要将无限周期化视为整个平面 Bargmann 空间上的有界算符。具体而言，按总次数 $d$ 对每个锚定种子作截断。在插入之前，次数 $d$ 的迹满足}
\begin{equation}
 \Tr(\sigma_{3,d})\le(1+a)^{-d-1},\qquad
 \Tr(\sigma_{4,d})\le(d+1)(1+a)^{-d-2}.
 \label{eq:B-seed-decay}
\end{equation}
\bi{Insertion multiplies these bounds by at most $3$ and $6$, respectively, by \eqref{eq:B-insertion-norms}. The raw-jet estimate of Lemma~\ref{lem:jets} bounds the squared norm of $C_R^{(k)}$ on a degree-$d$ block by $C_k(d+1)^{2k}$, uniformly in $R$ on the fixed torus. Consequently,}
{由式~\eqref{eq:B-insertion-norms}，插入映射至多将这两个上界分别乘以 $3$ 和 $6$。引理~\ref{lem:jets} 的原始 jet 估计将 $C_R^{(k)}$ 在次数 $d$ 块上的范数平方控制为 $C_k(d+1)^{2k}$，且在固定环面上对 $R$ 一致。因此，}
\begin{equation}
 \sum_{d=0}^{\infty}\sup_{R\in T}
  \Tr\!\left(C_R^{(k)}\widetilde\sigma_{k,d}
             C_R^{(k)\dagger}\right)<\infty,
 \qquad k=3,4.
 \label{eq:B-thermal-convergence}
\end{equation}
\bi{The fundamental cell has finite area. Thus finite-degree identities pass to the limit in trace norm on the physical few-body space, and the center-of-mass partial trace and coherent integrals may be interchanged with that limit. On the coherent side, boundedness of $\sup_R\|u_0^T(R)\|$ and Gaussian integrability also give direct domination. This proves \eqref{eq:B-one-spectator-transfer} and \eqref{eq:B-two-spectator-transfer} with all offsets retained.}
{基胞面积有限，故有限次数恒等式可在物理少体空间上以迹范数取极限，质心偏迹和相干积分也可与该极限交换。在相干积分一侧，$\sup_R\|u_0^T(R)\|$ 有界以及高斯可积性同样给出直接支配。由此证明了保留全部偏移量的式~\eqref{eq:B-one-spectator-transfer} 和式~\eqref{eq:B-two-spectator-transfer}。}

\bi{Every relative comparison below is diagonal in total relative degree, although its blocks may have nontrivial multiplicity. On each finite collection of blocks, take the positive difference of the two coefficients and apply the congruence defining $\Omega_k$. The resulting operator remains positive. The exponential domination just established then passes this inequality through a positive limit. Hence the thermal comparisons preserve order on the torus without defining $\Omega_k$ on arbitrary nondecaying infinite relative operators.}
{下文的相对比较均按总相对次数分块对角，但块内可以存在非平凡重数。在任意有限块集上，先取两个系数之差这一正矩阵，再施加定义 $\Omega_k$ 的合同变换，所得算符仍为正。刚建立的指数支配随后使该不等式通过正极限。因此，热核比较在环面上保持序关系，而无需对任意不衰减的无限相对算符定义 $\Omega_k$。}

\subsection{\bi{The negative tail and two finite comparison operators}{负尾与两个有限比较算符}}
\label{sec:B-comparisons}

\bi{Let $\xi_z=\xi_z^B$ be the normalized pair--spectator relative oscillator vector of degree $z$, as in the three-body calculation leading to \eqref{eq:boson-threebody}. Put}
{记 $\xi_z=\xi_z^B$ 为相对次数 $z$ 的规范化配对--旁观粒子相对振子向量，与导出式~\eqref{eq:boson-threebody} 的三体计算中的定义一致。令}
\begin{align}
 T_9&=\mathcal L_3\Omega_3
 \left(W_{3,\rm rel}
   \left[2\!\sum_{\substack{z\ge9\\z\ {\rm odd}}}
        2^{-z}|\xi_z\rangle\langle\xi_z|\right]
       W_{3,\rm rel}^\dagger\right),
 \label{eq:B-tail-definition}\\
 R_{k,D}&=\mathcal L_k\Omega_k
       \left(H_{k,\rm rel}^{\rm pl}\Pi_{\le D}^{(k)}\right),
 \qquad (k,D)=(3,8),(4,16).
 \label{eq:B-finite-comparisons}
\end{align}
\bi{The projections in \eqref{eq:B-finite-comparisons} retain every physical relative multiplicity at the indicated degree. They are not projections onto only one selected vector in each degree.}
{式~\eqref{eq:B-finite-comparisons} 中的投影保留所示次数内的全部物理相对重数，并非每个次数只选取一个向量的投影。}

\begin{proposition}[\bi{Bosonic comparison constants}{玻色比较常数}]
\label{prop:B-comparisons}
\bi{The operators just defined satisfy}
{上述算符满足}
\begin{align}
 T_9&\le c_9s(cS_4+H),
 &c_9&=30\left(\frac59\right)^9
       =\frac{19531250}{129140163},
 \label{eq:B-tail-bound}\\
 R_{3,8}&\le C_3^B s(cS_4+H),
 &C_3^B&=12\left(\frac98\right)^9
       =\frac{1162261467}{33554432},
 \label{eq:B-R3-bound}\\
 R_{4,16}&\le C_4^B s^2(2cS_4+H),
 &C_4^B&=272\left(\frac98\right)^{17}
       =\frac{283512088894331673}{140737488355328}.
 \label{eq:B-R4-bound}
\end{align}
\end{proposition}

\begin{proof}
\bi{For $r=9/10$, comparison of each relative eigenvalue gives}
{取 $r=9/10$，逐个比较相对本征值得}
\begin{equation}
 2\!\sum_{\substack{z\ge9\\z\ {\rm odd}}}
       2^{-z}|\xi_z\rangle\langle\xi_z|
 \le \frac2r\left(\frac1{2r}\right)^9r^{n_\xi+1}.
 \label{eq:B-tail-scalar}
\end{equation}
\bi{At $z=9$ the corresponding scalar coefficients agree; for higher odd $z$ their ratio decreases because $2r>1$, and the left-hand coefficient is zero in every other degree. Choosing $a=2/27$ gives $r=2/(2+3a)$. Apply the positive maps $W_{3,\rm rel}(\cdot)W_{3,\rm rel}^\dagger$, $\Omega_3$, and $\mathcal L_3$, in that order. Equation~\eqref{eq:B-three-thermal-bound} then proves \eqref{eq:B-tail-bound}, because $(2/(ar))(1/(2r))^9=30(5/9)^9$.}
{在 $z=9$ 时，相应的标量系数相等；在更高奇数次数上，由于 $2r>1$，其比值递减；其余次数的左侧系数为零。取 $a=2/27$ 即有 $r=2/(2+3a)$。依次施加正映射 $W_{3,\rm rel}(\cdot)W_{3,\rm rel}^\dagger$、$\Omega_3$ 和 $\mathcal L_3$。再由式~\eqref{eq:B-three-thermal-bound} 得式~\eqref{eq:B-tail-bound}，其中 $(2/(ar))(1/(2r))^9=30(5/9)^9$。}

\bi{For the finite three-body operator, take $a=1/12$ and $r=8/9$. On relative degrees $0\le z\le8$, $r^{z+1}\ge r^9$, so $\Pi_{\le8}^{\rm aux}\le r^{-9}r^{n_\xi+1}$. Since $W_3$ intertwines the center-of-mass oscillator and preserves total degree, \eqref{eq:B-insertion-norms} also gives}
{对有限三体算符，取 $a=1/12$、$r=8/9$。在相对次数 $0\le z\le8$ 上，$r^{z+1}\ge r^9$，因此 $\Pi_{\le8}^{\rm aux}\le r^{-9}r^{n_\xi+1}$。由于 $W_3$ 交织质心振子并保持总次数，式~\eqref{eq:B-insertion-norms} 还给出}
\begin{equation}
 W_{k,\rm rel}\Pi_{\le D}^{\rm aux}W_{k,\rm rel}^\dagger
   =H_{k,\rm rel}^{\rm pl}\Pi_{\le D}^{(k)},
 \qquad k=3,4.
 \label{eq:B-wide-cutoff}
\end{equation}
\bi{Applying the same positive maps and \eqref{eq:B-three-thermal-bound} gives \eqref{eq:B-R3-bound}, with $r^{-9}/a=12(9/8)^9$.}
{施加同样的正映射并使用式~\eqref{eq:B-three-thermal-bound}，即得式~\eqref{eq:B-R3-bound}，其常数为 $r^{-9}/a=12(9/8)^9$。}

\bi{For four particles, use the full auxiliary space $V\otimes\Sym^2U$ from \eqref{eq:B-insertion-maps}. On $n_1+n_2\le16$ with $n_1$ even, the smallest eigenvalue of $r_1^{n_1+1}r_2^{n_2+1}$ is $r_1r_2^{17}$: $r_2\le r_1<1$, so the minimum places all sixteen relative quanta in the second oscillator. Hence the auxiliary cutoff is bounded by $(r_1r_2^{17})^{-1}$ times the thermal coefficient. Apply \eqref{eq:B-wide-cutoff} and \eqref{eq:B-four-thermal-bound}, and set $a=1/16$. Then $r_1=16/17$, $r_2=8/9$, and}
{对四粒子，使用式~\eqref{eq:B-insertion-maps} 中的完整辅助空间 $V\otimes\Sym^2U$。在 $n_1+n_2\le16$ 且 $n_1$ 为偶数时，$r_1^{n_1+1}r_2^{n_2+1}$ 的最小本征值为 $r_1r_2^{17}$：因为 $r_2\le r_1<1$，最小值对应将全部十六个相对量子置于第二个振子。因此，辅助截断投影不大于热系数乘以 $(r_1r_2^{17})^{-1}$。使用式~\eqref{eq:B-wide-cutoff} 和式~\eqref{eq:B-four-thermal-bound}，并取 $a=1/16$。此时 $r_1=16/17$、$r_2=8/9$，且}
\begin{equation}
 \frac1{a^2r_1r_2^{17}}=272\left(\frac98\right)^{17},
 \label{eq:B-four-constant}
\end{equation}
\bi{which proves \eqref{eq:B-R4-bound}. All comparisons were established at the positive few-body coefficient level before lifting, so their constants do not depend on the number of particles.}
{从而证明式~\eqref{eq:B-R4-bound}。所有比较均在正少体系数层面先建立、随后才提升，因此其常数不依赖粒子数。}
\end{proof}

\subsection{\bi{The fixed 138-row certificate}{固定的 138 行证书}}
\label{sec:B-rows}

\bi{The finite input is the planar row certificate supplied with the bosonic candidate manuscript of Shen~\cite{Shen2026BosonicCandidate}. Only its fixed rational row data and the finite planar block inequalities are used here. No spherical gap theorem or finite-flux interval estimate is an assumption. Appendix~\ref{sec:B-certificate} gives the exact data decoding, occupation metric, relative-block construction, and rational positivity test.}
{有限输入取自 Shen 的玻色候选手稿~\cite{Shen2026BosonicCandidate} 所附平面行证书。本节只使用其固定有理行数据与有限平面分块不等式，不假设任何球面谱隙定理或有限磁通区间估计。附录~\ref{sec:B-certificate} 给出数据的精确解码、占据数度量、相对块构造以及有理数正性检验。}

\bi{For each row $\rho$, the channel $t_\rho$ belongs to $\{0,1,2,3\}$. The respective row counts are $31,37,36,34$. Define, using raw periodized jets,}
{每一行 $\rho$ 的通道 $t_\rho$ 取自 $\{0,1,2,3\}$，相应行数为 $31,37,36,34$。使用原始周期化 jets 定义}
\begin{align}
 F^\rho(R)
 &=\lambda_\rho B_{t_\rho}^B(R)
   +\sum_{\substack{p,j\ge0\\t_\rho\le p+j\le8}}
      \alpha_{pj}^\rho
      a_{p+j-t_\rho}(R)^\dagger a_j(R)B_p^B(R),
 \label{eq:B-rows}\\
 K&=\int_T\sum_\rho F^\rho(R)^\dagger F^\rho(R)\,
                      \dd\mu_2(R)\ge0.
 \label{eq:B-SOS}
\end{align}
\bi{The frozen integers divided by $10^{12}$ specify the rational numbers $y_0^\rho$ and $y_{pj}^\rho$. Their relation to the physical coefficients is}
{冻结整数除以 $10^{12}$ 后定义有理数 $y_0^\rho$ 和 $y_{pj}^\rho$，它们与物理系数的关系为}
\begin{equation}
 y_0^\rho=\lambda_\rho\kappa_{t_\rho},\qquad
 y_{pj}^\rho=\alpha_{pj}^\rho\kappa_p\sqrt{\frac{j!}{i!}},
 \quad i=p+j-t_\rho,\qquad
 \kappa_p^2=\frac{p!}{2^{p+1}},\quad\kappa_p>0.
 \label{eq:B-row-decoding}
\end{equation}
\bi{Thus all real algebraic row coefficients are uniquely fixed. On the plane, the canonical polynomial variables $x_j=\sqrt{j!}\,a_j^\dagger$ and $\partial_j=a_j/\sqrt{j!}$ satisfy $[\partial_i,x_j]=\delta_{ij}$. In these variables $B_p=\kappa_p\widehat B_p$, where $\widehat B_p=\sum_{a+b=p}\partial_a\partial_b$, and each row becomes the rational differential operator}
{因而所有实代数行系数都被唯一固定。在平面上，正则多项式变量 $x_j=\sqrt{j!}\,a_j^\dagger$、$\partial_j=a_j/\sqrt{j!}$ 满足 $[\partial_i,x_j]=\delta_{ij}$。在这些变量中，$B_p=\kappa_p\widehat B_p$，其中 $\widehat B_p=\sum_{a+b=p}\partial_a\partial_b$；每行于是成为有理微分算符}
\begin{equation}
 F^\rho=y_0^\rho\widehat B_{t_\rho}
      +\sum_{p,j}y_{pj}^\rho x_{p+j-t_\rho}\partial_j\widehat B_p.
 \label{eq:B-rational-row}
\end{equation}
\bi{Physical adjoints in these coordinates are taken with the occupation metric $\prod_j n_j!(j!)^{n_j}$, not with the Euclidean metric of the coefficient vector.}
{这些坐标中的物理伴随须相对于占据数度量 $\prod_j n_j!(j!)^{n_j}$ 计算，不能采用系数向量的欧氏度量。}

\bi{Suppress the row label temporarily and put $i=p+j-t$, $i'=p'+j'-t$. With canonical planar commutators, the normal-ordered contributions beyond $\lambda^2B_t^\dagger B_t$ are exactly}
{暂时省略行指标，并记 $i=p+j-t$、$i'=p'+j'-t$。使用平面正则对易关系时，除 $\lambda^2B_t^\dagger B_t$ 外的正规序贡献精确为}
\begin{align}
 \mathsf K_3^{\rm row}
 &=\lambda\sum_{p,j}\alpha_{pj}
   \left(B_t^\dagger a_i^\dagger a_jB_p+
             (B_t^\dagger a_i^\dagger a_jB_p)^\dagger\right)
 \notag\\[-2pt]
 &\quad+\sum_{p,j,p',j'}\alpha_{pj}\alpha_{p'j'}\delta_{ii'}
           B_p^\dagger a_j^\dagger a_{j'}B_{p'},
 \label{eq:B-row-three}\\
 \mathsf K_4^{\rm row}
 &=\sum_{p,j,p',j'}\alpha_{pj}\alpha_{p'j'}
       B_p^\dagger a_j^\dagger a_{i'}^\dagger a_i a_{j'}B_{p'}.
 \label{eq:B-row-four}
\end{align}
\bi{The first coefficient has total orbital degree at most eight; the second has total orbital degree at most sixteen. Let $\overline{\mathsf K}_{k,\rm rel}$ be their summed center-of-mass partial traces with the factor $2/k$ from \eqref{eq:B-torus-seed}. The exact finite certificate is}
{第一个系数的总轨道次数至多为八，第二个至多为十六。记 $\overline{\mathsf K}_{k,\rm rel}$ 为将这些系数求和后作质心偏迹，并乘以式~\eqref{eq:B-torus-seed} 中因子 $2/k$ 的结果。精确有限证书为}
\begin{align}
 S_{3,\rm rel}^{\rm pl}\Pi_{\le8}^{(3)}
   -\overline{\mathsf K}_{3,\rm rel}
   +\delta_B H_{3,\rm rel}^{\rm pl}\Pi_{\le8}^{(3)}&\ge0,
 \label{eq:B-certificate-three}\\
 \alpha_B S_{4,\rm rel}^{\rm pl}\Pi_{\le16}^{(4)}
   -\overline{\mathsf K}_{4,\rm rel}
   +\delta_B H_{4,\rm rel}^{\rm pl}\Pi_{\le16}^{(4)}&\ge0,
 \label{eq:B-certificate-four}\\
 \delta_B=10^{-6},\qquad
 \mu_B=\frac1{50},\qquad
 \alpha_B=1-c_9-\mu_B.&
 \notag
\end{align}
\bi{These are $9+17=26$ rational block tests, including the zero-dimensional degree-one blocks. Rational Schur elimination verifies positive semidefiniteness; in each block the exact kernel of $H_k^{\rm pl}$ is preserved, and the positive rank equals that of $H_k^{\rm pl}$. The equations define finite-dimensional statements with no limiting tolerance. Their construction and the block dimensions and ranks are given in Appendix~\ref{sec:B-certificate}.}
{它们对应 $9+17=26$ 个有理分块检验，包括次数一的零维块。有理 Schur 消元检验其半正定性；每块均保持 $H_k^{\rm pl}$ 的精确零核，且正秩等于 $H_k^{\rm pl}$ 的秩。这些都是不含极限容差的有限维命题，其构造与分块维数、秩列于附录~\ref{sec:B-certificate}。}

\bi{The two-body row cost is}
{行证书的两体代价为}
\begin{equation}
 \eta=\sum_\rho\lambda_\rho^2
 =\frac{764191703750823432978457}
        {1500000000000000000000000}.
 \label{eq:B-eta}
\end{equation}
\bi{Apply the positive periodization map and then the positive lift to \eqref{eq:B-certificate-three}--\eqref{eq:B-certificate-four}. Define}
{先对式~\eqref{eq:B-certificate-three}--\eqref{eq:B-certificate-four} 施加正周期化映射，再作正提升。定义}
\begin{equation}
 K_3^0=\mathcal L_3\Omega_3(\overline{\mathsf K}_{3,\rm rel}),\qquad
 K_4=\mathcal L_4\Omega_4(\overline{\mathsf K}_{4,\rm rel}),\qquad
 T_4=\mathcal L_4\Omega_4
      (S_{4,\rm rel}^{\rm pl}\Pi_{\le16}^{(4)}).
 \label{eq:B-transferred-rows}
\end{equation}
\bi{This transfer is a congruence of coefficient tensors, followed by integration. It therefore requires neither an inverse torus few-body Gram matrix nor a choice of basis for its varying kernel. The exact connected formula \eqref{eq:boson-threebody} has nonnegative even-degree tail and negative odd-degree tail $-T_9$. Consequently, the transferred inequalities imply}
{这一移植先对系数张量作合同变换，再积分。因此它既不需要环面少体 Gram 矩阵的逆，也不需要为其变化的零核选取基。精确连通公式~\eqref{eq:boson-threebody} 的偶数次数尾为非负，而奇数次数负尾为 $-T_9$。故移植后的不等式给出}
\begin{equation}
 S_3\ge K_3^0-\delta_B R_{3,8}-T_9,
 \qquad
 K_4\le\alpha_B T_4+\delta_B R_{4,16}.
 \label{eq:B-certificate-transferred}
\end{equation}

\subsection{\bi{The raw-jet commutator error in Loewner order}{原始 jet 对易误差的 Loewner 序控制}}
\label{sec:B-metric}

\bi{On the torus, $[a_i(R),a_{i'}(R)^\dagger]=g_{ii'}(R)$. In the normal ordering of a row square, the only change from \eqref{eq:B-row-three}--\eqref{eq:B-row-four} is the replacement of the single contraction $\delta_{ii'}$ by $g_{ii'}(R)$. The zero-contraction four-body coefficient is unchanged. The exact pair ladder identity integrates the two-body coefficient to $\eta H$, so}
{在环面上，$[a_i(R),a_{i'}(R)^\dagger]=g_{ii'}(R)$。对每个行平方作正规序时，相比式~\eqref{eq:B-row-three}--\eqref{eq:B-row-four}，唯一变化是单次收缩中的 $\delta_{ii'}$ 被 $g_{ii'}(R)$ 替换；零收缩的四体系数保持不变。精确配对梯级恒等式将两体系数积分为 $\eta H$，因此}
\begin{equation}
 K=\eta H+K_3^0+K_4+R_g.
 \label{eq:B-K-decomposition}
\end{equation}
\bi{Let}
{令}
\begin{align}
 e_8&=\sup_{R\in T}\|g_{0,\ldots,8}(R)-I\|,
 \notag\\
 C_{\rm rows}^B
 &=\sum_{\rho,p,j}(\alpha_{pj}^\rho)^2
   =\frac{31171576329140678412829557}
          {8750000000000000000000000},
 \qquad
 \kappa_B=6C_{\rm rows}^B e_8.
 \label{eq:B-row-error-constants}
\end{align}

\begin{lemma}[\bi{Two-sided commutator estimate}{双边对易误差估计}]
\label{lem:B-metric}
\begin{equation}
 -\kappa_B R_{3,8}\le R_g\le\kappa_B R_{3,8}.
 \label{eq:B-metric-bound}
\end{equation}
\end{lemma}

\begin{proof}
\bi{At a fixed center, write $X_{pj}=a_jB_p^B$ and, for each row, collect $Y_i^\rho=\sum_{p+j-t_\rho=i}\alpha_{pj}^\rho X_{pj}$. The local remainder is the Hermitian quadratic form}
{固定中心后，记 $X_{pj}=a_jB_p^B$，并对每行合并为 $Y_i^\rho=\sum_{p+j-t_\rho=i}\alpha_{pj}^\rho X_{pj}$。局部余项为 Hermitian 二次型}
\begin{equation}
 R_g(R)=\sum_{\rho,i,i'}
      (Y_i^\rho)^\dagger\bigl(g_{ii'}(R)-\delta_{ii'}\bigr)Y_{i'}^\rho.
 \label{eq:B-local-metric}
\end{equation}
\bi{The matrix inequality $-e_8I\le g-I\le e_8I$ is valid for all indices occurring here. Apply it to the vector of output states $Y_i^\rho\psi$ and then apply Cauchy--Schwarz to each finite coefficient sum. This gives the pointwise Loewner bounds}
{这里出现的所有指标均满足矩阵不等式 $-e_8I\le g-I\le e_8I$。将它用于输出态向量 $Y_i^\rho\psi$，再对各有限系数和使用 Cauchy--Schwarz，得到逐中心的 Loewner 序界}
\begin{equation}
 -\mathcal E(R)\le R_g(R)\le\mathcal E(R),\qquad
 \mathcal E(R)=e_8C_{\rm rows}^B
       \sum_{p+j\le8}(a_jB_p^B)^\dagger(a_jB_p^B).
 \label{eq:B-local-envelope}
\end{equation}
\bi{There is no estimate of $\|a_j\|$ in this step. To integrate the envelope, use the auxiliary orthonormal vectors $v_p\otimes e_j$. Its coefficient is $W_3P_{\mathrm{degree}\le8}^{\rm aux}W_3^\dagger$. A relative oscillator state of degree $z\le8$ has exactly $9-z$ center-of-mass descendants in this cutoff. Equation~\eqref{eq:B-torus-seed} therefore assigns it the partial-trace coefficient $(2/3)(9-z)\le6$. Intertwining with $W_3$, preserving positive order, and lifting now give}
{这一步没有使用 $\|a_j\|$ 的任何估计。为积分这一包络，采用辅助正交向量 $v_p\otimes e_j$，其系数即 $W_3P_{\mathrm{degree}\le8}^{\rm aux}W_3^\dagger$。次数为 $z\le8$ 的相对振子态，在该截断中恰有 $9-z$ 个质心后裔。因此，式~\eqref{eq:B-torus-seed} 赋予它偏迹系数 $(2/3)(9-z)\le6$。利用 $W_3$ 的交织性质、正序保持以及正提升，得}
\begin{equation}
 \int_T\mathcal E(R)\,\dd\mu_2(R)
      \le6e_8C_{\rm rows}^B R_{3,8}.
 \label{eq:B-envelope-integrated}
\end{equation}
\bi{Integrating both inequalities in \eqref{eq:B-local-envelope} proves \eqref{eq:B-metric-bound}. Thus the raw-jet metric error remains a three-body error with a particle-number-independent coefficient.}
{对式~\eqref{eq:B-local-envelope} 的两边分别积分，即证式~\eqref{eq:B-metric-bound}。因此，原始 jet 度量误差仍是一个系数不依赖粒子数的三体误差。}
\end{proof}

\subsection{\bi{Assembly of the bosonic inequality}{玻色不等式的汇总}}
\label{sec:B-assembly}

\bi{Use the exact decomposition $H^2=\mathcal L_2(h_T^2)+S_3+S_4$ from Proposition~\ref{prop:connected}. The pair almost-projector estimate gives $\mathcal L_2(h_T^2)\ge(1-\epsilon_2)H$. The relative frame bound of Lemma~\ref{lem:relative-frame}, with charges two and two and relative degrees up to sixteen, gives $T_4\le(1+\epsilon_4)S_4$. In terms of the raw pair descendants, that finite four-body coefficient is explicitly}
{使用命题~\ref{prop:connected} 的精确分解 $H^2=\mathcal L_2(h_T^2)+S_3+S_4$。配对近投影估计给出 $\mathcal L_2(h_T^2)\ge(1-\epsilon_2)H$。引理~\ref{lem:relative-frame} 中电荷二与电荷二、相对次数至多十六的相对框架界给出 $T_4\le(1+\epsilon_4)S_4$。用原始配对后裔表示，这一有限四体系数显式为}
\begin{align}
 C_d(R)&=\sum_{p=0}^d(-1)^{d-p}
       \sqrt{\binom dp}\,2^{-d/2}B_p^B(R)B_{d-p}^B(R),
 \notag\\
 T_4&=\int_T\sum_{d=0}^{16}C_d(R)^\dagger C_d(R)\,
                       \dd\mu_4(R)
       \le(1+\epsilon_4)S_4.
 \label{eq:B-four-frame}
\end{align}
\bi{Odd $d$ vanish because the pair annihilators commute. The estimate retains the error as a multiple of $S_4$, which is the positive term available in the bosonic argument.}
{由于配对湮灭算符彼此对易，奇数 $d$ 的项消失。此估计将误差保留为 $S_4$ 的倍数，正好对应玻色论证可用的正项。}

\bi{Substitute \eqref{eq:B-certificate-transferred} into the connected decomposition, then eliminate $K_3^0+K_4$ using \eqref{eq:B-K-decomposition}. Since $\alpha_B>0$, the four-body frame comparison is used in the correct direction. Finally use the upper bound on $R_g$ from Lemma~\ref{lem:B-metric}. The result is}
{将式~\eqref{eq:B-certificate-transferred} 代入连通分解，再用式~\eqref{eq:B-K-decomposition} 消去 $K_3^0+K_4$。由于 $\alpha_B>0$，四体框架比较具有所需方向。最后使用引理~\ref{lem:B-metric} 对 $R_g$ 的上界，得到}
\begin{align}
 H^2\ge{}&K+(1-\epsilon_2-\eta)H
       +[1-\alpha_B(1+\epsilon_4)]S_4
 \notag\\
 &-(\delta_B+\kappa_B)R_{3,8}
       -\delta_B R_{4,16}-T_9.
 \label{eq:B-master-before-payment}
\end{align}
\bi{Apply Proposition~\ref{prop:B-comparisons} to the three negative terms. This gives}
{对三个负项应用命题~\ref{prop:B-comparisons}，得}
\begin{equation}
 H^2\ge K+A_\Lambda H+B_\Lambda S_4,
 \label{eq:B-master}
\end{equation}
\begin{align}
 A_\Lambda
 &=1-\epsilon_2-\eta
       -s\bigl[c_9+(\delta_B+\kappa_B)C_3^B\bigr]
       -\delta_B C_4^B s^2,
 \label{eq:B-A-budget}\\
 B_\Lambda
 &=1-\alpha_B(1+\epsilon_4)
       -sc\bigl[c_9+(\delta_B+\kappa_B)C_3^B\bigr]
       -2\delta_B C_4^B s^2c.
 \label{eq:B-B-budget}
\end{align}
\bi{As the actual shortest nonzero lattice vector tends to infinity, the common frame estimates give $\epsilon_2,\epsilon_4,e_8\to0$ and $s,c\to1$, uniformly in twists and independently of the displayed lattice basis. Therefore,}
{当实际最短非零格矢趋于无穷时，公共框架估计给出 $\epsilon_2,\epsilon_4,e_8\to0$ 以及 $s,c\to1$，并且对边界 twists 一致，不依赖选取的格基。因此，}
\begin{align}
 A_\infty
 &=1-\eta-c_9-\delta_B(C_3^B+C_4^B)
 \notag\\
 &=\frac{23382014633024415475988526498535740455597}
         {69331597085638656000000000000000000000000}
   \approx0.33724904107059384>\frac13,
 \label{eq:B-A-infinity}\\
 B_\infty
 &=\mu_B-\delta_B(C_3^B+2C_4^B)
   =\frac{1121425355698906343}{70368744177664000000}
   \approx0.015936412803780887>0.
 \label{eq:B-B-infinity}
\end{align}
\bi{Since $K$ and $S_4$ are positive, every $\gamma<A_\infty$ admits a systole threshold for $H^2\ge\gamma H$. We next give a fully explicit sufficient threshold for $\gamma=1/3$.}
{由于 $K$ 和 $S_4$ 均为正，每个 $\gamma<A_\infty$ 都存在一个最短周期阈值，使 $H^2\ge\gamma H$ 成立。下面对 $\gamma=1/3$ 给出完全显式的充分阈值。}

\subsection{\bi{An explicit rational budget at systole twenty}{最短周期为二十时的显式有理预算}}
\label{sec:B-systole}

\bi{Assume $\ell(\Lambda)\ge L=20$ in the dimensionless coordinate $z=x/(\sqrt2\ell_B)$. Partition the nonzero lattice vectors into the shells $nL\le|\lambda|<(n+1)L$, $n\ge1$. Disks of radius $L/2$ about distinct lattice points have disjoint interiors. A shell is contained in the disk of radius $(n+3/2)L$ after adjoining these smaller disks, so its number of points is at most $(2n+3)^2$. This bound uses only the shortest lattice vector.}
{假设在无量纲坐标 $z=x/(\sqrt2\ell_B)$ 中，$\ell(\Lambda)\ge L=20$。将非零格矢按壳层 $nL\le|\lambda|<(n+1)L$（$n\ge1$）划分。以不同格点为中心、半径为 $L/2$ 的圆盘内部互不相交；将这些小圆盘附于一壳层后，整体包含在半径 $(n+3/2)L$ 的圆盘中，因此每壳格点数至多为 $(2n+3)^2$。这一估计只使用最短格矢。}

\bi{For the displacement estimate with $\nu=1$ and indices $d,e\le D$, the finite sum in the common jet bound has at most $D+1$ terms. For $r=|\lambda|\ge20$, each term is at most $D!r^{2D}e^{-r^2/2}$. A row-sum bound on the $(D+1)\times(D+1)$ matrix therefore gives}
{对 $\nu=1$、指标 $d,e\le D$ 的位移估计，公共 jet 界中的有限和至多含 $D+1$ 项。对 $r=|\lambda|\ge20$，每项不超过 $D!r^{2D}e^{-r^2/2}$。再对 $(D+1)\times(D+1)$ 矩阵使用行和界，得到}
\begin{equation}
 \mathrm{error}_D\le(D+1)^2D!
   \sum_{\lambda\ne0}|\lambda|^{2D}e^{-|\lambda|^2/2}.
 \label{eq:B-jet-shell-start}
\end{equation}
\bi{For $D\le16$, $r^{2D}e^{-r^2/2}$ decreases on $r\ge20$. The ratio of consecutive shell majorants is bounded by $2^{33}e^{-600}<1/2$: the shell-count ratio is less than two, the polynomial ratio is at most $2^{32}$, and the exponential ratio is at most $e^{-600}$. Thus the sum is bounded by twice its first-shell majorant, giving}
{当 $D\le16$ 时，$r^{2D}e^{-r^2/2}$ 在 $r\ge20$ 上递减。相邻壳层上界之比不超过 $2^{33}e^{-600}<1/2$：壳层计数之比小于二，多项式之比至多 $2^{32}$，指数之比至多 $e^{-600}$。因此该和不超过首壳上界的两倍，即}
\begin{equation}
 \sum_{\lambda\ne0}|\lambda|^{2D}e^{-|\lambda|^2/2}
       \le50\,20^{2D}e^{-200}.
 \label{eq:B-shell-bound}
\end{equation}
\bi{The elementary Taylor bound $e^{25/36}>\sum_{j=0}^{4}(25/36)^j/j!>2$ proves $\log2<25/36$. It follows that $e^{-200}<2^{-288}$, $e^{-100}<2^{-144}$, and $e^{-400}<2^{-576}$. Applying \eqref{eq:B-jet-shell-start} with $D=16$ to the two-center pair frame, and with $D=8$ to the raw-jet Gram matrix, yields}
{初等 Taylor 界 $e^{25/36}>\sum_{j=0}^{4}(25/36)^j/j!>2$ 证明 $\log2<25/36$，从而 $e^{-200}<2^{-288}$、$e^{-100}<2^{-144}$、$e^{-400}<2^{-576}$。将式~\eqref{eq:B-jet-shell-start} 的 $D=16$ 情形用于双中心配对框架、$D=8$ 情形用于原始 jet Gram 矩阵，得到}
\begin{align}
 \epsilon_4
 &\le\frac{50\cdot17^2\cdot16!\cdot20^{32}}{2^{288}}
      <2.612\times10^{-28},
 \label{eq:B-explicit-e4}\\
 e_8
 &\le\frac{50\cdot9^2\cdot8!\cdot20^{16}}{2^{288}}
      <2.152\times10^{-58}.
 \label{eq:B-explicit-e8}
\end{align}
\bi{The same shell argument applied directly to the zero-degree pair error and the diagonal charge-$j$ kernels gives}
{将同一壳层论证直接用于零次数配对误差与电荷 $j$ 核的对角元，得}
\begin{equation}
 \epsilon_2\le50\,2^{-144}<2.243\times10^{-42},\qquad
 |k_1(R)-1|\le50\,2^{-288},\qquad
 |k_2(R)-1|\le50\,2^{-576}.
 \label{eq:B-explicit-diagonals}
\end{equation}
\bi{In particular, $k_1$ is positive, as required in \eqref{eq:B-density-constants}. To keep the final calculation rational and simple, enlarge all five errors $\epsilon_2,\epsilon_4,e_8,|k_1-1|,|k_2-1|$ to the common allowance $e=10^{-8}$. Then}
{特别地，$k_1$ 为正，满足式~\eqref{eq:B-density-constants} 的要求。为使最终计算保持为简单的有理数计算，将五个误差 $\epsilon_2,\epsilon_4,e_8,|k_1-1|,|k_2-1|$ 全部放宽到共同上界 $e=10^{-8}$。此时}
\begin{equation}
 s\le s_e:=1+e,\qquad
 c\le c_e:=\frac{1+e}{(1-e)^2},\qquad
 \kappa_B\le\kappa_e:=6C_{\rm rows}^Be.
 \label{eq:B-rational-allowances}
\end{equation}
\bi{Every subtracted term in \eqref{eq:B-A-budget}--\eqref{eq:B-B-budget} is monotone in these nonnegative errors. Substitute $s_e,c_e,\kappa_e$ and $\epsilon_2=\epsilon_4=e$ to obtain rational lower bounds $\underline A$ and $\underline B$. Their exact values are recorded in Appendix~\ref{sec:B-certificate}; direct rational comparison gives}
{式~\eqref{eq:B-A-budget}--\eqref{eq:B-B-budget} 中的每个被减项都随这些非负误差单调增加。代入 $s_e,c_e,\kappa_e$ 和 $\epsilon_2=\epsilon_4=e$，得到有理下界 $\underline A$ 和 $\underline B$。其精确值列于附录~\ref{sec:B-certificate}；直接进行有理数比较得}
\begin{align}
 A_\Lambda\ge\underline A
 &\approx0.3372416256964786,\qquad \underline A>\frac{337241}{10^6}>\frac13,
 \label{eq:B-explicit-A}\\
 B_\Lambda\ge\underline B
 &\approx0.015928994442432054,\qquad \underline B>\frac{15928}{10^6}>0.
 \label{eq:B-explicit-B}
\end{align}
\bi{Equations~\eqref{eq:B-master}, \eqref{eq:B-explicit-A}, and \eqref{eq:B-explicit-B}, together with $K,S_4\ge0$, prove $H_B^2\ge H_B/3$ on every particle sector whenever $\ell(\Lambda)\ge20$. The threshold is sufficient and deliberately conservative; no optimization of the lattice budget is used.}
{由式~\eqref{eq:B-master}、\eqref{eq:B-explicit-A}、\eqref{eq:B-explicit-B} 以及 $K,S_4\ge0$，可知当 $\ell(\Lambda)\ge20$ 时，每个粒子数扇区均满足 $H_B^2\ge H_B/3$。这一阈值是刻意保守的充分条件，证明没有使用格点预算的优化。}

\bi{Finally, take the algebraic direct sum of finite-particle sectors as a common core. On each sector the preceding inequality excludes spectrum in $(0,1/3)$. The self-adjoint direct sum $H_B=\bigoplus_{N\ge0}H_{B,N}$ has the same spectral exclusion, and its spectral calculus gives the asserted full-Fock quadratic-form inequality. This step does not require $H_B$ to be a bounded operator on Fock space. At flux $M=2N$, the common zero-mode argument identifies its two-dimensional Laughlin kernel, so the bound is a gap above that kernel. For a square torus, $\ell(\Lambda)=\sqrt{\pi M}$, and on the even-flux sequence the sufficient condition holds for $N\ge64$.}
{最后，以有限粒子数扇区的代数直和为共同核心。上述不等式在每个扇区排除 $(0,1/3)$ 内的谱。自伴直和 $H_B=\bigoplus_{N\ge0}H_{B,N}$ 具有相同的谱排除性质，通过谱演算即得所述全 Fock 空间二次型不等式。这一步不要求 $H_B$ 在 Fock 空间上有界。在磁通 $M=2N$ 时，公共零模论证将零核识别为二维 Laughlin 空间，故该界给出零核上方的谱隙。对方形环面，$\ell(\Lambda)=\sqrt{\pi M}$，因此在偶磁通序列上，$N\ge64$ 即满足这一充分条件。}

\section{\bi{Laughlin zero modes and the geometric scope}{Laughlin零模与几何适用范围}}\label{sec:groundspaces}
\begin{proposition}[\bi{The zero-energy space at filling $1/m$}{填充$1/m$处的零能空间}]\label{prop:kernel}
\bi{Take $m=2$ for the bosonic $V_0$ interaction or $m=3$ for the fermionic $V_1$ interaction. At flux $M=mN$ and any fixed common single-particle twist, $\dim\ker H=m$. The kernel is exactly the torus Laughlin space.}
{对玻色$V_0$取$m=2$，对费米$V_1$取$m=3$。在磁通$M=mN$及任意固定共同单粒子twist下，有$\dim\ker H=m$，且零空间恰为环面Laughlin空间。}
\end{proposition}
\begin{proof}
\bi{We use the torus Laughlin--Jastrow theta construction of Haldane--Rezayi and Greiter--Schnells--Thomale~\cite{HaldaneRezayi1985Periodic,GreiterSchnellsThomale2016} for the factors and their multipliers. The following argument must additionally identify the kernel of the physical interaction and exhaust it, including the flat twist. This kernel argument is not included in the selected gap-root acceptance.}
{因子及其乘子使用Haldane--Rezayi和Greiter--Schnells--Thomale的环面Laughlin--Jastrow theta构造~\cite{HaldaneRezayi1985Periodic,GreiterSchnellsThomale2016}。下述论证还须识别并穷尽物理相互作用的核，包括平坦twist。这一核论证不在选定谱隙根的验收覆盖内。}
\bi{Positivity of~\eqref{eq:physical-H} makes zero energy equivalent to vanishing of the contact pair amplitude at every center, with the other particle coordinates arbitrary. In the bosonic case the wavefunction vanishes on every pair diagonal; exchange symmetry makes its local relative zero at least quadratic. In the fermionic case the wavefunction is already odd in the relative coordinate and the relative-degree-one contact derivative also vanishes, so its zero has order at least three. Analyticity promotes the integral zero condition to the pointwise statement.}
{式~\eqref{eq:physical-H}的正性使零能条件等价于在任意中心、其他粒子坐标任意时接触pair振幅均为零。玻色情形的波函数在每条pair对角线上消失；交换对称性使局部相对零点至少二阶。费米波函数本已对相对坐标为奇函数，而相对次数一的接触导数也消失，因此零点至少三阶。解析性把积分为零的条件提升为逐点条件。}

\bi{Put $E=\mathbb C/\Lambda$, choose $y_i=z_i/\omega_1$ and $\tau=\omega_2/\omega_1$, and divide the holomorphic factor of the wavefunction by}
{令$E=\mathbb C/\Lambda$，选$y_i=z_i/\omega_1$及$\tau=\omega_2/\omega_1$，把波函数的全纯部分除以}
\begin{equation}
 J_m(\boldsymbol y)=\prod_{i<j}\vartheta_1(y_i-y_j\mid\tau)^m.
\end{equation}
\bi{The quotient is holomorphic, including at intersections of pair diagonals, because its divisor has no negative component. The original line bundle on $E^N$ has Chern matrix $M I$. The Jastrow factor has matrix $m(NI-\mathbf1\mathbf1^{\mathsf T})$. At $M=mN$ the quotient therefore has matrix}
{商是全纯的，包括pair对角线的交点，因为其除子没有负分量。$E^N$上原线丛的Chern矩阵为$MI$，Jastrow因子的矩阵为$m(NI-\mathbf1\mathbf1^{\mathsf T})$。在$M=mN$时，商的矩阵因此为}
\begin{equation}
 mN I-m(NI-\mathbf1\mathbf1^{\mathsf T})=m\mathbf1\mathbf1^{\mathsf T}.
\end{equation}
\bi{The flat part also descends through the sum map, not just the Chern class. In a theta-function gauge the original holomorphic factor obeys $f(y_i+1)=e^{i\phi_1}f$ and $f(y_i+\tau)=e^{i\phi_2}e^{-i\pi M(2y_i+\tau)}f$. The theta multipliers show that, after division, a shift of any one coordinate has multiplier}
{不仅Chern类如此，平坦部分也通过求和映射下降。在theta函数规范中，原全纯部分满足$f(y_i+1)=e^{i\phi_1}f$及$f(y_i+\tau)=e^{i\phi_2}e^{-i\pi M(2y_i+\tau)}f$。Theta乘子表明，除掉Jastrow因子后，平移任意一个坐标的乘子均为}
\begin{equation}
 (-1)^{m(N-1)}e^{i\phi_1},\qquad
 (-1)^{m(N-1)}e^{i\phi_2}e^{-i\pi m(2Z+\tau)},\qquad Z=\sum_i y_i.
\end{equation}
\bi{Hence the quotient bundle is exactly $\Sigma^*L$, where $\Sigma:E^N\to E$ is summation and $L$ is a degree-$m$ line bundle with the indicated fixed flat twist. The fibers are connected compact copies of $E^{N-1}$. A holomorphic section of a pullback is constant along each fiber in a local trivialization, so}
{因此商线丛恰为$\Sigma^*L$，其中$\Sigma:E^N\to E$是求和映射，$L$是具有上述固定平坦twist的次数$m$线丛。其纤维是连通紧的$E^{N-1}$。拉回线丛的全纯截面在局部平凡化后沿每条纤维为常数，故}
\begin{equation}
 H^0(E^N,\Sigma^*L)\cong H^0(E,L),\qquad \dim H^0(E,L)=m.
\end{equation}
\bi{The last equality is Riemann--Roch on the elliptic curve for positive degree. Conversely, every such COM theta section multiplied by $J_m$ has the required statistics, boundary conditions, and pair zeros, so it is a zero mode. This proves both existence and exhaustion. The general torus theta construction is consistent with~\cite{HaldaneRezayi1985Periodic,GreiterSchnellsThomale2016}; the rectangular-torus pairing formulas and the fermionic threefold exhaustion statement in~\cite{OrtizNussinovDukelskySeidel2013} provide complementary background, not the proof for arbitrary shape and twist used here.}
{最后一个等式来自椭圆曲线上正次数线丛的Riemann--Roch定理。反过来，每个这样的COM theta截面乘以$J_m$都具有所需统计、边界条件及pair零点，因而是零模。这同时证明了存在性与穷尽性。一般环面theta构造与文献~\cite{HaldaneRezayi1985Periodic,GreiterSchnellsThomale2016}一致；文献~\cite{OrtizNussinovDukelskySeidel2013}中的矩形环面pairing公式与费米三重零模穷尽声明提供互补背景，但不替代这里对任意形状和twist的证明。}
\end{proof}

\bi{For a nonnegative self-adjoint operator, $H^2\ge\gamma H$ is equivalent by spectral calculus to $\operatorname{spec}(H)\subset\{0\}\cup[\gamma,\infty)$, or $H\ge\gamma(I-P_0)$ with $P_0$ the full kernel projection. This needs no uniqueness of the ground state. Applying Proposition~\ref{prop:kernel} proves Corollary~\ref{cor:fillings}. The torus relations are $M=2N$ and $M=3N$, not the spherical shifted relations $2(N-1)$ and $3(N-1)$.}
{对非负自伴算符，谱演算使$H^2\ge\gamma H$等价于$\operatorname{spec}(H)\subset\{0\}\cup[\gamma,\infty)$，亦即$H\ge\gamma(I-P_0)$，其中$P_0$是整个零空间的投影。这不要求基态唯一。应用命题~\ref{prop:kernel}即得推论~\ref{cor:fillings}。环面关系是$M=2N$及$M=3N$，不是球面的带shift关系$2(N-1)$和$3(N-1)$。}

\bi{At fixed shape the dimensionless systole grows as $\sqrt M$, uniformly on compact shape families. More general families are covered whenever their actual systole tends to infinity, even if their aspect ratio is not bounded. Conversely, a long displayed basis vector does not compensate for a short nonzero lattice combination. The proofs do not address fixed short circumferences. No claim uniform over arbitrary degenerating torus shapes follows from these estimates.}
{固定形状时，无量纲最短周期按$\sqrt M$增长，并在紧的形状族上一致成立。更一般的族只要实际最短周期趋于无穷，也在适用范围内，即使长宽比没有统一上界。相反，展示出的长格基向量无法补偿某个很短的非零格点组合。本证明不处理固定短周长。这些估计不能推出对任意退化环面形状都一致成立的结论。}

\section{\bi{Discussion}{讨论}}\label{sec:discussion}
\bi{The two proofs identify a common geometric mechanism with different statistical costs. Magnetic periodization preserves finite positive relative operators exactly; it does not preserve products of arbitrary plane operators. The missing product information is recovered through the full-COM pair frame and the explicit connected-kernel unfolding. This separation explains why local few-body sum-of-squares data can lead to particle-number-uniform torus estimates without a semidefinite program whose dimension grows with flux.}
{两条证明体现了共同几何机制与不同统计代价。磁周期化精确保留有限的正相对算符，但不保持任意平面算符的乘积。乘积中缺少的信息由全COM pair frame及明确的连通核展开补回。这种分离解释了为何局部少体平方和数据能够导出对粒子数一致的环面估计，而不需要维数随磁通增长的半正定规划。}

\bi{The bosonic comparison is not obtained by changing signs in the fermionic proof. Its coherent thermal operators first dominate relative cutoffs and the negative tail as positive few-body coefficients. Their periodizations become integrals over a torus anchor and unrestricted planar spectator offsets. Occupation inequalities then control those integrals by $H$ and $S_4$. Retaining a positive $S_4$ reserve is what permits arbitrarily large occupations. In contrast, Pauli bounds allow the fermionic finite left factors and tail to be paid for by $H$ alone.}
{玻色比较不是在费米证明中改几个符号得到的。其相干thermal算符先在正少体系数层支配相对截断与负尾，周期化后成为环面anchor与不受限平面spectator位移上的积分，再由占据数不等式用$H$及$S_4$控制。保留正$S_4$余量使任意大占据数成为可处理的。相反，Pauli界允许费米有限左因子与尾项都只由$H$支付。}

\bi{The constants are sufficient lower bounds in a fixed physical normalization, not estimates of the actual excitation energies. The bosonic threshold $\ell(\Lambda)\ge20$ is explicit but intentionally coarse. The fermionic threshold remains existential because the high-degree negative-tail control is taken by dominated convergence. A quantitative treatment of that tail would turn this part into an explicit finite-size statement. Improved planar certificates, including the stronger fermionic bounds in~\cite{Gao2026ImprovedGap}, are possible inputs to a further transfer, but their finite matrices and geometric budgets would have to be checked afresh; no such improved torus constant is asserted here.}
{这些常数是在固定物理归一化下的充分下界，不是实际激发能的估计。玻色阈值$\ell(\Lambda)\ge20$显式但有意保守。费米阈值仍是存在性的，因为高次数负尾通过支配收敛控制。若对该尾项作定量处理，就可把这一部分变为显式有限尺寸结论。更好的平面证书，包括文献~\cite{Gao2026ImprovedGap}中更强的费米下界，可能成为进一步迁移的输入，但必须重新检查相应有限矩阵与几何预算；本文不声称这些改进的环面常数已经成立。}

\bi{There are two scope boundaries. First, we do not replace the systole hypothesis by a claim of arbitrary shape uniformity, as discussed in Section~\ref{sec:groundspaces}. Second, this transfer is a quotient construction, not an orbital compression; Appendix~\ref{app:disk} gives an exact disk counterexample and identifies the failing frame estimate. The results also do not prove stability under general perturbations of the torus Hamiltonian. Establishing such stability would require further local-energy and perturbation estimates beyond the finite certificate comparison developed here.}
{适用范围有两条边界。第一，本文不将最短周期条件替换为任意形状的一致性声明，见第~\ref{sec:groundspaces}节。第二，本迁移是商空间构造，不是轨道压缩；附录~\ref{app:disk}给出精确disk反例，并指出失效的frame估计。本文也没有证明环面Hamiltonian在一般扰动下的稳定性。要得到这种稳定性，需要超出这里有限证书比较之外的局部能量与扰动估计。}

\section{\bi{Conclusion}{结论}}
\bi{The bounds below are the mathematical conclusion of the candidate argument, with the formal-verification scope specified in Appendix~\ref{app:reproducibility}. No priority claim is made.}
{下述界是本候选论证的数学结论，形式化验证范围见附录~\ref{app:reproducibility}。本文不作首创性声明。}
\bi{Fixed finite Fock-space certificates can be transferred to full bosonic $V_0$ and fermionic $V_1$ interactions on large-systole flat tori by combining positive periodization, parameter COM integration, and all-COM jet-frame bounds. The separate CCR and CAR comparisons give the gap estimates of Theorems~\ref{thm:boson} and~\ref{thm:fermion}; at $M=2N$ and $M=3N$ they apply above the entire two- and three-dimensional Laughlin ground spaces. The proof depends on large actual cycles and on the specified contact normalization, and does not extend by compression to a hard-cutoff disk.}
{通过结合正周期化、参数COM积分与全COM jet-frame界，可以把固定有限Fock空间证书迁移到大最短周期平坦环面上的完整玻色$V_0$和费米$V_1$相互作用。分别采用CCR和CAR比较，得到定理~\ref{thm:boson}与~\ref{thm:fermion}的谱隙估计；在$M=2N$及$M=3N$时，它们控制整个二维、三维Laughlin基态空间之上的谱。证明依赖实际周期足够大及明确的接触归一化，不能通过压缩延拓到硬截断disk。}

\appendix
\section{\bi{The complete finite fermionic certificate}{完整的有限费米证书}}
\label{app:F-certificate}

\bi{This appendix specifies all finite input data and all arithmetic required for \eqref{eq:F-three-certificate} and \eqref{eq:F-four-certificate}. The coefficients originate in \cite{OpenAI2026Fock}; the formulation below uses their planar oscillator values directly. The proof consists of finite sums of integers and rational numbers, followed by exact positivity tests. It neither estimates a many-particle spectrum nor introduces a size-dependent optimization problem.}{本附录给出式~\eqref{eq:F-three-certificate}和~\eqref{eq:F-four-certificate}所需的全部有限输入数据及算术。系数源自文献~\cite{OpenAI2026Fock}；下文直接使用其平面振子数值。证明由有限整数、有理数求和以及精确正性检验组成，既不估算多粒子能谱，也不引入随系统尺寸增长的优化问题。}

\subsection{\bi{All seven integer rows}{全部七行整数数据}}
\label{app:F-rows}

\bi{Let $P=10^7$. The first two integers of each row are $t_\rho$ and $\ell_\rho=P\lambda_\rho$. An entry written as $pj:a$ means $a_{pj}^{\rho}=a$, with $p$ and $j$ read as separate single-digit indices. Indented continuation lines belong to the preceding row. Every omitted entry is zero. Thus the following data, together with \eqref{eq:F-row-decode}, determine every coefficient without a fitted or rounded decimal input.}{令 $P=10^7$。每行的前两个整数为 $t_\rho$ 及 $\ell_\rho=P\lambda_\rho$。条目 $pj:a$ 表示 $a_{pj}^{\rho}=a$，其中 $p,j$ 应解析为两个单独的个位指标。缩进续行属于前一行，所有未列出的条目均为零。因此下列数据与式~\eqref{eq:F-row-decode}共同确定全部系数，无需拟合值或舍入的小数输入。}

\begingroup\scriptsize
\begin{verbatim}
0 7181518 05:-849651 07:-2958744 15:751659 16:-43084 17:-146216 32:720949 35:-127514
  42:222437 43:-262822 53:-117700 61:-55024 71:-166049
1 -112991 05:142772 06:-142342 07:527055 08:50742 16:197917 18:-18343 24:5764 27:-21183
  32:75832 35:-6115 45:2660 53:-9455 61:-27091 63:-3494 71:479 72:-15331
1 -4247311 05:-12207250 06:-14064191 07:-14063663 08:-15662679 16:483831 18:3987214
  24:-814558 27:888887 32:966462 35:145978 45:-171722 53:34350 61:389319 63:97603
  71:412447 72:-53786
2 3438746 14:6689266 15:6620205 16:6625065 17:4969272 18:4878238 27:-1387380 36:-586188
  37:344087 42:-1120634 45:-318530 46:559169 55:157455 61:1207103 63:211617 64:33240
  71:577001 73:-142578
5 -98645 05:-2719776 08:-8519913 14:-949523 17:-1952105 18:-1562606 28:-8919966 32:248826
  35:-1017623 37:3656005 38:-2081871 42:-33163 43:-305905 46:35677 47:-1320749
  48:-1040378 53:502094 55:1060633 58:-998003 61:-113617 63:184919 64:209239 65:-68198
  67:-266297 72:-606696 73:-563620 74:437345 76:-432189
5 -3160276 05:-1698129 08:10575875 14:-252576 17:-8002616 18:22592844 28:7370640
  32:289294 35:1400647 37:1987740 38:4030699 42:-1731401 43:-1034062 46:-1106371
  47:-1257051 48:1461735 53:1408149 55:1136847 58:1878501 61:-2838516 63:-734099
  64:1062416 65:11745 67:611312 72:-480776 73:20459 74:316925 76:37092
7 1441784 07:3441286 08:-6097155 37:-1002399 47:-251321 58:1143225 66:-17397 71:-1007464
  73:-662286 76:-344621 78:-1115284
\end{verbatim}
\endgroup

\bi{The row support satisfies $0\le p\le7$, $0\le j\le8$, and $0\le p+j-t_\rho\le8$. Direct integer addition gives}{各行支撑满足 $0\le p\le7$、$0\le j\le8$ 及 $0\le p+j-t_\rho\le8$。直接作整数加法得到}
\begin{equation}
 \sum_{\rho=1}^{7}\ell_\rho^2=93527408868499,
 \qquad \eta=\frac{93527408868499}{10^{14}}.
 \label{eq:F-eta-arithmetic}
\end{equation}

\subsection{\bi{Exact three-body arithmetic}{精确三体算术}}
\label{app:F-three-arithmetic}

\bi{For $c\in\{1,2\}$ and integers $0\le z,p\le T$, define the integer polynomial coefficient}{对 $c\in\{1,2\}$ 及整数 $0\le z,p\le T$，定义整数多项式系数}
\begin{equation}
 \mathsf U(c,z,T,p)=
 \sum_{h=\max(0,p+z-T)}^{\min(p,z)}
 (-1)^{z-h}\binom zh\binom{T-z}{p-h}c^{p-h}.
 \label{eq:F-U}
\end{equation}
\bi{Set $\mathsf U$ to zero whenever $z$ or $p$ lies outside $[0,T]$. Expanding \eqref{eq:F-coupling} at $u=1/\sqrt3$ gives}{当 $z$ 或 $p$ 超出 $[0,T]$ 时，将 $\mathsf U$ 定义为零。将式~\eqref{eq:F-coupling}在 $u=1/\sqrt3$ 时展开，得到}
\begin{equation}
 s^{(1/\sqrt3)}_{z,T,p}
 =\mathsf U(2,z,T,p)
       \sqrt{\frac{2^z\binom Tz}{3^T2^p\binom Tp}}.
 \label{eq:F-s-rational}
\end{equation}
\bi{Inserting this expression and \eqref{eq:F-row-decode} into \eqref{eq:F-ez} cancels every square root. Explicitly, let}{将此式及~\eqref{eq:F-row-decode}代入~\eqref{eq:F-ez}，全部平方根消去。具体地，令}
\begin{equation}
 X_{\rho,T}(z)=\sum_{p+j=T}a_{pj}^{\rho}\mathsf U(2,z,T,p).
 \label{eq:F-X-rational}
\end{equation}
\bi{Empty sums are zero, and the required rational expression is}{空求和取零，所需有理表达式为}
\begin{equation}
 P^2e_z=
 \sum_{\rho=1}^{7}\sum_{T=\max(z,t_\rho)}^{15}
 \frac{2^z\binom Tz}{3^T2^{t_\rho}\binom T{t_\rho}}
 X_{\rho,T}(z)
 \bigl[X_{\rho,T}(z)+2\ell_\rho\mathsf U(2,z,T,t_\rho)\bigr].
 \label{eq:F-e-rational}
\end{equation}
\bi{To see the cancellation directly, for a fixed $T$ the product $\alpha_{pj}^{\rho}s_{z,T,p}$ equals $a_{pj}^{\rho}/P$ times a factor independent of $p$. Its square is $2^z\binom Tz/(3^T2^{t_\rho}\binom T{t_\rho})$. The same factor multiplies $\lambda_\rho s_{z,T,t_\rho}$, which proves \eqref{eq:F-e-rational}.}{为直接看出根式的消去，固定 $T$ 后，乘积 $\alpha_{pj}^{\rho}s_{z,T,p}$ 等于 $a_{pj}^{\rho}/P$ 乘以一个与 $p$ 无关的因子，其平方为 $2^z\binom Tz/(3^T2^{t_\rho}\binom T{t_\rho})$。同一个因子也出现在 $\lambda_\rho s_{z,T,t_\rho}$ 中，故得到式~\eqref{eq:F-e-rational}。}

\bi{For each nonzero wedge sector, compute the exact margin}{对每个非零外积扇区，计算精确裕量}
\begin{equation}
 m_z=\frac32(3z-1)(-1/2)^z-e_z.
 \label{eq:F-margin}
\end{equation}
\bi{The following table gives the integer floors of $10^6m_z$. These are obtained by integer division after writing each rational number with a positive denominator.}{下表列出 $10^6m_z$ 的整数下取整值；将每个有理数写成分母为正的形式后，通过整数除法即可得到。}
\begin{center}
\small
\begin{tabular}{c|rrrrrrr}
$z$&2&4&5&6&7&8&9\\ \hline
$\lfloor10^6m_z\rfloor$&1298313&211317&249&250&249&250&250
\end{tabular}
\par\smallskip
\begin{tabular}{c|rrrrrr}
$z$&10&11&12&13&14&15\\ \hline
$\lfloor10^6m_z\rfloor$&250&249&250&250&6470&250
\end{tabular}
\end{center}
\bi{All thirteen floors are positive, and therefore every $m_z$ is strictly positive. This proves \eqref{eq:F-three-certificate} without floating-point sign decisions. The sectors $z=0,1,3$ are absent because their wedge maps vanish identically.}{十三个下取整值全部为正，因此每个 $m_z$ 均严格为正。这证明式~\eqref{eq:F-three-certificate}，且无需通过浮点数判断符号。$z=0,1,3$ 三个扇区未列入，因为其外积映射恒为零。}

\bi{For completeness, the constant controlling the torus Gram error is equally rational. Put $A_{\rho,i}=\sum_{p+j-t_\rho=i}(\alpha_{pj}^{\rho})^2$. Squaring \eqref{eq:F-row-decode} gives}{为完整起见，控制环面 Gram 误差的常数也同样是有理数。令 $A_{\rho,i}=\sum_{p+j-t_\rho=i}(\alpha_{pj}^{\rho})^2$，对式~\eqref{eq:F-row-decode}取平方得到}
\begin{equation}
 A_{\rho,i}=\frac1{P^2}
    \sum_{p+j-t_\rho=i}(a_{pj}^{\rho})^2
                  \frac{2^{p-t_\rho}t_\rho!i!}{p!j!}.
 \label{eq:F-A-rational}
\end{equation}
\bi{The nonzero-entry counts $n_\rho$, maximizing indices $i_\rho$, and maxima are}{非零项个数 $n_\rho$、取最大值的指标 $i_\rho$ 及相应最大值为}
\begin{center}
\small
\begin{tabular}{c|r|r|l}
$\rho$&$n_\rho$&$i_\rho$&$A_{\rho,i_\rho}$\\ \hline
1&12&8&$237052991863/390625000000$\\
2&16&6&$1154980099359/1400000000000000$\\
3&16&8&$4366100262613/15000000000000$\\
4&17&5&$90518611137713/600000000000000$\\
5&27&5&$289563620848499/5600000000000000$\\
6&27&2&$9393778357061/150000000000000$\\
7&10&8&$77741150041/6250000000000$
\end{tabular}
\end{center}
\bi{Comparison with the other eight rational values in each row verifies the maxima. Multiplying by $n_\rho$ and summing gives}{将每个最大值与该行其余八个有理值比较即可核实其最大性。乘以 $n_\rho$ 后求和，得到}
\begin{equation}
 C_\alpha^{\mathrm F}=\sum_{\rho=1}^{7}n_\rho A_{\rho,i_\rho}
  =\frac{297842353828763087}{16800000000000000}.
\end{equation}

\subsection{\bi{Rational reconstruction of all four-body blocks}{全部四体分块的有理重建}}
\label{app:F-four-arithmetic}

\bi{Fix $1\le D\le23$ and let all copy indices range over $\mathcal R_D$. Define}{固定 $1\le D\le23$，所有副本指标均在 $\mathcal R_D$ 中取值。定义}
\begin{equation}
 u_r=\frac1{r!(D-r)!},\qquad
 (E_D)_{rs}=\sqrt{u_ru_s}\,Y_{rs},\qquad
 (G_D)_{rs}=\sqrt{u_ru_s}\,Z_{rs}.
 \label{eq:F-scale}
\end{equation}
\bi{We now express $Y$ and $Z$ entirely over the rationals. First put}{下面将 $Y,Z$ 完全写成有理数矩阵。首先令}
\begin{equation}
 \mathsf V_r(D,T;p,j,k)
 =\mathsf U(1,r,j+k,j)\mathsf U(1,D-r,T-r,p),
 \label{eq:F-V}
\end{equation}
\bi{with value zero for $r>j+k$. Two expansions of \eqref{eq:F-coupling} at $u=1/\sqrt2$ give}{当 $r>j+k$ 时取零。对式~\eqref{eq:F-coupling}在 $u=1/\sqrt2$ 时展开两次，得到}
\begin{equation}
 S_r(D,T;p,j,k)=\mathsf V_r(D,T;p,j,k)
       \sqrt{2^{1-(j+k)-T+r}
                  \frac{j!k!p!\,u_r}{(T-D)!}}.
 \label{eq:F-S-rational}
\end{equation}
\bi{Consequently, the matrix defined by \eqref{eq:F-E} is determined by}{因此，式~\eqref{eq:F-E}定义的矩阵可由下式确定：}
\begin{align}
 P^2Y_{rs}=-\sum_{\rho=1}^{7}
  \sum_{\substack{(p,j),(q,l)\in\operatorname{supp}(a^\rho)\\T\ge D}}
 &a_{pj}^{\rho}a_{ql}^{\rho}
   \mathsf V_r(D,T;p,j,k)\mathsf V_s(D,T;q,l,i)\notag\\
 &\times2^{1-2T+p+q-t_\rho+(r+s)/2}
                    \frac{i!k!t_\rho!}{(T-D)!},
 \label{eq:F-Y-rational}
\end{align}
\bi{where $i=p+j-t_\rho$, $k=q+l-t_\rho$, and $T=p+j+k$. Both sums over input entries are ordered. To verify the square-root cancellation, multiply the two factors in \eqref{eq:F-S-rational} by the two coefficient factors in \eqref{eq:F-row-decode}. The factorial product under the combined radical is $(t_\rho!i!k!)^2u_ru_s/((T-D)!)^2$. Since $j+k=T-p$ and $l+i=T-q$, the exponent of two under that radical is $2+2p+2q-2t_\rho-4T+r+s$. Taking its positive square root and removing $\sqrt{u_ru_s}$ gives exactly \eqref{eq:F-Y-rational}.}{其中 $i=p+j-t_\rho$、$k=q+l-t_\rho$，以及 $T=p+j+k$。两个输入项的求和均为有序求和。为核查平方根消去，将式~\eqref{eq:F-S-rational}中的两个因子乘以式~\eqref{eq:F-row-decode}中的两个系数因子；合并根号内的阶乘乘积为 $(t_\rho!i!k!)^2u_ru_s/((T-D)!)^2$。由于 $j+k=T-p$、$l+i=T-q$，根号内二的指数为 $2+2p+2q-2t_\rho-4T+r+s$。取正平方根并除去 $\sqrt{u_ru_s}$ 后，恰得到式~\eqref{eq:F-Y-rational}。}

\bi{To reconstruct the Gram matrix, let $A$ be an increasing quadruple of nonnegative indices with sum $D+1$, and put $f(A)=\prod_{b\in A}b!$. For a choice $x<y$ from $A$, write the complementary indices as $j<k$ and set $p=x+y-1$. Let $\sigma(x,y,j,k)$ be the sign of sorting this ordered list into increasing order. Define}{为重建 Gram 矩阵，令 $A$ 为总和等于 $D+1$ 的非负整数递增四元组，并设 $f(A)=\prod_{b\in A}b!$。从 $A$ 中选取 $x<y$，将余下指标记为 $j<k$，并令 $p=x+y-1$。以 $\sigma(x,y,j,k)$ 表示将该有序列表重排为递增次序的符号。定义}
\begin{equation}
 L_r(A)=2^{(1-2D+r)/2}
 \sum_{\substack{x,y\in A,\ x<y\\\{j,k\}=A\setminus\{x,y\},\ j<k}}
 \sigma(x,y,j,k)(x-y)p!j!k!
                 \mathsf V_r(D,D;p,j,k).
 \label{eq:F-L-rational}
\end{equation}
\bi{Indeed, the coefficient of $e_x\wedge e_y$ in $v_p$ is $(x-y)\sqrt{p!/(2^px!y!)}$ when $x+y=p+1$. Multiplication by \eqref{eq:F-S-rational}, followed by the wedge sorting sign, shows that the coefficient of $w_{D,r}$ on the normalized Slater wedge $A$ is $\sqrt{u_r/f(A)}L_r(A)$. Taking inner products in this orthonormal wedge basis therefore gives}{事实上，当 $x+y=p+1$ 时，$v_p$ 在 $e_x\wedge e_y$ 上的系数为 $(x-y)\sqrt{p!/(2^px!y!)}$。将其乘以式~\eqref{eq:F-S-rational}并计入外积排序符号，可知 $w_{D,r}$ 在归一化 Slater 外积 $A$ 上的系数为 $\sqrt{u_r/f(A)}L_r(A)$。因此在此正交归一外积基中取内积，得到}
\begin{equation}
 Z_{rs}=\sum_{\substack{A\subset\mathbb Z_{\ge0},\ |A|=4\\\sum_{b\in A}b=D+1}}
                \frac{L_r(A)L_s(A)}{f(A)}.
 \label{eq:F-Z-rational}
\end{equation}
\bi{All powers of two in \eqref{eq:F-Y-rational} and \eqref{eq:F-L-rational} have integer exponents because $r,s$ are odd. Thus $Y$ and $Z$ are rational matrices. Empty quadruple sums, which occur at small $D$, are exactly zero.}{由于 $r,s$ 均为奇数，式~\eqref{eq:F-Y-rational}和~\eqref{eq:F-L-rational}中二的指数全部为整数，故 $Y,Z$ 均为有理矩阵。较小 $D$ 时出现的空四元组求和严格为零。}

\bi{Write $D_u=\operatorname{diag}(u_r)$ and $D_{\sqrt u}=\operatorname{diag}(\sqrt{u_r})$. Introduce rational matrices}{记 $D_u=\operatorname{diag}(u_r)$ 和 $D_{\sqrt u}=\operatorname{diag}(\sqrt{u_r})$，并引入有理矩阵}
\begin{equation}
 \mathsf B=\operatorname{diag}\bigl((2\mathbf1_{r=1}+\varepsilon_{\mathrm F})u_r\bigr)
                  -D_uYD_u,
 \qquad \mathsf M=Z\mathsf BZ.
 \label{eq:F-M-rational}
\end{equation}
\bi{Here $\mathsf M$ is a certificate matrix, distinct from the magnetic flux $M$. Substitution into \eqref{eq:F-scale} gives the exact congruence}{此处 $\mathsf M$ 为证书矩阵，与磁通量 $M$ 区分。代入式~\eqref{eq:F-scale}，得到精确合同关系}
\begin{equation}
 G_D(2\Pi_1-E_D+\varepsilon_{\mathrm F}I)G_D
      =D_{\sqrt u}\mathsf M D_{\sqrt u}.
 \label{eq:F-rational-congruence}
\end{equation}
\bi{Every $u_r$ is strictly positive. Hence proving $\mathsf M\ge0$ proves the desired certificate. This transformation only uses a positive diagonal congruence; it does not invert $G_D$ or discard its kernel.}{每个 $u_r$ 均严格为正，因此只需证明 $\mathsf M\ge0$ 即可得到所需证书。此变换仅使用正对角合同变换，既不求 $G_D$ 的逆，也不丢弃其核。}

\subsection{\bi{Exact positive-semidefiniteness tests}{精确半正定检验}}
\label{app:F-PSD}

\bi{The finite matrices in \eqref{eq:F-M-rational} can be tested by symmetric Schur elimination using rational numbers throughout. Order the copy indices increasingly and initialize $A=\mathsf M$. At each remaining index, read its diagonal entry $a$. If $a<0$, the matrix is not positive semidefinite. If $a=0$, every other entry in that row must be exactly zero; otherwise the associated $2\times2$ principal minor has determinant $-|b|^2<0$. A zero row and column may be removed without changing the rank. If $a>0$, write the remaining matrix as}{式~\eqref{eq:F-M-rational}中的有限矩阵可始终以有理数进行对称 Schur 消元检验。将副本指标递增排序，并初始化 $A=\mathsf M$。每一步读取当前对角元 $a$：若 $a<0$，则矩阵不是半正定的；若 $a=0$，则该行其余元素必须精确为零，否则相应 $2\times2$ 主子式的行列式为 $-|b|^2<0$。零行和零列可删除，秩保持不变。若 $a>0$，将当前矩阵写为}
\begin{equation}
 A=\begin{pmatrix}a&b^\dagger\\b&C\end{pmatrix}
  =\begin{pmatrix}1&0\\b/a&I\end{pmatrix}
    \begin{pmatrix}a&0\\0&C-bb^\dagger/a\end{pmatrix}
    \begin{pmatrix}1&b^\dagger/a\\0&I\end{pmatrix}.
 \label{eq:F-Schur}
\end{equation}
\bi{This invertible congruence reduces positivity to that of the smaller Schur complement $C-bb^\dagger/a$ and increases the accumulated rank by one. Iterating proves positive semidefiniteness if every pivot is positive or has an exactly zero remaining row. Rational subtraction and division determine the decision without a tolerance for small eigenvalues.}{该可逆合同变换将正性归结为更小的 Schur 补 $C-bb^\dagger/a$ 的正性，并使累计秩增加一。反复执行后，若每个主元均为正，或其为零且剩余行精确为零，则证明半正定性。有理数减法与除法确定所有判定，无需为小本征值设置容差。}

\bi{A second finite check can be obtained from characteristic coefficients. For a block of dimension $d=|\mathcal R_D|$, set $J_0=I_d$ and recursively compute}{另一种有限检验可由特征多项式系数得到。对维数为 $d=|\mathcal R_D|$ 的分块，令 $J_0=I_d$，并递归计算}
\begin{equation}
 X_b=\mathsf M J_{b-1},\qquad
 c_b=\frac{\operatorname{Tr}X_b}{b},\qquad
 J_b=c_bI_d-X_b,\qquad 1\le b\le d.
 \label{eq:F-characteristic-recurrence}
\end{equation}
\bi{Newton's identities, or direct coefficient comparison in the polynomial adjugate identity, give}{由 Newton 恒等式，或直接比较多项式伴随矩阵恒等式中的系数，可得}
\begin{equation}
 \det(xI_d+\mathsf M)=x^d+c_1x^{d-1}+\cdots+c_d.
 \label{eq:F-characteristic}
\end{equation}
\bi{For example, the recursion starts with $c_1=\operatorname{Tr}\mathsf M$ and $c_2=((\operatorname{Tr}\mathsf M)^2-\operatorname{Tr}\mathsf M^2)/2$. If all $c_b\ge0$, the polynomial in \eqref{eq:F-characteristic} is positive for every $x>0$. Since $\mathsf M$ is real symmetric, a negative eigenvalue would give a positive root, a contradiction. Thus nonnegative characteristic coefficients are an independent exact criterion for this calculation.}{例如，递归前两步为 $c_1=\operatorname{Tr}\mathsf M$ 及 $c_2=((\operatorname{Tr}\mathsf M)^2-\operatorname{Tr}\mathsf M^2)/2$。若所有 $c_b\ge0$，则式~\eqref{eq:F-characteristic}中的多项式在每个 $x>0$ 处均为正。由于 $\mathsf M$ 为实对称矩阵，若存在负本征值，就会给出一个正根，形成矛盾。因此特征系数非负为本计算提供了另一项精确判据。}

\bi{Evaluating the finite sums \eqref{eq:F-U}, \eqref{eq:F-Y-rational}, \eqref{eq:F-L-rational}, and \eqref{eq:F-Z-rational} at $\varepsilon_{\mathrm F}=3/10^6$ gives the following results. The rank is the number of positive Schur pivots. The sign string lists $c_1,\ldots,c_d$ in order; a plus means a strictly positive rational numerator with positive denominator, and zero means exact zero. Every block passes the Schur test, including its zero-row checks, and all characteristic coefficients are nonnegative.}{以 $\varepsilon_{\mathrm F}=3/10^6$ 计算有限和~\eqref{eq:F-U}、\eqref{eq:F-Y-rational}、\eqref{eq:F-L-rational}及~\eqref{eq:F-Z-rational}，得到下表。秩等于正 Schur 主元的个数；符号串依次列出 $c_1,\ldots,c_d$，其中加号表示分母为正时分子严格为正，零表示精确为零。每个分块均通过 Schur 检验及其零行检验，全部特征系数也均非负。}
\begin{center}
\small
\begin{tabular}{r|r|r|l}
$D$&$d$&$\operatorname{rank}(\mathsf M)$&\bi{Signs}{符号}\\ \hline
1&1&0&\texttt{0}\\
2&1&0&\texttt{0}\\
3&2&0&\texttt{00}\\
4&2&0&\texttt{00}\\
5&3&1&\texttt{+00}\\
6&3&0&\texttt{000}\\
7&4&1&\texttt{+000}\\
8&4&1&\texttt{+000}\\
9&5&2&\texttt{++000}\\
10&5&1&\texttt{+0000}\\
11&6&3&\texttt{+++000}\\
12&6&2&\texttt{++0000}\\
13&7&4&\texttt{++++000}\\
14&7&3&\texttt{+++0000}\\
15&8&5&\texttt{+++++000}\\
16&8&4&\texttt{++++0000}\\
17&9&6&\texttt{++++++000}\\
18&9&5&\texttt{+++++0000}\\
19&10&7&\texttt{+++++++000}\\
20&10&6&\texttt{++++++0000}\\
21&11&8&\texttt{++++++++000}\\
22&11&7&\texttt{+++++++0000}\\
23&12&9&\texttt{+++++++++000}
\end{tabular}
\end{center}
\bi{The matrices have dimension at most twelve, and every sum ranges over the printed integer rows and explicitly bounded indices. Equations~\eqref{eq:F-rational-congruence} and~\eqref{eq:F-Schur}, or alternatively \eqref{eq:F-characteristic}, therefore prove all twenty-three instances of \eqref{eq:F-four-certificate}. Their kernels are retained throughout.}{这些矩阵的维数至多为十二，全部求和仅涉及已列出的整数行和具有显式界的指标。因此由式~\eqref{eq:F-rational-congruence}与~\eqref{eq:F-Schur}，或采用式~\eqref{eq:F-characteristic}，便证明式~\eqref{eq:F-four-certificate}的全部二十三个实例，且全过程均保留其核。}

\bi{Finally, the finite arithmetic and the analytic tail bound give the exact limiting budget}{最后，有限算术与解析尾界给出精确的极限预算}
\begin{equation}
 1-\eta-\frac{61}{4096}-1222\varepsilon_{\mathrm F}
     =\frac{4616733319001}{10^{14}}>0.
\end{equation}
\bi{The extension from these finite matrices to the full torus Hamiltonian is supplied by the positive congruence and the $H$-relative estimates in Section~\ref{sec:F-proof}; finite arithmetic alone is not a substitute for those analytic steps.}{由这些有限矩阵推广至完整环面哈密顿量，依靠的是第~\ref{sec:F-proof}节中的保正合同变换及相对于 $H$ 的估计；有限算术本身不能替代这些解析步骤。}

\section{\bi{Exact definition and verification of the bosonic finite certificate}{玻色有限证书的精确定义与验证}}
\label{sec:B-certificate}

\bi{This appendix specifies the finite arithmetic input to \eqref{eq:B-certificate-three}--\eqref{eq:B-certificate-four}. The construction starts from rational row data, forms the physical normal-ordered coefficients, takes a finite center-of-mass partial trace, and tests twenty-six rational quadratic forms. Each step is stated explicitly so that the finite verification can be reproduced independently of the implementation that generated the rows. The analytic transfer to the torus and the occupation estimates are proved in Section~\ref{sec:B-bosonic}; they are not inferred from these finite tests.}
{本附录明确给出式~\eqref{eq:B-certificate-three}--\eqref{eq:B-certificate-four} 的有限算术输入。构造从有理行数据出发，形成物理正规序系数，取有限质心偏迹，最后检验二十六个有理二次型。以下逐步写明其定义，使有限验证可以独立于生成这些行的实现而复现。到环面的解析移植和占据数估计已在第~\ref{sec:B-bosonic} 节证明，并非从这些有限检验推断。}

\subsection{\bi{Frozen rows and canonical polynomial coordinates}{冻结行与正则多项式坐标}}
\label{sec:B-certificate-input}

\bi{The numerical input consists of the integers in the accompanying file \texttt{data/\allowbreak bosonic/\allowbreak exact\_\allowbreak shifted\_\allowbreak certificate\_\allowbreak rows.json}, taken from the fixed source accompanying Shen~\cite{Shen2026BosonicCandidate}. No fitted coefficient is recomputed for the torus. For channel $t=0,1,2,3$, the ordered list of columns is}
{数值输入由随附文件 \texttt{data/\allowbreak bosonic/\allowbreak exact\_\allowbreak shifted\_\allowbreak certificate\_\allowbreak rows.json} 中的整数构成，取自 Shen~\cite{Shen2026BosonicCandidate} 的固定附属来源。环面上不重新拟合任何系数。对通道 $t=0,1,2,3$，列按以下顺序排列：}
\begin{equation}
 \mathcal I_t=(\varnothing)\,\Vert\,
      \bigl((p,d-p):d=t,t+1,\ldots,8;\ p=0,1,\ldots,d\bigr).
 \label{eq:B-ordered-columns}
\end{equation}
\bi{Here $\varnothing$ is the first column, for the pair term, and $\Vert$ denotes concatenation. The file stores this list as \texttt{terms}, the rows as \texttt{integer\_\allowbreak scaled\_\allowbreak rows}, and the common denominator as \texttt{common\_\allowbreak denominator}. The complete sizes are}
{其中 $\varnothing$ 是配对项所在的首列，$\Vert$ 表示串接。文件分别用 \texttt{terms} 保存该列序、用 \texttt{integer\_\allowbreak scaled\_\allowbreak rows} 保存整数行、用 \texttt{common\_\allowbreak denominator} 保存共同分母。全部尺寸为}
\begin{equation}
 \begin{array}{c|rrrr}
 t&0&1&2&3\\ \hline
 |\mathcal I_t|&46&45&43&40\\
 \#\{\rho:t_\rho=t\}&31&37&36&34
 \end{array},\qquad \mathrm{denominator}=10^{12}.
 \label{eq:B-row-sizes}
\end{equation}
\bi{Dividing the stored integer $m_{\rho,a}$ by $10^{12}$ defines $y_{\rho,a}$, including $y_{\rho,\varnothing}=y_0^\rho$. Equation~\eqref{eq:B-row-decoding}, with the positive square root for each $\kappa_p$, uniquely decodes the physical real coefficients. Equivalently, the row itself is the rational polynomial operator \eqref{eq:B-rational-row}. These specifications fix the input without any floating-point rounding.}
{将存储整数 $m_{\rho,a}$ 除以 $10^{12}$ 定义 $y_{\rho,a}$，其中 $y_{\rho,\varnothing}=y_0^\rho$。式~\eqref{eq:B-row-decoding} 中各 $\kappa_p$ 取正平方根，便唯一确定物理实系数；等价地，每行就是式~\eqref{eq:B-rational-row} 的有理多项式算符。这些规定不经任何浮点舍入即固定了输入。}

\bi{For $k$ particles and total orbital degree $d$, use the lexicographically ordered basis}
{对 $k$ 个粒子和总轨道次数 $d$，采用按字典序排列的基}
\begin{equation}
 \mathcal B_{k,d}
  =\{A=(a_1,\ldots,a_k):0\le a_1\le\cdots\le a_k,
                              \ \textstyle\sum_\ell a_\ell=d\},
 \qquad |A]_x=\prod_{j\ge0}x_j^{n_j(A)}|0\rangle.
 \label{eq:B-rational-basis}
\end{equation}
\bi{The occupation $n_j(A)$ counts occurrences of $j$ in $A$. The Gram matrix in this basis is diagonal:}
{占据数 $n_j(A)$ 是 $j$ 在 $A$ 中出现的次数。此基的 Gram 矩阵为对角矩阵：}
\begin{equation}
 G_{k,d}=\operatorname{diag}(g_A:A\in\mathcal B_{k,d}),\qquad
 g_A=\prod_{j\ge0}n_j(A)!(j!)^{n_j(A)}.
 \label{eq:B-integer-metric}
\end{equation}
\bi{A coefficient matrix acts on column vectors. If $O$ is an operator on this degree block, its quadratic-form matrix is $G_{k,d}O$. Physical adjoints obey $\partial_j^*=x_j/j!$ and $x_j^*=j!\partial_j$, rather than $\partial_j^*=x_j$. For a map $F$ between two blocks, its adjoint matrix is $G_{\rm in}^{-1}F^{\mathsf T}G_{\rm out}$. Thus all adjoints and all form matrices can be computed over the rational numbers.}
{算符系数矩阵作用于列向量。若 $O$ 是该次数块上的算符，其二次型矩阵为 $G_{k,d}O$。物理伴随满足 $\partial_j^*=x_j/j!$ 和 $x_j^*=j!\partial_j$，而非 $\partial_j^*=x_j$。对两个分块间的映射 $F$，其伴随矩阵为 $G_{\rm in}^{-1}F^{\mathsf T}G_{\rm out}$。因此，所有伴随和二次型矩阵都可在有理数域上计算。}

\bi{For completeness, every monomial matrix entry is fixed by the elementary action}
{为使定义完整，每个单项式矩阵元由以下初等作用唯一确定：}
\begin{equation}
 x^\mu\partial^\nu\prod_jx_j^{n_j}
   =\left[\prod_j(n_j)_{\nu_j}\right]
       \prod_jx_j^{n_j-\nu_j+\mu_j},
 \qquad
 (n)_r=\begin{cases}n!/(n-r)!,&n\ge r,\\0,&n<r.\end{cases}
 \label{eq:B-monomial-action}
\end{equation}
\bi{The coefficient is zero whenever a required annihilation is unavailable. Ordered pair sums are retained as ordered sums, including both terms when the two orbital labels differ; this is required by $\widehat B_p=\sum_{a+b=p}\partial_a\partial_b$. No additional factorial is attached to the physical lift.}
{当所需湮灭无法完成时，系数为零。配对求和始终保留有序求和形式，两个轨道标号不同时两项都计入，这是 $\widehat B_p=\sum_{a+b=p}\partial_a\partial_b$ 的定义所要求的。物理提升不再附加额外阶乘。}

\subsection{\bi{Direct rational construction of the body coefficients}{体阶系数的直接有理构造}}
\label{sec:B-certificate-matrices}

\bi{It is sufficient to form the rational row Gram matrices}
{只需形成有理行 Gram 矩阵}
\begin{equation}
 \Gamma_{ab}^{(t)}
     =\sum_{\rho:t_\rho=t}y_{\rho,a}y_{\rho,b}
     =10^{-24}\sum_{\rho:t_\rho=t}m_{\rho,a}m_{\rho,b},
 \qquad a,b\in\mathcal I_t.
 \label{eq:B-row-gram}
\end{equation}
\bi{Write $P_p=\widehat B_p^*=\sum_{a+b=p}x_ax_b/(a!b!)$. For a noninitial column $a=(p,j)$ set $i=p+j-t$, and for $b=(p',j')$ set $i'=p'+j'-t$. Directly commuting $\partial_i$ through $x_{i'}$ gives the following rational normal-ordered operators:}
{记 $P_p=\widehat B_p^*=\sum_{a+b=p}x_ax_b/(a!b!)$。对非首列 $a=(p,j)$，令 $i=p+j-t$；对 $b=(p',j')$，令 $i'=p'+j'-t$。直接将 $\partial_i$ 与 $x_{i'}$ 对易，得到以下有理正规序算符：}
\begin{align}
 \mathsf M_3
 ={}&\sum_{t=0}^3\sum_{a=(p,j)\in\mathcal I_t\setminus\{\varnothing\}}
       \Gamma_{\varnothing a}^{(t)}
       \left(P_t x_i\partial_j\widehat B_p
                 +(P_t x_i\partial_j\widehat B_p)^*\right)
 \notag\\
 &+\sum_{t=0}^3\sum_{a=(p,j),\,b=(p',j')}
       \Gamma_{ab}^{(t)}\frac{i!}{j!}\,\delta_{ii'}
           P_p x_j\partial_{j'}\widehat B_{p'},
 \label{eq:B-rational-M3}\\
 \mathsf M_4
 ={}&\sum_{t=0}^3\sum_{a=(p,j),\,b=(p',j')}
       \Gamma_{ab}^{(t)}\frac{i!}{j!}
           P_p x_jx_{i'}\partial_i\partial_{j'}\widehat B_{p'}.
 \label{eq:B-rational-M4}
\end{align}
\bi{The double sums run over noninitial columns of $\mathcal I_t$. For example, $(x_i\partial_j\widehat B_p)^*=(i!/j!)P_px_j\partial_i$, which explains the factorial ratio in both formulas. These equations are precisely \eqref{eq:B-row-three}--\eqref{eq:B-row-four} after the change of variables, and they need no recursive subtraction between particle sectors. Each monomial in $\mathsf M_k$ has exactly $k$ creators and $k$ annihilators. Restricting it to $k$ particles therefore gives its physical $k$-body coefficient itself.}
{双重求和遍历 $\mathcal I_t$ 的非首列。例如，$(x_i\partial_j\widehat B_p)^*=(i!/j!)P_px_j\partial_i$，这解释了两式中的阶乘比。以上两式恰为式~\eqref{eq:B-row-three}--\eqref{eq:B-row-four} 在坐标变换后的形式，不需要在不同粒子数扇区间作递归扣除。$\mathsf M_k$ 的每个单项式恰有 $k$ 个创建和 $k$ 个湮灭算符，因此将其限制到 $k$ 粒子空间，就得到物理 $k$ 体系数本身。}

\bi{Let $M_{k,d}$ be the matrix obtained from \eqref{eq:B-rational-M3} or \eqref{eq:B-rational-M4} on $\mathcal B_{k,d}$ using \eqref{eq:B-monomial-action}. Its form matrix is $G_{k,d}M_{k,d}$. The three-body matrix vanishes for $d>8$ and the four-body matrix vanishes for $d>16$, because the creators' total orbital degrees in the displayed monomials obey these bounds. The two-body row contribution is fixed separately by}
{用式~\eqref{eq:B-monomial-action}，将式~\eqref{eq:B-rational-M3} 或式~\eqref{eq:B-rational-M4} 作用于 $\mathcal B_{k,d}$，所得矩阵记为 $M_{k,d}$，其二次型矩阵为 $G_{k,d}M_{k,d}$。三体矩阵在 $d>8$ 时为零，四体矩阵在 $d>16$ 时为零，因为上述单项式中创建算符的总轨道次数满足这些上界。两体行贡献则单独由下式固定：}
\begin{equation}
 \eta=\sum_{t=0}^3\Gamma_{\varnothing\varnothing}^{(t)}
                         \frac{2^{t+1}}{t!}.
 \label{eq:B-eta-rational}
\end{equation}

\bi{The planar target matrices have equally explicit rational definitions. On $\mathcal B_{k,d}$, let $H_{k,d}^{\rm pl}$ be the matrix of}
{平面目标矩阵也有同样显式的有理定义。在 $\mathcal B_{k,d}$ 上，令 $H_{k,d}^{\rm pl}$ 为下列算符的矩阵：}
\begin{equation}
 H_k^{\rm pl}=\sum_{p\ge0}\frac{p!}{2^{p+1}}P_p\widehat B_p.
 \label{eq:B-rational-H}
\end{equation}
\bi{Only $p\le d$ contribute. For three particles define $S_{3,d}^{\rm pl}=(H_{3,d}^{\rm pl})^2-H_{3,d}^{\rm pl}$. This is the connected coefficient because the planar pair operator is a projection and no four-body term acts on three particles. For four particles the zero-contraction coefficient on degree $d$ is the matrix of}
{只有 $p\le d$ 的项有贡献。对三粒子，定义 $S_{3,d}^{\rm pl}=(H_{3,d}^{\rm pl})^2-H_{3,d}^{\rm pl}$。由于平面配对算符是投影，且四体项在三粒子空间上不作用，它正是连通系数。对四粒子，次数 $d$ 的零收缩系数由下列算符的矩阵给出：}
\begin{equation}
 S_{4,d}^{\rm pl}
   =\left.\sum_{p+q=d}\frac{p!q!}{2^{d+2}}
                     P_pP_q\widehat B_p\widehat B_q
                    \right|_{\mathcal B_{4,d}}.
 \label{eq:B-rational-S4}
\end{equation}
\bi{Here $p+q=d$ because the four annihilators leave the vacuum. The ordered sum over $p,q$ is the zero-contraction sum from the square of $\sum_pB_p^\dagger B_p$, so its normalization agrees with the physical $S_4$.}
{这里必须有 $p+q=d$，因为四次湮灭后留下真空。对 $p,q$ 的有序求和就是 $\sum_pB_p^\dagger B_p$ 平方的零收缩求和，因此其归一化与物理 $S_4$ 一致。}

\subsection{\bi{Relative blocks and the exact center-of-mass trace}{相对分块与精确质心偏迹}}
\label{sec:B-certificate-relative}

\bi{On the $k$-particle space, the physical oscillator operators become the integer-coefficient polynomial operators}
{在 $k$ 粒子空间上，物理振子算符化为整数系数多项式算符}
\begin{equation}
 L_{\rm osc}=\sum_{j\ge1}j x_{j-1}\partial_j,
 \qquad L_{\rm osc}^\dagger=\sum_{j\ge0}x_{j+1}\partial_j,
 \qquad [L_{\rm osc},L_{\rm osc}^\dagger]=kI.
 \label{eq:B-rational-ladder}
\end{equation}
\bi{For each $z$, form the matrix of $L_{\rm osc}:\mathcal B_{k,z}\to\mathcal B_{k,z-1}$ by \eqref{eq:B-monomial-action}. A rational column matrix $U_{k,z}$ spanning its kernel defines the full relative space. A deterministic choice is the usual reduced-row-echelon kernel basis: scan columns in the order \eqref{eq:B-rational-basis}, use the first available pivot, and list the free-variable vectors in increasing column order. No orthogonalization is needed.}
{对每个 $z$，使用式~\eqref{eq:B-monomial-action} 形成 $L_{\rm osc}:\mathcal B_{k,z}\to\mathcal B_{k,z-1}$ 的矩阵。以有理列矩阵 $U_{k,z}$ 张成其核，即定义完整相对空间。可以确定性地选取通常的简化行阶梯形零空间基：按式~\eqref{eq:B-rational-basis} 的列序扫描，选取首个可用主元，并按列序递增排列自由变量基向量。此处不需要正交化。}

\bi{Write $D_{k,z}^{(m)}=(L_{\rm osc}^\dagger)^{m-z}U_{k,z}$ for the unnormalized descendants in total degree $m\ge z$. The oscillator relation implies that the normalized descendant is $D_{k,z}^{(m)}/\sqrt{k^{m-z}(m-z)!}$; the relative basis itself may remain nonorthonormal. Let $d_3=8$ and $d_4=16$. The relative form matrix of the averaged row coefficient is then}
{记 $D_{k,z}^{(m)}=(L_{\rm osc}^\dagger)^{m-z}U_{k,z}$ 为总次数 $m\ge z$ 的未归一化后裔。振子关系给出规范化后裔为 $D_{k,z}^{(m)}/\sqrt{k^{m-z}(m-z)!}$，而相对基本身可以继续保持非正交。令 $d_3=8$、$d_4=16$，则平均行系数的相对二次型矩阵为}
\begin{equation}
 \overline M_{k,z}
   =\frac2k\sum_{m=z}^{d_k}
       \frac{(D_{k,z}^{(m)})^{\mathsf T}
                 G_{k,m}M_{k,m}D_{k,z}^{(m)}}
            {k^{m-z}(m-z)!}.
 \label{eq:B-finite-partial-trace}
\end{equation}
\bi{This is the full center-of-mass partial trace, with the same factor $2/k$ as \eqref{eq:B-torus-seed}. Every center-of-mass descendant with nonzero anchored coefficient is included. The sum is finite because the anchored degree support is finite. The square-root descendant normalization has canceled between the two sides of the quadratic form, so every matrix entry remains rational.}
{这正是完整质心偏迹，其因子 $2/k$ 与式~\eqref{eq:B-torus-seed} 相同。所有具有非零锚定系数的质心后裔均被纳入。由于锚定次数支撑有限，该和是有限和。后裔归一化中的平方根在二次型两侧相消，因此每个矩阵元仍为有理数。}

\bi{The relative target forms are}
{相对目标二次型为}
\begin{equation}
 h_{k,z}=U_{k,z}^{\mathsf T}G_{k,z}H_{k,z}^{\rm pl}U_{k,z},
 \qquad
 s_{k,z}=U_{k,z}^{\mathsf T}G_{k,z}S_{k,z}^{\rm pl}U_{k,z}.
 \label{eq:B-relative-targets}
\end{equation}
\bi{Thus the twenty-six matrices that must be positive are exactly}
{因此，须检验为正的二十六个矩阵精确为}
\begin{align}
 E_{3,z}&=s_{3,z}-\overline M_{3,z}+\delta_Bh_{3,z},
                      &&z=0,\ldots,8,
 \label{eq:B-exact-E3}\\
 E_{4,z}&=\alpha_Bs_{4,z}-\overline M_{4,z}+\delta_Bh_{4,z},
                      &&z=0,\ldots,16.
 \label{eq:B-exact-E4}
\end{align}
\bi{All symbols on the right-hand sides have been defined by finite sums, rational arithmetic, and the fixed integer input. Positive semidefiniteness in this nonorthogonal relative basis is equivalent to the physical operator inequalities, because $U_{k,z}$ is an injective parametrization of the entire relative space.}
{右侧所有符号均已通过有限求和、有理数运算和固定整数输入定义。由于 $U_{k,z}$ 是整个相对空间的单射参数化，在这一非正交相对基中的半正定性等价于物理算符不等式。}

\subsection{\bi{Exact positivity, ranks, and zero kernels}{精确正性、秩与零核}}
\label{sec:B-certificate-positivity}

\bi{Each form matrix is first checked to be symmetric over the rational numbers. The positivity test repeatedly chooses a strictly positive diagonal entry $p$, permutes it to the first position, and replaces the remaining block $C$ by its Schur complement $C-bb^{\mathsf T}/p$. The justification is the exact congruence}
{首先在有理数域上检查每个二次型矩阵的对称性。正性检验反复选取一个严格为正的对角元 $p$，将其置换到首位，并将剩余块 $C$ 替换为 Schur 补 $C-bb^{\mathsf T}/p$。其依据是精确合同关系}
\begin{equation}
 \begin{pmatrix}p&b^{\mathsf T}\\ b&C\end{pmatrix}
 \sim
 \begin{pmatrix}p&0\\0&C-bb^{\mathsf T}/p\end{pmatrix}.
 \label{eq:B-Schur-congruence}
\end{equation}
\bi{At every stage, all diagonal entries are required to be nonnegative. If no positive diagonal entry remains, every remaining entry must be exactly zero. This last check is essential: a nonzero off-diagonal entry in a matrix with zero diagonal would prevent positive semidefiniteness. Termination with positive pivots and an exact zero block proves positive semidefiniteness by the displayed congruence, and the number of positive pivots is its rank. No numerical eigenvalue threshold is used.}
{每一步均要求所有对角元非负。若不再存在正对角元，则所有剩余矩阵元都必须精确为零。最后这一检验不可省略：零对角矩阵中的非零非对角元会破坏半正定性。依照上述合同关系，以正主元和精确零块结束即证明半正定性，而正主元个数就是矩阵秩。整个过程不使用数值本征值阈值。}

\bi{For each block, rational row reduction gives $\ker h_{k,z}$. Verify exactly that $E_{k,z}\ker h_{k,z}=0$. Independently apply the same positive-pivot test to $h_{k,z}$ and compare ranks. The resulting dimensions and ranks of all blocks are as follows; each rank shown is simultaneously the rank of $E_{k,z}$ and $h_{k,z}$.}
{对每个分块，用有理行消元计算 $\ker h_{k,z}$，并精确检验 $E_{k,z}\ker h_{k,z}=0$。再独立地对 $h_{k,z}$ 作同样的正主元检验并比较秩。全部分块的维数和秩如下，其中每个所列秩同时是 $E_{k,z}$ 与 $h_{k,z}$ 的秩。}
\begin{equation}
 \begin{array}{c|rrrrrrrrr}
 k=3,\ z&0&1&2&3&4&5&6&7&8\\ \hline
 \dim&1&0&1&1&1&1&2&1&2\\
 \operatorname{rank}&1&0&1&1&1&1&1&1&1
 \end{array}
 \label{eq:B-three-block-table}
\end{equation}
\begin{equation}
 \begin{array}{c|rrrrrrrrr}
 k=4,\ z&0&1&2&3&4&5&6&7&8\\ \hline
 \dim&1&0&1&1&2&1&3&2&4\\
 \operatorname{rank}&1&0&1&1&2&1&3&2&4
 \end{array}
 \label{eq:B-four-block-table-one}
\end{equation}
\begin{equation}
 \begin{array}{c|rrrrrrrr}
 k=4,\ z&9&10&11&12&13&14&15&16\\ \hline
 \dim&3&5&4&7&5&8&7&10\\
 \operatorname{rank}&3&5&4&6&5&7&6&8
 \end{array}
 \label{eq:B-four-block-table-two}
\end{equation}
\bi{The input includes both degree-one zero-dimensional blocks. Across the other blocks the procedure produces sixty-seven strictly positive rational pivots; the remaining Schur blocks are identically zero. The full fractions are included with the machine-readable certificate results. The zero-kernel checks and equal ranks show that each $E_{k,z}$ is strictly positive on the physical image of $H_{k,z}^{\rm pl}$ and vanishes on its kernel. In particular, this proves \eqref{eq:B-certificate-three}--\eqref{eq:B-certificate-four} as positive coefficient inequalities on the full relative multiplicity spaces.}
{输入包括次数一的两个零维块。在其余分块上，该过程产生六十七个严格为正的有理主元，剩余 Schur 块恒为零。全部分数随机器可读证书结果提供。零核检验和秩相等说明，每个 $E_{k,z}$ 在 $H_{k,z}^{\rm pl}$ 的物理像空间上严格为正，并在其零核上消失。特别地，这证明了式~\eqref{eq:B-certificate-three}--\eqref{eq:B-certificate-four} 在完整相对重数空间上的正系数不等式。}

\bi{The row constant used for the torus commutator error is also obtained directly from the same integers:}
{环面对易误差所用的行常数也直接从同一组整数得到：}
\begin{equation}
 C_{\rm rows}^B
  =\sum_{t=0}^3\sum_{\rho:t_\rho=t}
          \sum_{(p,j)\in\mathcal I_t\setminus\{\varnothing\}}
    (y_{\rho,pj})^2\frac{2^{p+1}(p+j-t)!}{p!j!}.
 \label{eq:B-Crows-rational}
\end{equation}
\bi{Equations~\eqref{eq:B-eta-rational} and \eqref{eq:B-Crows-rational} give exactly the fractions in \eqref{eq:B-eta} and \eqref{eq:B-row-error-constants}. Thus the positivity certificate and the cost of its noncanonical torus commutators are computed from one and the same fixed set of physical row coefficients.}
{式~\eqref{eq:B-eta-rational} 与式~\eqref{eq:B-Crows-rational} 精确给出式~\eqref{eq:B-eta} 和式~\eqref{eq:B-row-error-constants} 中的分数。因此，正性证书及其非正则环面对易关系的代价，均由同一组固定物理行系数计算。}

\subsection{\bi{Exact scalar arithmetic for the explicit threshold}{显式阈值的精确标量算术}}
\label{sec:B-certificate-budget}

\bi{The scalar budget uses only rational numbers after the shell estimates have replaced exponentials by powers of two. With $e=1/10^8$, $s_e=1+e$, $c_e=(1+e)/(1-e)^2$, and $\kappa_e=6C_{\rm rows}^Be$, define}
{壳层估计将指数函数替换为二的幂之后，标量预算只包含有理数。取 $e=1/10^8$、$s_e=1+e$、$c_e=(1+e)/(1-e)^2$、$\kappa_e=6C_{\rm rows}^Be$，定义}
\begin{align}
 \underline A
 &=1-e-\eta-s_e[c_9+(\delta_B+\kappa_e)C_3^B]
                      -\delta_BC_4^Bs_e^2,
 \label{eq:B-exact-budget-A-def}\\
 \underline B
 &=1-\alpha_B(1+e)-s_ec_e[c_9+(\delta_B+\kappa_e)C_3^B]
                      -2\delta_BC_4^Bs_e^2c_e.
 \label{eq:B-exact-budget-B-def}
\end{align}
\bi{Their exact values are}
{其精确值为}
\begin{align}
 \underline A
 &=\frac{8183525179652905540939777339971061478405006776240354085927}
         {24266058979973529600000000000000000000000000000000000000000},
 \label{eq:B-exact-A20}\\
 \underline B
 &=\frac{19326695545052423240596864889482025919460640111495553399803410151}
         {1213302924732617621356765299867648000000000000000000000000000000000}.
 \label{eq:B-exact-B20}
\end{align}
\bi{Multiplication by the positive denominators verifies $\underline A>337241/10^6>1/3$ and $\underline B>15928/10^6>0$. If decimal lower bounds are desired, valid downward truncations are $\underline A\ge0.3372416256964785$ and $\underline B\ge0.015928994442432054$. Rounded decimal displays are approximations only; the exact fractions, rather than floating-point rounding, determine the signs in the proof.}
{乘以正分母即可检验 $\underline A>337241/10^6>1/3$ 和 $\underline B>15928/10^6>0$。若需要十进制下界，可使用向下截断 $\underline A\ge0.3372416256964785$ 和 $\underline B\ge0.015928994442432054$。舍入后的小数显示仅表示近似，证明中决定符号的是精确分数，而非浮点舍入。}

\bi{For reference, the rational shell bounds used before the common allowance are}
{为便于核对，使用共同误差上界之前的有理壳层界为}
\begin{align}
 b_4&=\frac{50\cdot17^2\cdot16!\cdot20^{32}}{2^{288}},
 &b_8&=\frac{50\cdot9^2\cdot8!\cdot20^{16}}{2^{288}},
 &b_2&=\frac{50}{2^{144}}.
 \label{eq:B-rational-shell-values}
\end{align}
\bi{Each is strictly less than $1/10^8$, as are $50/2^{288}$ and $50/2^{576}$ for the diagonal kernels. Consequently, the common substitution in \eqref{eq:B-exact-budget-A-def}--\eqref{eq:B-exact-budget-B-def} enlarges every error in the safe direction. This finite arithmetic closes the analytic estimate; it does not replace the positive thermal transfer, the connected-kernel identity, or the all-center frame bounds used to obtain that estimate.}
{以上三者均严格小于 $1/10^8$，对角核的两个界 $50/2^{288}$ 和 $50/2^{576}$ 也如此。因此，式~\eqref{eq:B-exact-budget-A-def}--\eqref{eq:B-exact-budget-B-def} 中的共同代入将每个误差都沿安全方向放大。这一有限算术完成了解析估计，但并不替代推导该估计所需的正热核移植、连通核恒等式和全质心框架界。}

\section{\bi{Why hard orbital truncation does not follow from the transfer}{为何硬轨道截断不由本迁移推出}}\label{app:disk}
\bi{A hard disk cutoff is not a magnetic-periodic quotient. Let $P_m$ retain the planar orbitals $e_0,\ldots,e_m$ and compress the fermionic relative-degree-one projector to $\bigwedge^2P_mU$. The highest total-degree sector $2m-1$ contains only $e_{m-1}\wedge e_m$. The pair coefficient~\eqref{eq:F-pair} therefore gives its exact eigenvalue}
{硬disk截断不是磁周期商空间。令$P_m$保留平面轨道$e_0,\ldots,e_m$，把费米相对次数一的投影子压缩到$\bigwedge^2P_mU$。最高总次数$2m-1$扇区只含$e_{m-1}\wedge e_m$。Pair系数~\eqref{eq:F-pair}因此给出其精确本征值}
\begin{equation}
 E_m=\frac{(2m-2)!}{2^{2m-2}(m-1)!m!}
 =\frac1m\frac{\binom{2m-2}{m-1}}{4^{m-1}}
 \sim\frac1{\sqrt\pi m^{3/2}}>0.
\end{equation}
\bi{For example, $E_{12}=29393/2097152$. Hence a particle-number and cutoff-uniform positive gap for all these disk compressions is false already with two particles. Compression does not preserve a squared lower bound: for any many-body projection $P$,}
{例如$E_{12}=29393/2097152$。因此，这些disk压缩模型对粒子数及截断一致的正谱隙在两粒子层面就不成立。压缩不保持平方下界：对任意多体投影$P$，}
\begin{equation}
 (PHP)^2=PH^2P-PH(I-P)HP.
\end{equation}
\bi{The omitted term is nonnegative and cannot be discarded in a lower-bound proof. In the present method the exact failure is visible in the pair map: the compressed disk map has a highest-COM singular value squared $E_m$, so its Gram operator is not uniformly close to the identity, even on its nonzero support. Local disk jets also cease to be uniformly close to planar jets outside the droplet, and there is no lattice-image identity~\eqref{eq:gram-unfold}. The full-COM estimate and unfolding, rather than agreement on a few low orbitals, distinguish the torus construction.}
{被遗漏项非负，在下界证明中不能删除。在本文方法中，具体失效处可直接由pair映射看出：压缩后的disk映射在最高COM处有奇异值平方$E_m$，故其Gram算符即使限制到非零支撑，也不一致趋近恒等算符。Disk局部jets在液滴外同样不再一致接近平面jets，且不存在格点image恒等式~\eqref{eq:gram-unfold}。区分环面构造的是全COM估计与展开，而非少数低轨道上的一致。}

\section{\bi{Reproducibility}{可复现性}}\label{app:reproducibility}

\bi{The seven fermionic rows come from OpenAI~\cite[Appendix~A]{OpenAI2026Fock}; the 138 bosonic rows come from Shen~\cite{Shen2026BosonicCandidate}, at the fixed versions linked in the bibliography. The torus argument uses the bosonic source's twenty-six planar blocks, not its 156 finite-sphere interval checks. The retained data and exact constructions are specified in the certificate appendices.}
{费米七行数据来自OpenAI~\cite[附录A]{OpenAI2026Fock}，玻色138行数据来自Shen~\cite{Shen2026BosonicCandidate}；固定版本见参考文献链接。环面论证使用玻色来源的26个平面块，不使用其156个有限球面区间检查。保留的数据及精确构造见证书附录。}

\bi{The accompanying Python programs reconstruct the finite coefficient forms and scalar budgets using exact rational arithmetic. The bosonic reconstruction checks all twenty-six planar blocks and their zero kernels; the fermionic reconstruction checks the twenty-three four-body blocks, finite three-body margins, and scalar gap budget. Singular positive semidefinite blocks are tested with exact zero-row and kernel checks, without numerical regularization.}
{随附Python程序以精确有理算术重建有限系数型和标量预算。玻色重建检查全部26个平面块及其零核；费米重建检查23个四体块、有限三体余量及标量谱隙预算。奇异半正定块使用精确零行和核检查，不作数值正则化。}

\bi{These finite arithmetic checks do not establish the magnetic-periodization identities, the COM trace interchange, the physical normalization correspondence, or the occupation bounds. Those analytic steps are supplied by the candidate arguments in the main text. Selected formal endpoint checks do not constitute formal verification of every theorem and corollary here; the detailed scope is recorded in the accompanying repository's \texttt{VERIFICATION.md}.}
{这些有限算术检查不证明磁周期化恒等式、COM偏迹交换、物理归一化对应关系或占据数界。正文的候选论证提供这些解析环节。选定形式化终点的检查不构成对本文每项定理和推论的形式化验证；详细范围记录在随附仓库的\texttt{VERIFICATION.md}中。}

\bi{Finite-size normalization diagnostics gave a fermionic three-body residual about $7.21\times10^{-12}$ at $M=9$, and bosonic positive energies about $0.700809668$ at $(M,N)=(6,3)$ and $0.586180647$ at $(8,4)$. These calculations are diagnostics, not inputs to either uniform bound.}
{有限尺寸归一化诊断在$M=9$时给出约$7.21\times10^{-12}$的费米三体残差；玻色$(M,N)=(6,3)$和$(8,4)$时的正能量约为$0.700809668$和$0.586180647$。这些计算属于诊断，不是两项一致下界的输入。}

\bi{Both language editions share equations, constants, assumptions, and labels through one paired-prose source, compiled with XeLaTeX and BibTeX. Build commands, exact data, and executable finite checks accompany the manuscript. AI contributions and human roles are documented in the repository's \texttt{AI\_CONTRIBUTIONS.md}, following the explicit disclosure practice of Shen~\cite{Shen2026AIContributions}.}
{两个语言版本通过同一配对文本源文件共享方程、常数、假设及标签，使用XeLaTeX和BibTeX编译。构建命令、精确数据和可执行有限检查随稿提供。AI贡献及人类角色记录在仓库的\texttt{AI\_CONTRIBUTIONS.md}中，披露方式参考Shen的明确声明~\cite{Shen2026AIContributions}。}

\bibliographystyle{plain}
\bibliography{references}
\end{document}
